\chapter{Monstrous moonshine: McKay--Thompson series, Hauptmoduln and the statement}\label{ch:16}

\providecommand{\chsixteenH}{\mathbb{H}}
\providecommand{\chsixteenR}{\mathbb{R}}
\providecommand{\chsixteenM}{\mathbb{M}}
\providecommand{\chsixteenSL}{\mathrm{SL}_2(\mathbb{Z})}
\providecommand{\chsixteenSLR}{\mathrm{SL}_2(\mathbb{R})}
\providecommand{\chsixteenGLR}{\mathrm{GL}_2(\mathbb{R})}
\providecommand{\chsixteenP}{\mathbb{P}^1}
\providecommand{\chsixteenX}{X(\Gamma)}
\providecommand{\chsixteenT}{T_g}
\providecommand{\chsixteenmat}[4]{\left(\begin{smallmatrix}#1&#2\\#3&#4\end{smallmatrix}\right)}
\providecommand{\chsixteenlean}[1]{\texttt{\directlua{ local s = "\luaescapestring{\detokenize{#1}}" for piece, sep in string.gmatch(s, "([^_.]*)([_.]?)") do if not (piece == "" and sep == "") then tex.sprint(-2, piece .. sep) if not (sep == "") then tex.sprint(string.char(92) .. "allowbreak{}") end end end}}}
\chaptermark{Monstrous moonshine: the statement}
\begingroup\setlength{\emergencystretch}{3em}\sloppy

This chapter documents the release of the library named in the front matter; every Lean quotation below is checked by
\texttt{scripts/textbook/sync.py} to occur verbatim (up to whitespace and comments) in the source files of that
release.  It
opens the part of the book devoted to \emph{monstrous moonshine}: the Conway--Norton conjecture, proved by
Borcherds in 1992, and the action of the Monster on the moonshine module $\Vnat$.  Chapters~\ref{ch:02}
to~\ref{ch:14} constructed $\Vnat$ as a vertex operator algebra.  The present chapter states, in Lean and on paper,
what it means for $\Vnat$ to exhibit \emph{moonshine}.  Chapters~\ref{ch:17}--\ref{ch:20} prove it for every
automorphism of $\Vnat$; Chapters~\ref{ch:21}--\ref{ch:22} identify the automorphism group with the Monster.

\paragraph{The result in one sentence.}  The library proves, with no hypotheses and only Lean's three standard
axioms (\lean{propext}, \lean{Classical.choice}, \lean{Quot.sound}; Chapter~\ref{ch:01}, \S\ref{ch01:sec:trust}), the theorem
\begin{equation*}
  \lean{VOA.Moonshine.ConwayNorton.conwayNorton : ConwayNortonConjecture},
\end{equation*}
which says: \emph{for every automorphism $g$ of the moonshine vertex operator algebra $\Vnat$, the McKay--Thompson
series $T_g(\tau)=q^{-1}+\sum_{n\ge1}\tr(g\,|\,\Vnat_{n+1})\,q^n$ is the normalised Hauptmodul of a genus-zero group
$\Gamma_g\le\chsixteenSLR$ that contains $\Gamma_0(N)$ for some $N\ge1$ and is commensurable with $\chsixteenSL$.}
That the automorphism group of $\Vnat$ is the Monster is a separate matter, treated in Chapters~\ref{ch:21}
and~\ref{ch:22}: there the library proves, again with no hypotheses, that $\Aut(\Vnat)$ is isomorphic to a subgroup
of $\mathrm{AutO}$, an independently constructed Lean model of the Monster taken from the AxiomMath library (a Lean
formalisation of parts of the classification of finite simple groups, CFSG), which contains a certain generated
subgroup $\mathrm{Gmat}$ (all of this is explained in \S\ref{ch16:sec:gmat}), and,
\emph{conditionally on one named hypothesis} (\lean{GmatEqAutO}: $\mathrm{Gmat}=\mathrm{AutO}$, a theorem of the
AxiomMath CFSG library that is not available on this library's toolchain), that $\Aut(\Vnat)\cong\mathrm{AutO}$ and
that monstrous moonshine holds for every element of $\mathrm{AutO}$.

\paragraph{What this chapter does.}  Its job is to make the statement checkable.  (The ``Fidelity Protocol'' mentioned below is the book's
fixed format for presenting a definition or theorem: statement, verbatim Lean, dictionary, proof on paper, physical
meaning, machine status; Chapter~\ref{ch:01}, \S\ref{ch01:sec:boxes}.)  It
\begin{enumerate}[label=(\arabic*)]
\item tells the story of moonshine and lays out the logical dependencies of Chapters~\ref{ch:16}--\ref{ch:22}
  (\S\ref{ch16:sec:history}--\S\ref{ch16:sec:roadmap});
\item defines every object in the statement: the named vertex operator algebra \lean{VnatVOA}, its automorphism
  group \lean{AutVnat}, graded traces and the McKay--Thompson series, the upper half-plane and its cusps, the
  normalised Hauptmodul of a genus-zero group, moonshine groups, and the conjecture itself, each under the Fidelity
  Protocol (\S\ref{ch16:sec:objects}--\S\ref{ch16:sec:statement});
\item presents the auxiliary notions on which the proof runs: complete replicability, replicable families and
  rationality of traces (\S\ref{ch16:sec:replication});
\item proves, on paper and in Lean, the step from $\Gamma_1(N)$ to $\Gamma_0(N)$ that converts the output of the
  genus-zero argument of Chapter~\ref{ch:19} into the Conway--Norton statement (\S\ref{ch16:sec:gammazero});
\item records the final theorem and the audit modules that check it (\S\ref{ch16:sec:final}), the recorded traps
  (\S\ref{ch16:sec:traps}), the referee's checklist and the references.
\end{enumerate}
The most important section for a referee is \S\ref{ch16:sec:statement}, the clause-by-clause comparison of the Lean
proposition \lean{ConwayNortonConjecture} with the conjecture of Conway and Norton.  As everywhere in this book, our
concern is the \emph{meaning} of the Lean statements; proofs are explained mathematically, and we do not describe how
the Lean proof scripts are organised.

%======================================================================
\section{Monstrous moonshine: the story}\label{ch16:sec:history}
%======================================================================

\subsection{Two numbers}

In 1978 John McKay noticed that
\begin{equation}\label{ch16:eq:mckay}
  196884 \;=\; 1 + 196883 .
\end{equation}
The left-hand side is the first nontrivial coefficient of the normalised elliptic modular function
\begin{equation}\label{ch16:eq:J}
\begin{aligned}
  J(\tau) \;=\; j(\tau)-744 \;&=\; q^{-1} + 196884\,q + 21493760\,q^2 + 864299970\,q^3\\
  &\qquad + 20245856256\,q^4+\cdots,\qquad q=e^{2\pi i\tau},
\end{aligned}
\end{equation}
the function of which every \emph{modular function} for $\chsixteenSL$ is a rational function; here a modular function is
a $\chsixteenSL$-invariant meromorphic function on the upper half-plane $\chsixteenH=\{\tau\in\C:\operatorname{Im}\tau>0\}$
that is also meromorphic at the cusp, i.e.\ has a $q$-expansion with only finitely many negative powers of $q$
(Chapter~\ref{ch:05}, Theorem~\ref{ch05:thm:haupt}; without the cusp condition the claim is false, e.g.\ for
$e^{J}$).  The
right-hand side involves $196883$, the dimension of the smallest nontrivial irreducible complex representation of
the \emph{Monster} $\chsixteenM$, the largest of the twenty-six sporadic finite simple groups, whose existence had been
predicted by Fischer and Griess in 1973 and which Griess constructed in 1980--82 as the automorphism group of a
commutative non-associative algebra of dimension $196884$, the \emph{Griess algebra}.  Its order is
\begin{equation*}
  |\chsixteenM| = 2^{46}\cdot3^{20}\cdot5^{9}\cdot7^{6}\cdot11^{2}\cdot13^{3}\cdot17\cdot19\cdot23\cdot29\cdot31\cdot41
  \cdot47\cdot59\cdot71\approx 8.08\times10^{53}.
\end{equation*}
John Thompson observed that the coincidence continues: if $1$, $196883$, $21296876$, $842609326,\dots$ are the
dimensions of the irreducible representations of $\chsixteenM$ in increasing order, then
\begin{equation*}
  21493760 = 1+196883+21296876,\qquad 864299970 = 2\cdot1+2\cdot196883+21296876+842609326,
\end{equation*}
and so on.  Thompson proposed that there should be a natural \emph{graded} representation
$V=\bigoplus_{n\ge0}V_n$ of $\chsixteenM$ whose graded dimension $\sum_n(\dim V_n)\,q^{n-1}$ is $J$, and that the
\emph{graded traces}
\begin{equation}\label{ch16:eq:Tg-intro}
  T_g(\tau) \;=\; \sum_{n\ge0}\tr(g\,|\,V_n)\,q^{n-1}, \qquad g\in\chsixteenM,
\end{equation}
now called the \emph{McKay--Thompson series}, should also be interesting functions.

\subsection{The conjecture of Conway and Norton}

Conway and Norton (1979) computed the first coefficients of $T_g$ for all $194$ conjugacy classes of $\chsixteenM$
from the (then conjectural) character table and the (then conjectural) decomposition of the first few $V_n$, and
found that every one of them appears to be the \emph{normalised Hauptmodul} of a \emph{genus-zero} subgroup of
$\chsixteenSLR$.  The terms will be defined carefully in \S\ref{ch16:sec:haupt}, and the congruence subgroups
$\Gamma_0(N)=\{\chsixteenmat abcd\in\chsixteenSL: N\mid c\}$ and $\Gamma_1(N)\subseteq\Gamma_0(N)$ (also $a\equiv d\equiv1$
mod $N$) in \S\ref{ch16:sec:analytic}.  Two subgroups $H,K$ of a group are \emph{commensurable} if $H\cap K$ has finite
index in both $H$ and $K$.  For the moment: a discrete group
$\Gamma\le\chsixteenSLR$ commensurable with $\chsixteenSL$ acts on $\chsixteenH$, the quotient $\Gamma\backslash
\chsixteenH$ can be compactified by adding finitely many \emph{cusps} to a compact Riemann surface
$\chsixteenX$, and $\Gamma$ has \emph{genus zero} if $\chsixteenX$ is a sphere.  In that case there is an isomorphism
$\chsixteenX\to\chsixteenP(\C)$ sending the cusp $\infty$ to $\infty$, and it can be normalised uniquely so that its
expansion at $\infty$ is $q^{-1}+0+O(q)$; that function is the \emph{normalised Hauptmodul} (German:
``principal modulus'').  For example $J$ is the normalised Hauptmodul of $\chsixteenSL$, and
\begin{equation}\label{ch16:eq:T2B}
  T_{2B}(\tau)=\Bigl(\frac{\eta(\tau)}{\eta(2\tau)}\Bigr)^{24}+24
  = q^{-1}+276\,q-2048\,q^2+11202\,q^3-49152\,q^4+\cdots
\end{equation}
is the normalised Hauptmodul of $\Gamma_0(2)$, while
\begin{equation}\label{ch16:eq:T2A}
  T_{2A}(\tau)=\Bigl(\frac{\eta(\tau)}{\eta(2\tau)}\Bigr)^{24}+24+2^{12}\Bigl(\frac{\eta(2\tau)}{\eta(\tau)}\Bigr)^{24}
  = q^{-1}+4372\,q+96256\,q^2+1240002\,q^3+\cdots
\end{equation}
is the normalised Hauptmodul of the larger group $\Gamma_0(2)+$ generated by $\Gamma_0(2)$ and the Fricke
involution $\tau\mapsto-1/(2\tau)$.  (Here $\eta(\tau)=q^{1/24}\prod_{n\ge1}(1-q^n)$ is Dedekind's eta function, and
the coefficients were checked numerically for this chapter by expanding the products; $276=1+275$ and
$4372=1+4371$, where $275$ and $4371$ are the values of the $196883$-dimensional character on the classes $2B$
and~$2A$.)  The \emph{Conway--Norton conjecture}, in the form proved by Borcherds, reads:

\begin{quote}\itshape
There is a graded representation $V=\bigoplus_{n\ge0}V_n$ of the Monster with graded dimension $J$ such that, for
every $g\in\chsixteenM$, the McKay--Thompson series $T_g$ is the normalised Hauptmodul of a genus-zero group
$\Gamma_g\le\chsixteenSLR$ which contains $\Gamma_0(N)$ for some $N$ and is commensurable with $\chsixteenSL$.
\end{quote}
Conway and Norton said much more: they proposed for every class an explicit group $\Gamma_g$ and an explicit
Hauptmodul.  Their groups have the form ``$\Gamma_0(n|h)+e,f,\dots$'', where $n$ is the order of $g$, $h\mid\gcd(24,n)$,
\begin{equation*}
  \Gamma_0(n|h)=\Bigl\{\chsixteenmat{a}{b/h}{cn}{d}: a,b,c,d\in\Z,\ ad-bcn/h=1\Bigr\}
\end{equation*}
(conjugate to $\Gamma_0(n/h)$ by $\tau\mapsto h\tau$, and containing $\Gamma_0(nh)$), and ``$+e,f,\dots$'' means that
certain Atkin--Lehner involutions $W_e,W_f,\dots$ (for exact divisors $e,f,\dots$ of $n/h$; the Fricke involution is
$W_{n/h}$) are adjoined.  These groups lie between $\Gamma_0(N)$, $N=nh$, and its normaliser in $\chsixteenSLR$.  This
notation is for orientation only; nothing in the library uses it. there are $171$ distinct functions among the $194$
classes.  The Lean statement of this library is the \emph{qualitative} form displayed above; the relation between
the two is discussed in detail in \S\ref{ch16:sec:fidelity}.  Two striking consequences: the \emph{supersingular
primes} $2,3,5,7,11,13,17,19,23,29,31,41,47,59,71$ (the primes $p$ for which $\Gamma_0(p)+$ has genus zero, as Ogg
had observed in 1974) are exactly the primes dividing $|\chsixteenM|$; and all the coefficients of all $T_g$ are
integers.

\subsection{The moonshine module and Borcherds' proof}

Frenkel, Lepowsky and Meurman (1984, 1988) constructed Thompson's representation $V$: the \emph{moonshine module}
$\Vnat$, a vertex operator algebra of central charge $24$ obtained from the Leech lattice by a $\Z/2$-orbifold, with
$\Vnat_1=0$, $\dim\Vnat_2=196884$ and graded dimension $J$, and they showed that its automorphism group is the
Monster.  This library constructs $\Vnat$ in Chapters~\ref{ch:02}--\ref{ch:14}; the headline is the theorem
\lean{VOA.Moonshine.moonshine} of Chapter~\ref{ch:14}.

Borcherds (1992) proved the Conway--Norton conjecture for FLM's $\Vnat$.  His argument has five steps, and the
library follows it closely for the first four:
\begin{enumerate}[label=(B\arabic*)]
\item Tensor $\Vnat$ with the vertex algebra $V_{I\!I_{1,1}}$ of the even unimodular Lorentzian lattice of rank
  two; the result has central charge $26$.
\item The \emph{no-ghost theorem} of string theory (Goddard--Thorn) identifies the space of physical states of
  momentum $(m,n)\in I\!I_{1,1}$, $(m,n)\ne0$, with $\Vnat_{mn+1}$, $\chsixteenM$-equivariantly.  The physical states
  modulo null states form a Lie algebra, the \emph{Monster Lie algebra} $\mathfrak m$.
\item $\mathfrak m$ is a Borcherds--Kac--Moody (generalised Kac--Moody) algebra, and its Weyl--Kac--Borcherds
  denominator identity is the product formula
  \begin{equation}\label{ch16:eq:denominator}
    p^{-1}\prod_{m>0,\;n\in\Z}(1-p^mq^n)^{c(mn)} = J(p)-J(q),\qquad J(q)=\sum_n c(n)q^n .
  \end{equation}
\item Taking $g$-traces in the homology of a free Lie subalgebra of $\mathfrak m$ gives \emph{twisted denominator
  identities}, which say that the $T_g$ are \emph{completely replicable} (\S\ref{ch16:sec:replication}).
\item A completely replicable function whose first few coefficients agree with those of the conjectured Hauptmodul
  is that Hauptmodul; Borcherds verified the agreement class by class.
\end{enumerate}
Step (B5) is a finite case check that needs the character table of $\chsixteenM$.  Cummins and Gannon (1997) showed
how to replace it by a uniform argument: \emph{completely replicable functions with suitable growth are Hauptmoduln of
genus-zero groups}.  The library uses a case-free route of this kind (Chapter~\ref{ch:19}), re-proved and corrected
where needed, so that no character table, no list of $171$ groups and no case check enter the formal proof.

\begin{physics}
\textbf{Moonshine as a statement about a conformal field theory.}  $\Vnat$ is the chiral algebra of a holomorphic
two-dimensional conformal field theory with central charge $c=24$; physically, $24$ free chiral bosons compactified
on the Leech torus $\chsixteenR^{24}/\Leech$ and orbifolded by $x\mapsto-x$.  The Monster is its \emph{global symmetry group}:
it commutes with the stress tensor and preserves all operator product expansions.  The graded dimension
$\sum_n\dim\Vnat_n\,q^{n-1}=\tr_{\Vnat}q^{L_0-c/24}$ is the torus partition function $Z(\tau)$ (the shift by
$c/24=1$ is the Casimir energy), and its $\chsixteenSL$-invariance is the modular invariance of the theory on a torus
of modulus $\tau$.  The McKay--Thompson series is the \emph{twisted partition function}
\begin{equation*}
  T_g(\tau)=Z\bigl[\,{}^{\;1}_{\;g}\,\bigr](\tau)=\tr_{\Vnat}\bigl(g\,q^{L_0-1}\bigr),
\end{equation*}
the path integral on the torus with the symmetry defect $g$ inserted along the time circle.  Modular
transformations move the defect: $\tau\mapsto-1/\tau$ exchanges ``$g$ along time'' with ``$g$ along space'', which is
the partition function of the \emph{$g$-twisted sector}.  This is why one expects each $T_g$ to be invariant under a
subgroup $\Gamma_g$ of finite index (the subgroup that brings the defect back to itself), as in the orbifold
constructions of Dixon--Ginsparg--Harvey (1988) and of Dong--Li--Mason.  The \emph{genus-zero} property is much
stronger and has no generally accepted physical derivation; the most suggestive heuristics connect it to the
uniqueness of $\Vnat$ (Tuite 1995) and to three-dimensional gravity (Witten 2007).  We label these as heuristics: none
of them is used by the library.

\textbf{Borcherds' proof as string theory.}  Steps (B1)--(B3) are the quantisation of the bosonic string in $26$
dimensions with $24$ of them compactified on the moonshine orbifold and two forming a Lorentzian torus $I\!I_{1,1}$.
The physical states of the string at momentum $r$ (the BRST cohomology at ghost number one, or the old covariant
quantisation: Virasoro-primary weight-one states modulo null states) form the Lie algebra $\mathfrak m$; the no-ghost
theorem says they are counted by the transverse ($24$-dimensional) excitations.  The denominator identity
\eqref{ch16:eq:denominator} is then a statement about the spectrum of this string theory, and the twisted
denominator identities are its versions with a symmetry twist.
\end{physics}

%======================================================================
\section{The logical structure of the moonshine part}\label{ch16:sec:roadmap}
%======================================================================

\subsection{The dependency graph}

The part ``Monstrous moonshine'' (Chapters~\ref{ch:16}--\ref{ch:22}) documents the files of the library whose names
begin with \texttt{MoonshineCN} (written by Fedor Popov's side of the project) and \texttt{MoonshineLR} (Leonardo
Rastelli's side), the files \texttt{Aut*.lean} and the folder \texttt{MonsterGriess/}.  The chapters are ordered for
reading, \emph{not} in logical order: the dependencies form a graph, not a chain.  In particular, the theorems of
Chapter~\ref{ch:19} that produce complete replicability are \emph{conditional} on two inputs that are proved in
Chapter~\ref{ch:20}; and the theorem for $\Aut(\Vnat)$ needs only Chapters~\ref{ch:17}--\ref{ch:20} (and this chapter),
while Chapters~\ref{ch:21}--\ref{ch:22} are needed only to identify $\Aut(\Vnat)$ with the Monster.

The nodes are the following; ``uses'' lists the nodes whose results are used.
\begin{enumerate}[label=(\Alph*)]
\item\label{ch16:node:A} \textbf{$\Vnat$} (Chapters~\ref{ch:02}--\ref{ch:14}): the vertex operator algebra
  \lean{VnatVOA}, $c=24$, $\Vnat_1=0$, graded dimension $J$, simple.  \lean{AutVnat} and the McKay--Thompson series are
  defined in this chapter.
\item\label{ch16:node:B} \textbf{$V_{I\!I_{1,1}}$ and $\Vnat\otimes V_{I\!I_{1,1}}$} (Chapter~\ref{ch:17}): the vertex
  algebra of the hyperbolic plane (not a vertex operator algebra: its grading is not bounded below) and the tensor
  theory of central charge $26$ with its invariant form and Virasoro action.  Uses (A).
\item\label{ch16:node:C} \textbf{No-ghost theorem and $\mathfrak m$} (Chapter~\ref{ch:18}): physical states modulo null
  states, the root spaces $\mathfrak m_{(m,n)}\cong\Vnat_{mn+1}$, the bracket and invariant form, and $\mathfrak m$ as
  Borcherds data \emph{given} generation by its simple root spaces.  Uses (B).
\item\label{ch16:node:D} \textbf{Product identity} (Chapter~\ref{ch:19}): the denominator formula
  \eqref{ch16:eq:denominator}, proved analytically.  Uses (A).
\item\label{ch16:node:E} \textbf{Simple-root generation of $\mathfrak m$} (Chapter~\ref{ch:20}): a dimension sandwich
  between $\mathfrak m$ and the algebra generated by its simple root spaces.  Uses (C), (D).
\item\label{ch16:node:F} \textbf{Finite order} (Chapter~\ref{ch:20}): every automorphism of $\Vnat$ has finite order
  (rigid Ising vectors); also $\Aut(\Vnat)$ is finite.  Uses (A).
\item\label{ch16:node:G} \textbf{Twisted denominator identities and complete replicability} (Chapter~\ref{ch:19}):
  \lean{CompletelyReplicable}, as a theorem with hypotheses ``generation'' and ``finite order''.  Uses (C), and the
  hypotheses are discharged by (E), (F).
\item\label{ch16:node:H} \textbf{Genus zero, $\Gamma_1$ form} (Chapter~\ref{ch:19}): the case-free argument, from complete
  replicability to \lean{ConwayNortonGammaOne}.  Uses (F), (G) (through the abstract interface of
  \S\ref{ch16:sec:replication}).
\item\label{ch16:node:I} \textbf{From $\Gamma_1$ to $\Gamma_0$} (this chapter, \S\ref{ch16:sec:gammazero}), and the
  \textbf{assembly} \lean{conwayNorton} (Chapter~\ref{ch:20}; recorded in \S\ref{ch16:sec:final}).  Uses (F), (G), (H).
\item\label{ch16:node:J} \textbf{The Griess algebra and AxiomMath's Monster} (Chapter~\ref{ch:21}): an explicit
  isomorphism $\Vnat_2\cong\C\otimes_\Q\mathrm{SpaceB}$ of algebras with forms, where $\mathrm{SpaceB}$ is AxiomMath's
  rational Griess algebra and $\mathrm{AutO}$ its automorphism group.  Uses (A).
\item\label{ch16:node:K} \textbf{The Monster acts on $\Vnat$} (Chapter~\ref{ch:22}): every element of
  $\mathrm{Gmat}=\langle\mathrm{Gvis},\tau_B\rangle\le\mathrm{AutO}$ lifts to $\Vnat$; $\Aut(\Vnat)$ is isomorphic to the
  subgroup of liftable elements (via a commutant lemma); given \lean{GmatEqAutO}, $\Aut(\Vnat)\cong\mathrm{AutO}$ and
  monstrous moonshine holds for the Monster.  Uses (J), (I), and for rationality of traces (H).
\end{enumerate}
As a list of edges (``$X\Rightarrow Y$'': $Y$ uses $X$):
\begin{plaincode}
(A) ⇒ (B) ⇒ (C)                    V♮, the c = 26 theory, the Monster Lie algebra 𝔪
(A) ⇒ (D)                          product identity for J(p) − J(q)
(C), (D) ⇒ (E)                     𝔪 is generated by its simple root spaces
(A) ⇒ (F)                          every automorphism of V♮ has finite order
(C), (E), (F) ⇒ (G)                complete replicability of every T_g
(F), (G) ⇒ (H)                     genus zero, Γ₁ form
(F), (G), (H) ⇒ (I)                Γ₁ ⇒ Γ₀ and conwayNorton        [Chapters 16–20 only]
(A) ⇒ (J)                          V♮₂ ≅ ℂ ⊗ SpaceB
(H), (I), (J) ⇒ (K)                the Monster; GmatEqAutO          [Chapters 21–22]
\end{plaincode}
The Lean names of the terms used here (``rigid Ising vectors'', ``Borcherds data'', ``simple root spaces'', ``commutant
lemma'') are defined in the chapters indicated.

\subsection{The final theorems of the part}

The following table lists the final theorems informally; the precise Lean statements, with their fidelity checks,
are in the chapters indicated.  ``No hypotheses'' means: the Lean theorem has no explicit hypotheses, and its
\texttt{\#print axioms} output is \lean{[propext, Classical.choice, Quot.sound]}.
\begin{center}\small
\begin{longtable}{p{0.33\textwidth}p{0.38\textwidth}p{0.2\textwidth}}
\toprule
Lean name & informal content & hypotheses; where\\
\midrule
\endhead
\chsixteenlean{ConwayNorton.conwayNorton} & for every $g\in\Aut(\Vnat)$, $T_g$ is the normalised Hauptmodul of a group
  $\Gamma_g\supseteq\Gamma_0(N)$ commensurable with $\chsixteenSL$ (genus zero follows on paper,
  Prop.~\ref{ch16:prop:haupt}) & none; statement here, proof Ch.~\ref{ch:20}\\
\chsixteenlean{ConwayNorton.conwayNortonGammaOne_holds} & the same with $\Gamma_1(N)$ in place of $\Gamma_0(N)$, for every $g$
  of finite order & none; Ch.~\ref{ch:19}, \ref{ch:22}\\
\chsixteenlean{ConwayNorton.M8.autVnat_isOfFinOrder}, \chsixteenlean{ConwayNorton.M8.finite_autVnat} & every automorphism of $\Vnat$ has
  finite order; $\Aut(\Vnat)$ is finite & none; Ch.~\ref{ch:20}\\
\chsixteenlean{ConwayNorton.rationalTraces} & every coefficient of every $T_g$ is rational & none; Ch.~\ref{ch:22}\\
\chsixteenlean{ConwayNorton.conwayNorton_gmat} & every $g\in\mathrm{Gmat}$ is the restriction to $\Vnat_2\cong\C\otimes
  \mathrm{SpaceB}$ of an automorphism $h$ of $\Vnat$, and $T_h$ is as in \chsixteenlean{conwayNorton} & none; Ch.~\ref{ch:22}\\
\chsixteenlean{ConwayNorton.nonempty_autVnat_mulEquiv_liftable} & $\Aut(\Vnat)$ is isomorphic to the subgroup of $\mathrm{AutO}$
  of elements that act on $\Vnat$ & none; Ch.~\ref{ch:22}\\
\chsixteenlean{ConwayNorton.monsterEquivAutVnat} & a group isomorphism $\mathrm{AutO}\cong\Aut(\Vnat)$ compatible with the
  Griess algebras & \textbf{\chsixteenlean{GmatEqAutO}}; Ch.~\ref{ch:22}\\
\chsixteenlean{ConwayNorton.monstrousMoonshine} & monstrous moonshine for AxiomMath's Monster: an isomorphism
  $\varphi:\mathrm{AutO}\to\Aut(\Vnat)$ as above such that $T_{\varphi(g)}$ is a normalised Hauptmodul as in
  \chsixteenlean{conwayNorton}, for every $g$ & \textbf{\chsixteenlean{GmatEqAutO}}; Ch.~\ref{ch:22}\\
\bottomrule
\end{longtable}
\end{center}
All names are in the namespace \lean{VOA.Moonshine}.

\subsection{The one named hypothesis, and why it is a hypothesis}\label{ch16:sec:gmat}

The Monster enters the library through a port of definitions from the AxiomMath CFSG library (a Lean formalisation
of parts of the classification of finite simple groups, on Lean 4.35): the rational Griess algebra
$\mathrm{SpaceB}$ (a $\Q$-vector space of dimension $196884$ with a commutative product and a form), its
automorphism group $\mathrm{AutO}\le\mathrm{GL}_\Q(\mathrm{SpaceB})$, the ``visible'' subgroup $\mathrm{Gvis}$ (of
shape $2^{1+24}.\mathrm{Co}_1$), the extra element $\tau_B$, and the subgroup $\mathrm{Gmat}$ that they generate.  The
definitions are copied verbatim from AxiomMath's commit \texttt{69cb972f}; Chapter~\ref{ch:21} presents them.  The
proposition
\begin{leancode}
def GmatEqAutO : Prop :=
  Generated.Gmat = AutO
\end{leancode}
(\leanpath{VOA/Moonshine/MoonshineLRMonsterHeadline.lean}, \lean{VOA.Moonshine.ConwayNorton.GmatEqAutO}) says that the
generated group is the whole automorphism group.  It is a theorem of the AxiomMath library,
\lean{CFSG.Blueprint.NativeSporadic.MonsterModel.gmat_eq_griess}, at that commit.  Its proof passes through the
Stroth--Leon--Sims recognition of the Baby Monster; the project measured its import closure at about $16{,}000$ files
and $208$\,MB on a newer toolchain, so it cannot be ported to this
library's Lean~4.30 / Mathlib~v4.30 at reasonable cost.  It therefore enters as an \emph{explicit hypothesis}, never
as an axiom, of exactly two final theorems, \lean{monsterEquivAutVnat} and \lean{monstrousMoonshine}.

Three points must be kept in mind.
\begin{enumerate}[label=(\alph*)]
\item A theorem with the hypothesis \lean{(hM : GmatEqAutO)} is \textbf{conditional}, however clean its
  \texttt{\#print axioms} output; \texttt{\#print axioms} does not report explicit hypotheses
  (Chapter~\ref{ch:01}, \S\ref{ch01:sec:uncond}).  This book never calls such a theorem unconditional.
\item The project checked that the ported definitions agree with AxiomMath's up to definitional equality
  (\texttt{Notes/monster\_aut/CFSG\_PORT\_FIDELITY.md}; one definition, \lean{SpinLattice.phi0}, is an unbundled
  form of the original), so that \lean{GmatEqAutO} is the statement AxiomMath proves.  The project did \emph{not}
  check the axiom footprint of AxiomMath's proof inside AxiomMath's library.
\item That $\mathrm{AutO}$ is \emph{the Monster} (finite, simple, of the order displayed above) is likewise a
  theorem of AxiomMath's library; only the definitions are ported (the docstring of
  \leanpath{VOA/Moonshine/MonsterGriess/AutO.lean} says so).  Inside LeanVOA, ``the Monster'' means the Lean group
  \lean{AutO}.
\end{enumerate}
The Conway--Norton theorem \lean{conwayNorton} itself does \emph{not} depend on \lean{GmatEqAutO}: it is a statement
about $\Aut(\Vnat)$, whatever that group is.  How large $\Aut(\Vnat)$ is, and why the theorem is therefore not
vacuous, is discussed in \S\ref{ch16:sec:fidelity}.

\begin{physics}
Physically, the distinction is between ``the moonshine CFT has the McKay--Thompson twisted partition functions of
every one of its global symmetries as genus-zero Hauptmoduln'' (proved without hypotheses) and ``its global symmetry
group is precisely the Monster as constructed by AxiomMath'' (proved given one group-theoretic identification).  The
first is a statement about the theory; the second identifies its symmetry group with an independently constructed
finite group.
\end{physics}

\subsection{What is not formalised}

For the record, the following parts of the classical story are \emph{not} statements of the library.
\begin{itemize}
\item The identification of $\Gamma_g$ with the group of the Conway--Norton table, and of $N$ with $o(g)\,h$; the Lean
  statement is qualitative (\S\ref{ch16:sec:fidelity}).
\item Any statement about conjugacy classes or the character table of the Monster, the $194$ classes, the $171$
  functions.
\item Integrality of the coefficients of $T_g$ (rationality is proved: \lean{rationalTraces}, Chapter~\ref{ch:22}).
\item The uniqueness of $\Vnat$ among holomorphic $c=24$ vertex operator algebras with $V_1=0$ (an open problem in
  general), and the identification of the library's $\Vnat$ with FLM's beyond the paper argument recorded in
  Chapter~\ref{ch:14}.
\item Generalised moonshine (Norton), modularity of twisted modules, and the physical heuristics above.
\end{itemize}


%======================================================================
\section{The objects of the statement}\label{ch16:sec:objects}
%======================================================================

The Lean statement \lean{ConwayNortonConjecture} mentions, after all definitions are unfolded: the vertex operator
algebra \lean{VnatVOA}, its automorphisms, the trace of a linear map on a finite-dimensional space, power series, the
upper half-plane with its $q$-parameter, the action of $\chsixteenGLR$ on $\chsixteenH$ and on $\chsixteenP(\chsixteenR)$,
cusps, the congruence subgroups $\Gamma_0(N)$, $\Gamma_1(N)$, commensurability, and holomorphy.  This section
introduces each of them, with its Lean form.

\subsection{The moonshine module as a named vertex operator algebra}

Chapter~\ref{ch:14} ended with the theorem \lean{VOA.Moonshine.moonshine}, an existential statement: \emph{there is} a
vertex operator algebra structure on the carrier \lean{Vnat} with the listed properties, and its vertex map, vacuum and
conformal vector are pinned to explicit definitions.  To speak about automorphisms one needs the structure itself as a
named object.

\begin{definition}[The moonshine module and its automorphism group]\label{ch16:def:vnatvoa}
$\Vnat$ denotes the vertex operator algebra on the space $V_\Leech^+\oplus(V_\Leech^T)^+$ of
Chapters~\ref{ch:10}--\ref{ch:14}, with the block vertex map, vacuum and conformal vector constructed there.  An
\emph{automorphism} of $\Vnat$ is a linear bijection $g:\Vnat\to\Vnat$ with
\begin{equation*}
  g\bigl(u_{(n)}v\bigr)=(gu)_{(n)}(gv)\quad(u,v\in\Vnat,\ n\in\Z),\qquad g\mathbf 1=\mathbf 1,\qquad g\omega=\omega .
\end{equation*}
They form a group $\Aut(\Vnat)$ under composition.
\end{definition}
\begin{leancode}
noncomputable def VnatVOA : VertexOperatorAlgebra Vnat := vnatVOA vnat_twistedLocality

theorem VnatVOA_grading : VnatVOA.grading = gradingNat := rfl

theorem VnatVOA_grading_one : VnatVOA.grading 1 = ⊥ := gradingNat_one

theorem VnatVOA_finrank_two : Module.finrank ℂ (VnatVOA.grading 2) = 196884 :=
  finrank_gradingNat_two

abbrev AutVnat : Type := Aut VnatVOA
\end{leancode}
(\leanpath{VOA/Moonshine/AutVnat.lean}, namespace \lean{VOA.Moonshine.AutMonster}.)  The automorphism structure is
\begin{leancode}
structure Aut (W : VertexOperatorAlgebra 𝒱) where
  /-- The underlying `ℂ`-linear equivalence `α : 𝒱 ≃ₗ 𝒱`. -/
  toLinearEquiv : 𝒱 ≃ₗ[ℂ] 𝒱
  /-- `α` intertwines every product: `α(u₍ₙ₎v) = (αu)₍ₙ₎(αv)`. -/
  map_smul' : ∀ (u : 𝒱) (n : ℤ) (v : 𝒱),
    toLinearEquiv (W.smul u n v) = W.smul (toLinearEquiv u) n (toLinearEquiv v)
  /-- `α` fixes the vacuum: `α(𝟙) = 𝟙`. -/
  map_vacuum : toLinearEquiv W.vacuum = W.vacuum
  /-- `α` fixes the conformal vector: `α(ω) = ω`. -/
  map_omega : toLinearEquiv W.conformalVector = W.conformalVector
\end{leancode}
(\leanpath{VOA/Core/Automorphism.lean}, \lean{VOA.Aut}; it was presented in Chapter~\ref{ch:L04},
Definition~\ref{l04:def:aut}).

\begin{dictionary}
\textbf{\lean{VnatVOA}.}  \lean{vnatVOA htt} is the vertex operator algebra of Chapter~\ref{ch:10} built from the
block vertex map, whose only hypothesis is the twisted$\times$twisted locality \lean{htt}; Chapter~\ref{ch:14}
proved that locality, and \lean{vnat_twistedLocality} extracts it from the existential theorem
\lean{moonshine}.  So \lean{VnatVOA} is \emph{the} object of \lean{moonshine}, not a new construction: its vertex map,
vacuum and conformal vector are \lean{Ynat wbarL}, \lean{vacNat}, \lean{omegaNat} (the theorems
\lean{VnatVOA_Y}, \lean{VnatVOA_vacuum}, \lean{VnatVOA_conformalVector} in the same file state this by \lean{rfl}).
\lean{VnatVOA.grading n} is the weight-$n$ subspace $\Vnat_n$ (a \lean{Submodule ℂ Vnat}), equal by definition to
\lean{gradingNat n}, the grading of Chapter~\ref{ch:10}; \lean{VnatVOA_grading_one} says $\Vnat_1=0$ (\lean{⊥} is
the zero subspace) and \lean{VnatVOA_finrank_two} says $\dim\Vnat_2=196884$.

\textbf{\lean{Aut}.}  A term of \lean{Aut W} bundles a $\C$-linear equivalence \lean{toLinearEquiv} with three
proofs.  \lean{W.smul u n v} is the $n$-th product $u_{(n)}v$ (Chapter~\ref{ch:L03}); \lean{map_smul'} is
$g(u_{(n)}v)=(gu)_{(n)}(gv)$ for every $n\in\Z$, \lean{map_vacuum} is $g\mathbf 1=\mathbf 1$, \lean{map_omega} is
$g\omega=\omega$.  This is exactly the textbook definition of a VOA automorphism.  Fixing $\omega$ makes $g$ commute
with every $L_n=\omega_{(n+1)}$, hence preserve each $\Vnat_n$ (the $L_0$-eigenspaces); this is proved as
\lean{GriessAut.aut_map_grading} and restated below as \lean{vnatAut_preserves_grading}.  The group structure on
\lean{Aut W} (composition, identity, inverse) is defined in the same file; \lean{g * h} is ``apply $h$ first''.

\textbf{\lean{AutVnat}} is an \lean{abbrev}, i.e.\ a transparent name: \lean{AutVnat} \emph{is} \lean{Aut VnatVOA}.
The docstring's phrase ``the target is \lean{AutVnat ≃* 𝕄}'' records an intention; the identification with
the Monster is the subject of Chapter~\ref{ch:22} and is conditional on \lean{GmatEqAutO}.

\textbf{Could this be vacuous?}  If the structure on \lean{Vnat} were degenerate, $\Aut$ could be enormous and
traces meaningless; but \lean{VnatVOA} is simple, has $c=24$, $\Vnat_1=0$ and graded dimension $J$
(Chapter~\ref{ch:14}, \S\ref{ch14:sec:decode}).  If, at the other extreme, \lean{AutVnat} were trivial, the
conjecture would collapse to its $g=1$ case; Chapter~\ref{ch:22} proves that \lean{AutVnat} contains an isomorphic
copy of AxiomMath's group $\mathrm{Gmat}$, which contains $2^{1+24}.\mathrm{Co}_1$.
\end{dictionary}
\begin{physics}
\lean{VnatVOA} is the chiral algebra of the Monster CFT; \lean{AutVnat} is its group of global (internal) symmetries:
the operators that commute with the stress tensor ($g\omega=\omega$), fix the vacuum, and preserve every operator
product expansion ($g(u_{(n)}v)=(gu)_{(n)}(gv)$ is the statement that $g$ intertwines the OPE coefficients).
$\Vnat_1=0$ says the theory has no conserved currents, hence no continuous symmetries; that is why $\Aut(\Vnat)$ can
be (and, by Chapter~\ref{ch:20}, is) finite.
\end{physics}
\begin{status}
\lean{VnatVOA}, \lean{AutVnat}: definitions with no hypotheses.  The theorems quoted are unconditional and are used by
\lean{conwayNorton}, whose \texttt{\#print axioms} is \lean{[propext, Classical.choice, Quot.sound]}
(\S\ref{ch16:sec:final}).  The file \leanpath{VOA/Moonshine/AutVnat.lean} is documented in Chapter~\ref{ch:20}.
\end{status}

\subsection{The upper half-plane, the \texorpdfstring{$q$}{q}-parameter and the action of \texorpdfstring{$\mathrm{GL}_2(\mathbb{R})$}{GL2(R)}}\label{ch16:sec:analytic}

We collect the Mathlib notions that occur in the statement, each explained in plain words.
\begin{itemize}
\item \lean{ℍ} (Mathlib's \lean{UpperHalfPlane}) is the type of complex numbers with positive imaginary part; a term
  \lean{τ : ℍ} coerces to a complex number \lean{(τ : ℂ)}.
\item \lean{GL (Fin 2) ℝ} is the group $\chsixteenGLR$ of invertible real $2\times2$ matrices; \lean{γ 0 0},
  \lean{γ 0 1}, \lean{γ 1 0}, \lean{γ 1 1} are the entries $a,b,c,d$ of $\gamma=\chsixteenmat abcd$, and
  \lean{γ.det} its determinant (a unit of $\chsixteenR$, coerced to a real number by \lean{(γ.det : ℝ)}).  It acts on
  \lean{ℍ} by \lean{γ • τ}; for $\det\gamma>0$ this is the M\"obius action $\tau\mapsto(a\tau+b)/(c\tau+d)$ (Mathlib
  extends it to negative determinant by complex conjugation, which never occurs below).  The scalar matrices act
  trivially; in particular $-1$ acts trivially.
\item \lean{OnePoint ℝ} is $\chsixteenP(\chsixteenR)=\chsixteenR\cup\{\infty\}$, on which $\chsixteenGLR$ acts by the same M\"obius
  formula: $\gamma\cdot x=(ax+b)/(cx+d)$ and $\gamma\cdot\infty=a/c$ (or $\infty$ if $c=0$).
\item \lean{SL(2, ℤ)} is $\chsixteenSL$, and \lean{mapGL ℝ : SL(2, ℤ) →* GL (Fin 2) ℝ} the inclusion.  The notation
  \lean{𝒮ℒ} (scoped in \lean{MatrixGroups}) is Mathlib's name for its image, the subgroup $\chsixteenSL$ of
  $\chsixteenGLR$.  A subgroup of \lean{SL(2, ℤ)} is coerced to a subgroup of \lean{GL (Fin 2) ℝ} by mapping along
  \lean{mapGL ℝ}.
\item \lean{Gamma0 N} and \lean{Gamma1 N} are Mathlib's congruence subgroups of $\chsixteenSL$:
  \begin{equation*}
    \Gamma_0(N)=\Bigl\{\chsixteenmat abcd: c\equiv0\ (N)\Bigr\},\qquad
    \Gamma_1(N)=\Bigl\{\chsixteenmat abcd: a\equiv d\equiv1,\ c\equiv0\ (N)\Bigr\},
  \end{equation*}
  with congruences in \lean{ZMod N}.  (For $N=0$, ``mod $0$'' means equality, a degenerate case that the statement
  excludes by requiring $0<N$.)
\item \lean{Subgroup.HasDetOne Γ} (a Mathlib class, used as a proposition) says every element of \lean{Γ} has
  determinant~$1$: $\Gamma\le\chsixteenSLR$.
\item \lean{Subgroup.Commensurable H K} is Mathlib's commensurability:
  \lean{H.relIndex K ≠ 0 ∧ K.relIndex H ≠ 0}, where \lean{H.relIndex K} is the index $[K:H\cap K]$ and the value $0$
  encodes ``infinite''.  So $H\cap K$ has finite index in both $H$ and $K$.
\item \lean{IsCusp c Γ} (Mathlib) holds if some \lean{g ∈ Γ} is \emph{parabolic} (not scalar, with discriminant
  $(\operatorname{tr}g)^2-4\det g=0$) and fixes \lean{c}.  For a group commensurable with $\chsixteenSL$ the cusps are
  exactly the points of $\chsixteenP(\Q)=\Q\cup\{\infty\}$.
\item \lean{atImInfty} is the filter ``$\operatorname{Im}\tau\to\infty$'' on \lean{ℍ}: its neighbourhoods are the
  horoballs $\{\operatorname{Im}\tau>R\}$, so \lean{Tendsto f atImInfty (𝓝 w)} means $f(\tau)\to w$ as
  $\operatorname{Im}\tau\to\infty$ \emph{uniformly in $\operatorname{Re}\tau$}.
\item \lean{MDifferentiable 𝓘(ℂ) 𝓘(ℂ) f} says $f:\chsixteenH\to\C$ is complex-differentiable at every point of the
  complex manifold $\chsixteenH$, i.e.\ holomorphic.
\item \lean{Subgroup.adjoinNegOne Γ} is $\pm\Gamma=\{g: g\in\Gamma\ \text{or}\ -g\in\Gamma\}$.
\end{itemize}
The two local definitions of the statement file in this area are the cusp $\infty$ and the $q$-parameter.
\begin{leancode}
abbrev cuspInfty : OnePoint ℝ := OnePoint.infty

noncomputable def qPar (τ : ℍ) : ℂ := Function.Periodic.qParam 1 (τ : ℂ)
\end{leancode}
(\leanpath{VOA/Moonshine/MoonshineCNStatement.lean}.)
\begin{dictionary}
\lean{OnePoint.infty} is the point $\infty\in\chsixteenP(\chsixteenR)$.  Mathlib's \lean{Function.Periodic.qParam h z} is
$\exp(2\pi i z/h)$; with $h=1$, \lean{qPar τ} $=e^{2\pi i\tau}=q$.  Since $|q|=e^{-2\pi\operatorname{Im}\tau}<1$, $q$
never vanishes and $q\to0$ as $\operatorname{Im}\tau\to\infty$.  Width $1$ is the right choice for all
McKay--Thompson series: their invariance groups contain the translation $\tau\mapsto\tau+1$ and no smaller translation
(\S\ref{ch16:sec:haupt}).
\end{dictionary}
\begin{physics}
$q=e^{2\pi i\tau}$ is the exponentiated modulus of the torus on which the partition function is computed: $\tau$ is
the ratio of the two periods, and $q^{L_0-c/24}$ propagates states once around the time circle.  The cusp $\infty$
($q\to0$) is the limit of a very long torus, where the partition function is dominated by the vacuum, $Z\sim q^{-1}$.
\end{physics}
\begin{status}
Definitions; both printed in \leanpath{VOA/Moonshine/MoonshineCNAudit3.lean} with standard axioms.
\end{status}

%======================================================================
\section{Graded traces and the McKay--Thompson series}\label{ch16:sec:trace}
%======================================================================

\subsection{The action on a homogeneous piece, and its trace}

Every automorphism $g$ of a vertex operator algebra $V$ maps each weight space $V_n$ to itself, and $V_n$ is finite
dimensional; so $g$ has a trace on $V_n$.

\begin{definition}[Graded action and graded trace]\label{ch16:def:gradedtrace}
For a vertex operator algebra $V$, an integer $n$ and $g\in\Aut(V)$, let $g|_{V_n}\in\End(V_n)$ be the restriction,
and $\tr(g\,|\,V_n)$ its trace.  The map $g\mapsto g|_{V_n}$ is a monoid homomorphism $\Aut(V)\to\End(V_n)$.
\end{definition}
\begin{leancode}
noncomputable def gradedAutAction (V : VertexOperatorAlgebra M) (n : ℤ) :
    Aut V →* Module.End ℂ (V.grading n) where
  toFun g := (GriessAut.autRestrict g n).toLinearMap
  map_one' := by ext x; rfl
  map_mul' g h := by ext x; rfl

noncomputable def gradedAutTrace (V : VertexOperatorAlgebra M) (g : Aut V) (n : ℤ) : ℂ :=
  LinearMap.trace ℂ (V.grading n) (gradedAutAction V n g)
\end{leancode}
(\leanpath{VOA/Moonshine/MoonshineCNTrace.lean}, namespace \lean{VOA.Moonshine.ConwayNorton}.)
\begin{dictionary}
The carrier \lean{M} of \lean{V} is a section variable of the file, \lean{{M : Type u} [AddCommGroup M] [Module ℂ M]}:
an arbitrary complex vector space.  \lean{V.grading n} is the weight-$n$ subspace $V_n$, regarded as a type (a
submodule coerced to a type);
\lean{Module.End ℂ (V.grading n)} is $\End_\C(V_n)$.  \lean{GriessAut.autRestrict g n} is the restriction of $g$ to
$V_n$ as a linear equivalence $V_n\to V_n$ (defined in \leanpath{VOA/Core/GriessAutomorphism.lean} using
\lean{aut_map_grading}, the fact that $g$ maps $V_n$ onto $V_n$); \lean{.toLinearMap} forgets invertibility.
\lean{→*} means monoid homomorphism: \lean{map_one'} and \lean{map_mul'} say $1\mapsto\mathrm{id}$ and
$gh\mapsto g|_{V_n}\circ h|_{V_n}$ (both hold by definition, \lean{rfl}).  \lean{LinearMap.trace ℂ (V.grading n)} is
Mathlib's trace; it is the usual trace whenever the space is finite dimensional, which $V_n$ is by the axioms of
\lean{VertexOperatorAlgebra} (the instance \lean{V.finrank n}).  (On an infinite-dimensional space Mathlib's trace is
defined to be $0$; this junk value never arises here.)  The definition is the textbook one.
\end{dictionary}
\begin{physics}
$\tr(g\,|\,V_n)$ counts the states of energy $n$ weighted by their $g$-charge: in a basis of $g$-eigenvectors with
eigenvalues $\lambda_i$ it is $\sum_i\lambda_i$.  Summed against $q^{n-c/24}$ it becomes the partition function with
a $g$-twist along the time direction.
\end{physics}
\begin{status}
Definitions; printed in \leanpath{VOA/Moonshine/MoonshineCNAudit.lean} with standard axioms.
\end{status}

\begin{proposition}[Basic properties of graded traces]\label{ch16:prop:traces}
For every vertex operator algebra $V$, $n\in\Z$ and $g,h\in\Aut(V)$:
\begin{enumerate}[label=(\alph*)]
\item $\tr(1\,|\,V_n)=\dim V_n$;
\item $\tr(hgh^{-1}\,|\,V_n)=\tr(g\,|\,V_n)$ (traces are class functions);
\item for $V=\Vnat$: every $g\in\Aut(\Vnat)$ maps $\Vnat_n$ onto $\Vnat_n$.
\end{enumerate}
\end{proposition}
\begin{leancode}
theorem gradedAutTrace_one (V : VertexOperatorAlgebra M) (n : ℤ) :
    gradedAutTrace V 1 n = (Module.finrank ℂ (V.grading n) : ℂ)

theorem gradedAutTrace_conjugate (V : VertexOperatorAlgebra M) (g h : Aut V) (n : ℤ) :
    gradedAutTrace V (h * g * h⁻¹) n = gradedAutTrace V g n

theorem vnatAut_preserves_grading (g : AutVnat) (n : ℤ) :
    (VnatVOA.grading n).map g.toLinearEquiv.toLinearMap = VnatVOA.grading n
\end{leancode}
(\leanpath{VOA/Moonshine/MoonshineCNTrace.lean}.)
\begin{dictionary}
In (a), \lean{Module.finrank ℂ (V.grading n)} is $\dim_\C V_n$, cast to $\C$.  In (b), \lean{h * g * h⁻¹} is
composition in \lean{Aut V}.  In (c), \lean{(VnatVOA.grading n).map g.toLinearEquiv.toLinearMap} is the image
$g(\Vnat_n)$, and the equality of submodules says $g(\Vnat_n)=\Vnat_n$.  No hypotheses beyond the objects; all three
statements are the textbook ones.
\end{dictionary}
\paragraph{Proof on paper.}  (a) The identity has trace equal to the dimension.  (b) $\tr(ABA^{-1})=\tr(B)$ on a
finite-dimensional space, applied to the restrictions (restriction is a homomorphism).  (c) $g$ commutes with
$L_0=\omega_{(1)}$ because $g(\omega_{(1)}v)=(g\omega)_{(1)}(gv)=\omega_{(1)}(gv)$; hence it preserves each
$L_0$-eigenspace $\Vnat_n$, and so does $g^{-1}$, so $g(\Vnat_n)=\Vnat_n$. \qed
\begin{physics}
(b) is the statement that a partition function with a symmetry twist depends only on the conjugacy class of the
twist; (c) is energy conservation for a symmetry commuting with the Hamiltonian $L_0-c/24$.
\end{physics}
\begin{status}
Unconditional; printed in \leanpath{VOA/Moonshine/MoonshineCNAudit.lean}, standard axioms.
\end{status}

\subsection{The McKay--Thompson series as a power series}

\begin{definition}[McKay--Thompson series, regular part]\label{ch16:def:mckayThompson}
For $g\in\Aut(\Vnat)$ let
\begin{equation*}
  \mathrm{MT}_g \;=\;\sum_{n\ge0}\tr\bigl(g\,|\,\Vnat_{n+1}\bigr)\,X^n\;\in\;\C[[X]] .
\end{equation*}
Thus the McKay--Thompson series is $T_g(\tau)=q^{-1}+\mathrm{MT}_g(q)$, and $q\,T_g=1+q\,\mathrm{MT}_g(q)$.
\end{definition}
\begin{leancode}
noncomputable def mckayThompson (g : AutVnat) : PowerSeries ℂ :=
  PowerSeries.mk fun n => gradedAutTrace VnatVOA g ((n : ℤ) + 1)

noncomputable def mckayThompsonShift (g : AutVnat) : PowerSeries ℂ :=
  1 + PowerSeries.X * mckayThompson g
\end{leancode}
(\leanpath{VOA/Moonshine/MoonshineCNTrace.lean}.)
\begin{dictionary}
\lean{PowerSeries.mk a} is the formal power series $\sum_{n\ge0}a(n)X^n$; the $n$-th coefficient of
\lean{mckayThompson g} is \lean{gradedAutTrace VnatVOA g ((n : ℤ) + 1)} $=\tr(g\,|\,\Vnat_{n+1})$.  \textbf{The index
shift by one is the convention to remember}: the coefficient of $X^n$ is the trace on weight $n+1$, so that $X$
plays the role of $q$ in $T_g-q^{-1}$.  The vacuum (weight $0$) contributes the $q^{-1}$, which is not part of
\lean{mckayThompson}; it is added separately in \lean{mckayThompsonFun} below.  \lean{mckayThompsonShift g}
$=1+X\cdot$\lean{mckayThompson g} is $q\,T_g$, in the same normalisation as the library's \lean{JShift}
$=qJ$ (Chapter~\ref{ch:05}, Definition~\ref{ch05:def:jshift}); its constant term $1$ is $\tr(g|\Vnat_0)=1$.
\end{dictionary}
\begin{physics}
$q\,T_g=\tr_{\Vnat}(g\,q^{L_0})$, the twisted partition function with the Casimir factor $q^{-c/24}=q^{-1}$ removed;
the constant term $1$ is the unique vacuum, which every symmetry fixes.
\end{physics}
\begin{status}
Definitions; printed in \leanpath{VOA/Moonshine/MoonshineCNAudit.lean}, standard axioms.
\end{status}

\begin{proposition}[First properties of $\mathrm{MT}_g$]\label{ch16:prop:mtprops}
For all $g,h\in\Aut(\Vnat)$ and $n\in\N$:
\begin{enumerate}[label=(\alph*)]
\item the $n$-th coefficient of $\mathrm{MT}_g$ is $\tr(g\,|\,\Vnat_{n+1})$;
\item the constant coefficient vanishes: $\tr(g\,|\,\Vnat_1)=0$;
\item $\mathrm{MT}_{hgh^{-1}}=\mathrm{MT}_g$;
\item $\mathrm{MT}_1=\sum_{n\ge0}[q^{n+1}](qJ)\,X^n$, i.e.\ the coefficients of $\mathrm{MT}_1$ are those of $J$;
  equivalently $q\,T_1=qJ$;
\item $\mathrm{MT}_1\ne qJ$ (a reminder that the two series live in different normalisations).
\end{enumerate}
\end{proposition}
\begin{leancode}
@[simp] theorem mckayThompson_coeff (g : AutVnat) (n : ℕ) :
    PowerSeries.coeff n (mckayThompson g) =
      gradedAutTrace VnatVOA g ((n : ℤ) + 1)

theorem mckayThompson_coeff_zero (g : AutVnat) :
    PowerSeries.coeff 0 (mckayThompson g) = 0

theorem mckayThompson_conjugate (g h : AutVnat) :
    mckayThompson (h * g * h⁻¹) = mckayThompson g

theorem mckayThompson_one :
    mckayThompson 1 = PowerSeries.mk (fun n => PowerSeries.coeff (n + 1) JShift)

theorem mckayThompsonShift_one : mckayThompsonShift 1 = JShift

theorem mckayThompson_one_ne_JShift : mckayThompson 1 ≠ JShift
\end{leancode}
(\leanpath{VOA/Moonshine/MoonshineCNTrace.lean}.)
\begin{dictionary}
\lean{PowerSeries.coeff n F} is the coefficient of $X^n$ in $F$.  \lean{JShift} is the library's formal
$qJ=1+196884q^2+21493760q^3+\cdots$ (Chapter~\ref{ch:05}); its coefficient of $q^{n+1}$ is $c(n)$, the coefficient
of $q^n$ in $J$.  So (d) says: $\tr(1|\Vnat_{n+1})=\dim\Vnat_{n+1}=c(n)$ for all $n\ge0$, which is the graded
dimension theorem of Chapter~\ref{ch:14} in the present normalisation.  (e) is a deliberate ``negative'' lemma: the
constant coefficients are $0$ and $1$.  It records the index shift as a theorem so that it cannot be silently
confused.  All statements are the textbook ones.
\end{dictionary}
\paragraph{Proof on paper.}  (a) is the definition.  (b) $\Vnat_1=0$, so $g|_{\Vnat_1}$ is the zero map on the zero
space.  (c) Proposition~\ref{ch16:prop:traces}(b) coefficientwise.  (d) $\tr(1|\Vnat_{n+1})=\dim\Vnat_{n+1}$, and
$\dim\Vnat_m=[q^m]\,\mathtt{JShift}$ for all $m\ge0$ (\lean{finrank_gradingNat_eq_coeff_JShift},
Chapter~\ref{ch:05}).  (e) The constant coefficients are $0$ (by (b)) and $1$. \qed
\begin{physics}
(b) says the Monster CFT has no weight-one states: no currents, so the twisted partition functions have no constant
term.  This is why the Hauptmoduln are \emph{normalised} with constant term $0$.
\end{physics}
\begin{status}
Unconditional; all six printed in \leanpath{VOA/Moonshine/MoonshineCNAudit.lean}, standard axioms.
\end{status}

\subsection{Trace bounds}

To evaluate $T_g$ at a point of $\chsixteenH$ one needs the series to converge.  The trace of a finite-order linear
map is bounded by the dimension, and $J$ converges on $\chsixteenH$; that is the whole argument.

\begin{theorem}[Trace bounds]\label{ch16:thm:tracebounds}
\begin{enumerate}[label=(\alph*)]
\item If $f$ is a linear endomorphism of finite order of a finite-dimensional complex vector space $M$, then
  $|\tr f|\le\dim M$.
\item If $g$ is an automorphism of finite order of a vertex operator algebra $V$, then $|\tr(g|V_n)|\le\dim V_n$ for
  all $n$.
\item If $g\in\Aut(\Vnat)$ has finite order, then $|[X^n]\,\mathrm{MT}_g|\le|[X^n]\,\mathrm{MT}_1|=c(n)$ for all
  $n\ge0$.
\end{enumerate}
\end{theorem}
\begin{leancode}
theorem finiteOrderTrace_norm_le {M : Type*} [AddCommGroup M] [Module ℂ M]
    [FiniteDimensional ℂ M] (f : Module.End ℂ M) (hf : IsOfFinOrder f) :
    ‖LinearMap.trace ℂ M f‖ ≤ Module.finrank ℂ M

theorem gradedAutTrace_norm_le {M : Type*} [AddCommGroup M] [Module ℂ M]
    (V : VertexOperatorAlgebra M) (g : Aut V) (hg : IsOfFinOrder g) (n : ℤ) :
    ‖gradedAutTrace V g n‖ ≤ Module.finrank ℂ (V.grading n)

theorem mckayThompson_coeff_norm_le (g : AutVnat) (hg : IsOfFinOrder g) (n : ℕ) :
    ‖PowerSeries.coeff n (mckayThompson g)‖ ≤
      ‖PowerSeries.coeff n (mckayThompson 1)‖
\end{leancode}
(\leanpath{VOA/Moonshine/MoonshineCNTraceBounds.lean}.)
\begin{dictionary}
\lean{IsOfFinOrder f} (Mathlib) means $f^N=1$ for some $N\ge1$.  \lean{‖z‖} is the absolute value of a complex number.
\lean{[FiniteDimensional ℂ M]} is the hypothesis $\dim M<\infty$.  In (b) the finite-dimensionality of $V_n$ comes from
the vertex operator algebra axioms.  In (c) the right-hand side is $|\tr(1|\Vnat_{n+1})|=\dim\Vnat_{n+1}$.  The
hypothesis ``finite order'' is essential (a non-unitary automorphism of infinite order could have large traces), and
it is discharged for \emph{every} $g\in\Aut(\Vnat)$ by \lean{autVnat_isOfFinOrder} in Chapter~\ref{ch:20}.  The file
docstring stresses that no invariant Hermitian metric is used.
\end{dictionary}
\paragraph{Proof on paper.}  (a) If $f^N=1$ then every eigenvalue $\lambda$ satisfies $\lambda^N=1$, so $|\lambda|=1$.
Over $\C$ the characteristic polynomial splits and $\tr f$ is the sum of its $\dim M$ roots (with multiplicity); the
triangle inequality gives $|\tr f|\le\dim M$.  (b) Apply (a) to $g|_{V_n}$, which has finite order because
restriction is a homomorphism.  (c) Apply (b) to $\Vnat_{n+1}$ and use Proposition~\ref{ch16:prop:traces}(a). \qed
\begin{physics}
A finite symmetry acts on each energy level by a unitary matrix in a suitable basis, so twisting can only reduce the
magnitude of the partition-function coefficients: $|\tr(g|V_n)|\le\dim V_n$.  In particular $T_g$ converges wherever
the untwisted partition function $J$ does.
\end{physics}
\begin{status}
Unconditional (with the finite-order hypothesis explicit); printed in \leanpath{VOA/Moonshine/MoonshineCNAudit3.lean},
standard axioms.
\end{status}

\subsection{The McKay--Thompson series as a function on the upper half-plane}

\begin{definition}[McKay--Thompson function]\label{ch16:def:mtfun}
For $g\in\Aut(\Vnat)$ and $\tau\in\chsixteenH$, with $q=e^{2\pi i\tau}$,
\begin{equation*}
  T_g(\tau)\;=\;q^{-1}+\sum_{n\ge0}\tr\bigl(g\,|\,\Vnat_{n+1}\bigr)\,q^n
  \;=\;q^{-1}+\sum_{n\ge1}\tr\bigl(g\,|\,\Vnat_{n+1}\bigr)\,q^n .
\end{equation*}
\end{definition}
\begin{leancode}
noncomputable def mckayThompsonFun (g : AutVnat) (τ : ℍ) : ℂ :=
  (qPar τ)⁻¹ + ∑' n : ℕ, PowerSeries.coeff n (mckayThompson g) * qPar τ ^ n
\end{leancode}
(\leanpath{VOA/Moonshine/MoonshineCNStatement.lean}.)
\begin{dictionary}
\lean{(qPar τ)⁻¹} is $q^{-1}$.  \lean{∑' n : ℕ, a n} is Mathlib's \emph{unconditional infinite sum} (\lean{tsum}):
the sum of the family when it is summable, and \textbf{$0$ by convention when it is not} (a junk value,
Chapter~\ref{ch:01}, \S\ref{ch01:sec:junk}).  The summand is $\tr(g|\Vnat_{n+1})\,q^n$ by
Proposition~\ref{ch16:prop:mtprops}(a).  The $n=0$ term is $\tr(g|\Vnat_1)=0$, so the two displayed formulas agree;
\lean{mckayThompsonFun_eq_qSeriesFun} below makes this precise in Lean.

\textbf{Does the junk value matter?}  No, for two independent reasons.  (i) For $g$ of finite order the series is
absolutely summable at every $\tau\in\chsixteenH$: by Theorem~\ref{ch16:thm:tracebounds}(c) it is dominated by
$\sum_n c(n)|q|^n$, which converges because the $q$-expansion of $J$ converges on the open unit disc
(\lean{hasSum_JShift_laurent}, Chapter~\ref{ch:05}); and every $g$ has finite order (Chapter~\ref{ch:20}).  (ii) Even
without that, the junk value could only make the conjecture \emph{false}, never trivially true: where the sum is not
summable the function would equal $q^{-1}$, and $q^{-1}$ is not invariant under any $\chsixteenmat10N1$, $N\ge1$,
which every moonshine group contains (it is in $\Gamma_0(N)$).
\end{dictionary}
\begin{physics}
$T_g(\tau)=\tr_{\Vnat}\bigl(g\,q^{L_0-1}\bigr)$, the twisted torus partition function, as an honest function of the
modulus $\tau$.
\end{physics}
\begin{status}
Definition; printed in \leanpath{VOA/Moonshine/MoonshineCNAudit3.lean}, standard axioms.
\end{status}

\begin{theorem}[$T_1=J$]\label{ch16:thm:T1}
For every $\tau\in\chsixteenH$, $T_1(\tau)=j(\tau)-744$, where $j=E_4^3/\Delta$ is the analytic $j$-invariant.
\end{theorem}
\begin{leancode}
theorem mckayThompsonFun_one :
    mckayThompsonFun 1 = fun τ => jInvariant τ - 744
\end{leancode}
(\leanpath{VOA/Moonshine/MoonshineCNStatement.lean}.)
\begin{dictionary}
\lean{mckayThompsonFun 1} is $T_1$ for the identity automorphism \lean{1 : AutVnat}.  \lean{jInvariant} is the
analytic $j(\tau)=E_4(\tau)^3/\Delta(\tau)$ of Chapter~\ref{ch:L14}, Definition~\ref{l14:def:jinv} (file
\leanpath{VOA/Modularity/JInvariantAnalytic.lean}).  The equation is an equality of functions $\chsixteenH\to\C$.
This is the analytic counterpart of Proposition~\ref{ch16:prop:mtprops}(d); it is the bridge between the formal
graded dimension of Chapters~\ref{ch:05}--\ref{ch:14} and the analytic objects of the conjecture.
\end{dictionary}
\paragraph{Proof on paper.}  By Proposition~\ref{ch16:prop:mtprops}(d), the summand of $T_1$ is $c(n)q^n$ with
$c(n)=[q^{n+1}]\,\mathtt{JShift}$.  The theorem \lean{hasSum_JShift_laurent} (Chapter~\ref{ch:05}) says that
$\sum_{n\ge0}[q^{n+1}]\mathtt{JShift}\cdot q^n$ converges to $j(\tau)-744-q^{-1}$; it is proved from the
$q$-expansion of $E_4^3/\Delta$.  Adding $q^{-1}$ gives the claim. \qed
\begin{physics}
The untwisted partition function of the Monster CFT is $J=j-744$: the unique modular invariant function with the
vacuum pole $q^{-1}$ and no constant term (no currents).
\end{physics}
\begin{status}
Unconditional; printed in \leanpath{VOA/Moonshine/MoonshineCNAudit3.lean}, standard axioms.
\end{status}

\begin{example}[Two McKay--Thompson series, on paper]\label{ch16:ex:T2}
Under any isomorphism of $\Aut(\Vnat)$ with the Monster, the classes $2A$ and $2B$ give \eqref{ch16:eq:T2A} and
\eqref{ch16:eq:T2B}.  For instance, $\Vnat_2$ decomposes as $1\oplus196883$, so $\tr(g|\Vnat_2)=1+\chi_{196883}(g)$,
which is $4372$ for $2A$ and $276$ for $2B$.  For $2B$ there is a concrete representative, FLM's involution $z$,
which acts by $+1$ on $V_\Leech^+$ and by $-1$ on $(V_\Leech^T)^+$ (it is an automorphism because the block vertex
map respects the $\Z/2$-grading by sectors).  Its trace on $\Vnat_2$ is
$\dim(V_\Leech^+)_2-\dim(V_\Leech^T)^+_2=98580-98304=276$ (Chapter~\ref{ch:05} gives both dimensions).  The library
does not state class names; these examples are for orientation only.
\end{example}


%======================================================================
\section{Hauptmoduln and genus zero, without Riemann surfaces}\label{ch16:sec:haupt}
%======================================================================

\subsection{Modular curves, cusps and genus: the classical picture}

Let $\Gamma\le\chsixteenSLR$ be a subgroup commensurable with $\chsixteenSL$.  Such a group is discrete and acts
properly discontinuously on $\chsixteenH$.  The quotient $Y(\Gamma)=\Gamma\backslash\chsixteenH$ is a non-compact
Riemann surface: at a point $\tau_0$ whose stabiliser in $\Gamma/\{\pm1\}$ has order $e>1$ (an \emph{elliptic point})
the local coordinate is $\bigl((\tau-\tau_0)/(\tau-\bar\tau_0)\bigr)^e$.  The \emph{cusps} of $\Gamma$ are the points of
$\chsixteenP(\chsixteenR)$ fixed by parabolic elements of $\Gamma$; for our groups they are the points of $\chsixteenP(\Q)$, and
there are finitely many $\Gamma$-orbits of them.  Adding one point for each orbit gives a compact Riemann surface
\begin{equation*}
  \chsixteenX=\Gamma\backslash\bigl(\chsixteenH\cup\chsixteenP(\Q)\bigr).
\end{equation*}
Near the cusp $\infty$, if the stabiliser $\Gamma_\infty$ contains $\pm\chsixteenmat1h01$ with $h>0$ minimal (the
\emph{width} $h$), the local coordinate is $e^{2\pi i\tau/h}$; at another cusp $s=\gamma\infty$ one first moves $s$ to
$\infty$ with $\gamma^{-1}$.  The \emph{genus} of $\Gamma$ is the genus of $\chsixteenX$.

A \emph{modular function} for $\Gamma$ is a meromorphic function on $\chsixteenX$; equivalently, a $\Gamma$-invariant
meromorphic function on $\chsixteenH$ that is meromorphic in the local coordinate at every cusp.  If $\chsixteenX$
has genus zero, there is an isomorphism of Riemann surfaces $\chsixteenX\to\chsixteenP(\C)$, and it can be chosen to
send the cusp $\infty$ to $\infty$.  Such an isomorphism is a \emph{Hauptmodul}; every modular function for $\Gamma$ is
a rational function of it.  If the width of $\infty$ is $1$, the Hauptmodul has a simple pole $a\,q^{-1}$ at $\infty$,
and replacing $f$ by $(f-b)/a$ we may arrange $f=q^{-1}+0+O(q)$; with this normalisation it is unique (two
Hauptmoduln differ by an automorphism of $\chsixteenP(\C)$ fixing $\infty$, i.e.\ by $z\mapsto az+b$, and the
expansion $q^{-1}+0+\cdots$ forces $a=1$, $b=0$), the \emph{normalised Hauptmodul}.  For all these classical facts
(compactness of $\chsixteenX$, local coordinates at elliptic points and cusps, discreteness of groups commensurable
with $\chsixteenSL$, cusps $=\chsixteenP(\Q)$) see Diamond--Shurman, Chapters~2 and~3.

\begin{example}[Genus-zero and not genus-zero groups]\label{ch16:ex:genus}
$\chsixteenSL$ has genus zero, with normalised Hauptmodul $J=j-744$.  $\Gamma_0(N)$ has genus zero exactly for
$N\in\{1,\dots,10,12,13,16,18,25\}$; for instance the Hauptmodul of $\Gamma_0(2)$ is $(\eta(\tau)/\eta(2\tau))^{24}+24$.
$\Gamma_0(11)$ has genus one ($X_0(11)$ is an elliptic curve), so it has no Hauptmodul; but the larger group
$\Gamma_0(11)+$ obtained by adjoining the Fricke involution $\tau\mapsto-1/(11\tau)$ has genus zero, and its normalised
Hauptmodul is $T_{11A}$.  Ogg's observation: $\Gamma_0(p)+$ has genus zero exactly for the fifteen primes dividing
$|\chsixteenM|$.  Note that $\Gamma_0(p)+$ is not contained in $\chsixteenSL$: the Fricke involution is
$\frac1{\sqrt p}\chsixteenmat0{-1}p0\in\chsixteenSLR$, with irrational entries.  This is why the conjecture is
formulated in $\chsixteenSLR$, and the library works in $\chsixteenGLR$.
\end{example}

The library's Mathlib (v4.30) has congruence subgroups, cusps and the action of $\chsixteenGLR$, but no compactified
modular curves as Riemann surfaces and no notion of genus.  The statement file therefore encodes ``normalised
Hauptmodul of a genus-zero group'' \emph{without Riemann surfaces}: it asks that $f$ induce a \emph{bijection}
$\chsixteenX\to\chsixteenP(\C)$ with the right behaviour at the cusps, and leaves to the mathematics (below) the
observation that such a bijection is automatically an isomorphism of Riemann surfaces.

\subsection{Values at cusps}

\begin{definition}[Cusp data]\label{ch16:def:cuspdata}
Let $c\in\chsixteenP(\chsixteenR)$, $\Gamma\le\chsixteenGLR$, $f:\chsixteenH\to\C$ and $w\in\C$.
\begin{enumerate}[label=(\alph*)]
\item $\gamma\in\chsixteenGLR$ \emph{sends $\infty$ to $c$} if $\det\gamma>0$ and $\gamma\cdot\infty=c$.
\item $f$ \emph{has value $w$ at the cusp $c$} if for every $\gamma$ sending $\infty$ to $c$,
  $f(\gamma\tau)\to w$ as $\operatorname{Im}\tau\to\infty$, uniformly in $\operatorname{Re}\tau$.
\item $c$ is an \emph{other cusp} of $\Gamma$ if it is a cusp of $\Gamma$ and is not in the $\Gamma$-orbit of $\infty$.
\end{enumerate}
\end{definition}
\begin{leancode}
def SendsInftyTo (γ : GL (Fin 2) ℝ) (c : OnePoint ℝ) : Prop :=
  0 < (γ.det : ℝ) ∧ γ • cuspInfty = c

def HasCuspValue (f : ℍ → ℂ) (c : OnePoint ℝ) (w : ℂ) : Prop :=
  ∀ γ : GL (Fin 2) ℝ, SendsInftyTo γ c → Tendsto (fun τ : ℍ => f (γ • τ)) atImInfty (𝓝 w)

def IsOtherCusp (Γ : Subgroup (GL (Fin 2) ℝ)) (c : OnePoint ℝ) : Prop :=
  IsCusp c Γ ∧ ∀ γ ∈ Γ, γ • cuspInfty ≠ c
\end{leancode}
(\leanpath{VOA/Moonshine/MoonshineCNStatement.lean}.)
\begin{dictionary}
\lean{(γ.det : ℝ)} is $\det\gamma$; \lean{γ • cuspInfty} is $\gamma\cdot\infty=a/c$ (or $\infty$ if $c=0$).
\lean{Tendsto (fun τ => f (γ • τ)) atImInfty (𝓝 w)} is ``$f(\gamma\tau)\to w$ as $\operatorname{Im}\tau\to\infty$,
uniformly in $\operatorname{Re}\tau$'' (\S\ref{ch16:sec:analytic}).  \lean{IsCusp c Γ} is Mathlib's: $c$ is fixed by a
parabolic element of $\Gamma$.

\textbf{``For every $\gamma$'' versus ``for some $\gamma$''.}  Two matrices $\gamma,\gamma'$ of positive determinant
sending $\infty$ to $c$ differ by an element fixing $\infty$: $\gamma'=\gamma\beta$ with
$\beta\tau=a\tau+b$, $a>0$.  Since $\operatorname{Im}(a\tau+b)\to\infty$ uniformly exactly when
$\operatorname{Im}\tau\to\infty$, the limit along $\gamma$ exists iff it exists along $\gamma'$, with the same value.
So the universal quantifier in \lean{HasCuspValue} is equivalent to an existential one; it simply avoids choosing a
matrix.  The condition is the usual one: $f$ is bounded near $c$ and extends continuously (indeed, for invariant
holomorphic $f$, holomorphically in the cusp coordinate, by Riemann's removable singularity theorem) with value $w$.
\end{dictionary}
\begin{physics}
A cusp of $\Gamma$ is a degeneration limit of the torus, seen through a modular transformation in $\Gamma$; the other
cusps correspond to long tori in which a twisted sector propagates.  ``$f$ has a finite value at the other cusps''
says that only the cusp $\infty$ carries the vacuum pole $q^{-1}$: in the other degeneration limits the partition
function is dominated by states of non-negative effective energy.
\end{physics}
\begin{status}
Definitions; printed in \leanpath{VOA/Moonshine/MoonshineCNAudit3.lean}, standard axioms.
\end{status}

\subsection{The normalised Hauptmodul}

\begin{definition}[Normalised Hauptmodul]\label{ch16:def:hauptmodul}
Let $\Gamma\le\chsixteenGLR$ and $f:\chsixteenH\to\C$.  $f$ is a \emph{normalised Hauptmodul for $\Gamma$} if:
\begin{enumerate}[label=(H\arabic*)]
\item\label{ch16:H1} $f$ is holomorphic on $\chsixteenH$ \emph{(no poles in the interior)};
\item\label{ch16:H2} $f(\gamma\tau)=f(\tau)$ for all $\gamma\in\Gamma$, $\tau\in\chsixteenH$ \emph{($f$ descends to
  $\Gamma\backslash\chsixteenH$)};
\item\label{ch16:H3} $f(\tau)-q^{-1}\to0$ as $\operatorname{Im}\tau\to\infty$ \emph{(a simple pole at the cusp $\infty$,
  normalised $q^{-1}+0+O(q)$)};
\item\label{ch16:H4} $f$ has a (finite) value at every other cusp of $\Gamma$ \emph{(no poles at the other cusps:
  so $F$ below is well defined and takes the value $\infty$ only at the cusp $\infty$)};
\item\label{ch16:H5} if $f(\tau)=f(\sigma)$ then $\sigma=\gamma\tau$ for some $\gamma\in\Gamma$ \emph{($F$ is injective
  on interior points)};
\item\label{ch16:H6} a value of $f$ at an other cusp is not attained on $\chsixteenH$ \emph{($F$ never gives an
  interior point and a cusp the same value)};
\item\label{ch16:H7} if two other cusps $c,c'$ have the same value, then $c'=\gamma c$ for some $\gamma\in\Gamma$
  \emph{($F$ is injective on cusps)};
\item\label{ch16:H8} every $w\in\C$ is attained on $\chsixteenH$ or is the value of $f$ at an other cusp \emph{($F$
  is surjective)}.
\end{enumerate}
Here $F:\chsixteenX\to\chsixteenP(\C)$ is the map induced by $f$: $F(\Gamma\tau)=f(\tau)$ on interior points, $F$ of the
cusp $\infty$ is $\infty$, and $F$ of another cusp is the value of $f$ there.  So \ref{ch16:H1}--\ref{ch16:H4} say that
$f$ is a modular function for $\Gamma$ with a single, simple, normalised pole at $\infty$; \ref{ch16:H5}--\ref{ch16:H7}
say that $F$ is injective (on interior points, between interior points and cusps, and on cusps); \ref{ch16:H8} says $F$
is surjective.  Example~\ref{ch16:ex:gamma02} checks all eight clauses for $\Gamma_0(2)$.
A subgroup is \emph{of genus zero} (with a cusp of width one at $\infty$) if it has a normalised Hauptmodul.
\end{definition}
\begin{leancode}
structure IsNormalizedHauptmodul (Γ : Subgroup (GL (Fin 2) ℝ)) (f : ℍ → ℂ) : Prop where
  holomorphic : MDifferentiable 𝓘(ℂ) 𝓘(ℂ) f
  invariant : ∀ γ ∈ Γ, ∀ τ : ℍ, f (γ • τ) = f τ
  normalized : Tendsto (fun τ : ℍ => f τ - (qPar τ)⁻¹) atImInfty (𝓝 0)
  cuspFinite : ∀ c, IsOtherCusp Γ c → ∃ w, HasCuspValue f c w
  injective : ∀ τ σ : ℍ, f τ = f σ → ∃ γ ∈ Γ, γ • τ = σ
  cuspValue_not_attained :
    ∀ c w, IsOtherCusp Γ c → HasCuspValue f c w → ∀ τ : ℍ, f τ ≠ w
  cuspValue_injective : ∀ c c' w, IsOtherCusp Γ c → IsOtherCusp Γ c' →
    HasCuspValue f c w → HasCuspValue f c' w → ∃ γ ∈ Γ, γ • c = c'
  surjective : ∀ w : ℂ, (∃ τ : ℍ, f τ = w) ∨ ∃ c, IsOtherCusp Γ c ∧ HasCuspValue f c w

def IsGenusZero (Γ : Subgroup (GL (Fin 2) ℝ)) : Prop :=
  ∃ f : ℍ → ℂ, IsNormalizedHauptmodul Γ f
\end{leancode}
(\leanpath{VOA/Moonshine/MoonshineCNStatement.lean}.)
\begin{dictionary}
\textbf{Field by field.}  A \lean{structure ... : Prop} is a conjunction of named clauses
(Chapter~\ref{ch:01}, \S\ref{ch01:sec:structures}).
\begin{itemize}
\item \lean{holomorphic} is \ref{ch16:H1}.
\item \lean{invariant} is \ref{ch16:H2}, with $\gamma$ acting by \lean{γ • τ}.
\item \lean{normalized} is \ref{ch16:H3}: \lean{f τ - (qPar τ)⁻¹ → 0} along \lean{atImInfty}.  For a $1$-periodic
  holomorphic $f$ (which is the case in the conjecture, see below) this says exactly that the Fourier expansion of $f$
  is $q^{-1}+\sum_{n\ge1}a_nq^n$: a simple pole with residue normalised to $1$, \emph{zero constant term}, and no other
  negative powers.
\item \lean{cuspFinite} is \ref{ch16:H4}.
\item \lean{injective}, \lean{cuspValue_not_attained}, \lean{cuspValue_injective} are \ref{ch16:H5}--\ref{ch16:H7}:
  together they say that the induced map $F:\chsixteenX\to\chsixteenP(\C)$, defined by $F(\Gamma\tau)=f(\tau)$,
  $F(\Gamma\infty)=\infty$ and $F(\Gamma c)=$ the value at $c$, is \emph{injective}.  (Points of $\chsixteenH$ never go
  to $\infty$ because $f$ is $\C$-valued, and the other cusps have finite values by \ref{ch16:H4}; so the only point
  mapped to $\infty$ is the cusp $\infty$.)
\item \lean{surjective} is \ref{ch16:H8}: $F$ is \emph{surjective} ($\infty$ is hit by the cusp $\infty$).
\end{itemize}
\lean{IsGenusZero Γ} is \lean{∃ f, IsNormalizedHauptmodul Γ f}.  It is \emph{not} used in the statement of the
conjecture; it records the terminology.

\textbf{Same as the textbook notion?}  For $\Gamma\le\chsixteenSLR$ commensurable with $\chsixteenSL$ and containing
the translation $\tau\mapsto\tau+1$ (both hold in the conjecture), \lean{IsNormalizedHauptmodul Γ f} holds if and only
if $\chsixteenX$ has genus zero, the width of $\infty$ is $1$, and $f$ is its normalised Hauptmodul.  The proof is
Proposition~\ref{ch16:prop:haupt} below.  It is a paper argument: the library contains neither Riemann surfaces nor
genus, so this equivalence is a statement \emph{about} the definition, not a Lean theorem.  This is the central
translation step of the chapter, and the referee should check it.

\textbf{How the definition could have been wrong, and why it is not.}
\begin{itemize}
\item \emph{Too weak, by choosing $\Gamma$ too small?}  Field \lean{injective} prevents this: the fibres of $f$ on
  $\chsixteenH$ must be single $\Gamma$-orbits.  If $f(\delta\tau)=f(\tau)$ for some $\delta\in\chsixteenSLR$ and all
  $\tau$, then for each $\tau$ some $\gamma_\tau\in\Gamma$ has $\gamma_\tau\tau=\delta\tau$.  When $\Gamma$ is countable (as it is in the
  conjecture, being commensurable with $\chsixteenSL$), one $\gamma$ works for uncountably many $\tau$; then
  $\gamma^{-1}\delta\in\chsixteenSLR$ fixes at least two points of $\chsixteenH$, and a non-scalar M\"obius map fixes
  at most one, so $\gamma^{-1}\delta=\pm1$.  Hence $\pm\Gamma$ is the \emph{full} invariance group of $f$ in
  $\chsixteenSLR$: the group is determined by the function, up to $\pm1$.
\item \emph{Too weak, by allowing a pole of higher order, or poles at other cusps, or in $\chsixteenH$?}  Excluded by
  \lean{normalized}, \lean{cuspFinite} and $f$ being $\C$-valued and holomorphic.
\item \emph{Too weak, by allowing $F$ to be of higher degree?}  Excluded by injectivity: a degree-$d$ map with $d\ge2$
  identifies $d$ points of a generic fibre.
\item \emph{Too strong, by demanding something false of $J$?}  No: $J$ satisfies all eight clauses for
  $\Gamma=\chsixteenSL$, by Proposition~\ref{ch16:prop:haupt} and the classical fact that $X(\chsixteenSL)$ has genus
  zero with Hauptmodul $J$ (Chapter~\ref{ch:05}, Theorem~\ref{ch05:thm:haupt}).
\item \emph{Vacuous because cusps are mis-specified?}  The cusps are Mathlib's parabolic fixed points; for groups
  commensurable with $\chsixteenSL$ they are $\chsixteenP(\Q)$, as they should be.
\end{itemize}
\end{dictionary}

\begin{physics}
A normalised Hauptmodul is the partition function of a chiral theory whose symmetry group is $\Gamma$. Being a
Hauptmodul means it is a single coordinate on the moduli space $\Gamma\backslash\mathbb H^{*}$ of tori with
$\Gamma$-structure. Normalisation $q^{-1}+0+O(q)$ says the theory has a unique ground state of energy $-c/24=-1$ and
no weight-one currents. In moonshine, $T_g$ is the torus partition function of $\Vnat$ twisted by the symmetry $g$ in
the time direction. The genus-zero property says this twisted partition function alone determines the whole modular
curve on which it lives.
\end{physics}

\begin{proposition}[What the eight fields mean; on paper]\label{ch16:prop:haupt}
Let $\Gamma\le\chsixteenSLR$ be commensurable with $\chsixteenSL$ and contain $\chsixteenmat1101$.
\begin{enumerate}[label=(\alph*)]
\item If $f$ satisfies \ref{ch16:H1}--\ref{ch16:H8}, then the width of $\infty$ is $1$, the map
  $F:\chsixteenX\to\chsixteenP(\C)$ is an isomorphism of Riemann surfaces, $\chsixteenX$ has genus zero, and $f$ is
  the normalised Hauptmodul of $\Gamma$.
\item Conversely, if $\chsixteenX$ has genus zero and $\infty$ has width $1$, the normalised Hauptmodul satisfies
  \ref{ch16:H1}--\ref{ch16:H8}.
\end{enumerate}
\end{proposition}
\paragraph{Proof on paper.}  (a) \emph{Width.}  If $\Gamma$ contained $\tau\mapsto\tau+h$ with $0<h<1$, then
$f(\tau+h)=f(\tau)$ while $f(\tau)=e^{-2\pi i\tau}+o(1)$; comparing, $(e^{-2\pi ih}-1)e^{-2\pi i\tau}\to0$, which
forces $e^{2\pi ih}=1$, impossible.  So the width is $1$ and $q$ is the local coordinate at $\infty$.
\emph{Holomorphy of $F$.}  On $Y(\Gamma)$: $f$ is holomorphic and invariant, hence a holomorphic function of the local
coordinate (at an elliptic point $\tau_0$ of order $e$, an invariant holomorphic function is a power series in
$(\tau-\tau_0)^e/(\tau-\bar\tau_0)^e$).  At $\infty$: by \ref{ch16:H2} and periodicity $f=\sum_n a_nq^n$, and by
\ref{ch16:H3} $f=q^{-1}+O(q)$, so $1/F=q+O(q^2)$ is holomorphic with a simple zero.  At another cusp $c=\gamma\infty$:
$f\circ\gamma$ is invariant under the (conjugated) stabiliser and tends to $w$ by \ref{ch16:H4}, so by Riemann's
removable singularity theorem it is holomorphic in the cusp coordinate with value $w$.  So $F$ is a holomorphic map of
compact Riemann surfaces.  \emph{Bijectivity.}  Injectivity was explained in the dictionary; surjectivity is
\ref{ch16:H8} together with $F(\infty)=\infty$.  A bijective holomorphic map between Riemann surfaces is a
biholomorphism (an injective holomorphic map has nowhere vanishing derivative in local coordinates).  Hence
$\chsixteenX\cong\chsixteenP(\C)$ has genus zero, and $F$ is a Hauptmodul with the normalisation $q^{-1}+0+O(q)$.

(b) Let $F:\chsixteenX\to\chsixteenP(\C)$ be an isomorphism with $F(\infty)=\infty$, normalised to $q^{-1}+0+O(q)$
(possible because the width is $1$), and $f=F\circ\pi$ on $\chsixteenH$.  Then \ref{ch16:H1}--\ref{ch16:H4} hold
because $F$ is holomorphic with its only pole at $\infty$; \ref{ch16:H5}--\ref{ch16:H8} hold because $F$ is
bijective. \qed

\begin{example}[The eight clauses for $\Gamma_0(2)$; on paper]\label{ch16:ex:gamma02}
Let $\Gamma=\Gamma_0(2)$ and $f=t+24$, $t(\tau)=(\eta(\tau)/\eta(2\tau))^{24}=\Delta(\tau)/\Delta(2\tau)$, so that
$f=T_{2B}$ of \eqref{ch16:eq:T2B}.
\emph{The quotient.}  $\Gamma_0(2)$ has index $3$ in $\chsixteenSL$.  A cusp $a/c$ (in lowest terms; $\infty=1/0$) is
$\Gamma_0(2)$-equivalent to $\infty$ if $c$ is even and to $0$ if $c$ is odd, so there are two cusp orbits: $\infty$, of
width $1$ (the stabiliser contains $\tau\mapsto\tau+1$), and $0$, of width $2$ (its stabiliser is conjugated by
$S=\chsixteenmat0{-1}10$ to the translations by multiples of $2$).  There is one elliptic point, $(1+i)/2$, of
order~$2$.  The other cusps in the sense of Definition~\ref{ch16:def:cuspdata} are exactly the rationals $a/c$ with
$c$ odd.
\begin{itemize}
\item \ref{ch16:H1}: $\eta$ has no zeros on $\chsixteenH$, so $t$ is holomorphic.
\item \ref{ch16:H2}: $\Delta(\tau)$ and $\Delta(2\tau)$ are both modular forms of weight $12$ for $\Gamma_0(2)$, so
  their quotient is invariant.
\item \ref{ch16:H3}: $t=q^{-1}\prod_{n\ge1}(1+q^n)^{-24}=q^{-1}-24+276q+\cdots$, so $f=q^{-1}+0+276q+\cdots$.
\item \ref{ch16:H4}: the cusp $0$ is $S\infty$ ($\det S=1>0$).  From $\Delta(-1/\tau)=\tau^{12}\Delta(\tau)$ one gets
  $t(S\tau)=2^{12}\Delta(\tau)/\Delta(\tau/2)=2^{12}q^{1/2}(1+O(q^{1/2}))\to0$, so $f$ has the value $24$ at the cusp
  $0$, and hence (by invariance) at every other cusp.  The half-integral power $q^{1/2}$ is the width-$2$ coordinate.
\item \ref{ch16:H6}: $f(\tau)=24$ would mean $\Delta(\tau)=0$, which never happens; so the cusp value $24$ is not
  attained on $\chsixteenH$.
\item \ref{ch16:H7}: there is only one orbit of other cusps, so any two other cusps are $\Gamma$-equivalent.
\item \ref{ch16:H5}, \ref{ch16:H8}: the induced $F:X_0(2)\to\chsixteenP(\C)$ is a meromorphic function on a compact
  Riemann surface whose only pole is the simple pole at $\infty$ (by \ref{ch16:H1}, \ref{ch16:H3}, \ref{ch16:H4}).  A
  non-constant meromorphic function takes every value in $\chsixteenP(\C)$ exactly as many times as it has poles,
  counted with multiplicity; here once.  So $F$ is bijective: every $w\ne24$ is attained at exactly one orbit in
  $\Gamma\backslash\chsixteenH$ (\ref{ch16:H5}, \ref{ch16:H8}), and $24$ is attained only at the cusp $0$.  (In
  particular $X_0(2)$ has genus zero.)
\end{itemize}
This is the pattern of every application: \ref{ch16:H1}--\ref{ch16:H4} are local checks, and the global clauses
\ref{ch16:H5}--\ref{ch16:H8} follow from ``exactly one simple pole''.  For the McKay--Thompson series no such
explicit formula is available in general, which is why the library proves the global clauses by the case-free
argument of Chapter~\ref{ch:19}.
\end{example}

\begin{physics}
Think of $\chsixteenX$ as a moduli space: the space of tori (of modulus $\tau$) carrying a fixed symmetry twist,
modulo those modular transformations in $\Gamma$ that preserve the twist.  Its cusps are the degeneration limits
(infinitely long tori).  ``Genus zero with Hauptmodul $f$'' says that this moduli space is a sphere and that the
twisted partition function $f$ is a global coordinate on it, with the vacuum pole $q^{-1}$ at the one degeneration
where the untwisted vacuum propagates and finite values at all others.  Consequently any other quantity of the theory
that is invariant under $\Gamma$, holomorphic on $\chsixteenH$ and has at most poles at the degeneration limits is a
rational function of $f$ (a polynomial if its only pole is at $\infty$); this is the twisted analogue of the
statement that a $c=24$ holomorphic partition function is $J+\text{const}$.  It is a rigid condition: twisted
partition functions of a general orbifold are invariant under congruence groups of positive genus.
\end{physics}
\begin{status}
\lean{IsNormalizedHauptmodul}, \lean{IsGenusZero}: definitions, printed in
\leanpath{VOA/Moonshine/MoonshineCNAudit3.lean}, standard axioms.  Proposition~\ref{ch16:prop:haupt}: argued on paper
only (it is about the meaning of the definition and needs Riemann surfaces, which the library does not have).
\end{status}

%======================================================================
\section{Moonshine groups and the Conway--Norton statement}\label{ch16:sec:statement}
%======================================================================

\subsection{Moonshine groups}

\begin{definition}[Moonshine group]\label{ch16:def:moonshinegroup}
For $N\in\N$, $\Gamma_0(N)$ and $\Gamma_1(N)$ are regarded as subgroups of $\chsixteenGLR$.  A subgroup
$\Gamma\le\chsixteenGLR$ is a \emph{moonshine group of level $N$} if $N\ge1$, $\Gamma\le\chsixteenSLR$,
$\Gamma_0(N)\le\Gamma$, and $\Gamma$ is commensurable with $\chsixteenSL$ (i.e.\ $\Gamma\cap\chsixteenSL$ has finite
index in both $\Gamma$ and $\chsixteenSL$).
\end{definition}
\noindent\textbf{Terminology.}  ``Moonshine group'' is this chapter's (and the library's) local name, taken from the Lean
definition \lean{IsMoonshineGroup}.  In the literature the phrase is used differently (for instance for the genus-zero
groups of the Conway--Norton list, or for the groups of ``moonshine type'' of Conway--McKay--Sebbar); a reader
consulting Gannon or Conway--Norton will not find this exact definition there.
\begin{leancode}
abbrev gamma0GL (N : ℕ) : Subgroup (GL (Fin 2) ℝ) := ((Gamma0 N : Subgroup SL(2, ℤ)) : Subgroup (GL (Fin 2) ℝ))
abbrev gamma1GL (N : ℕ) : Subgroup (GL (Fin 2) ℝ) := ((Gamma1 N : Subgroup SL(2, ℤ)) : Subgroup (GL (Fin 2) ℝ))

def IsMoonshineGroup (Γ : Subgroup (GL (Fin 2) ℝ)) (N : ℕ) : Prop :=
  0 < N ∧ Γ.HasDetOne ∧ gamma0GL N ≤ Γ ∧ Subgroup.Commensurable Γ 𝒮ℒ
\end{leancode}
(\leanpath{VOA/Moonshine/MoonshineCNStatement.lean}.)
\begin{dictionary}
\textbf{Decoding.}  \lean{((Gamma0 N : Subgroup SL(2, ℤ)) : Subgroup (GL (Fin 2) ℝ))} is the image of Mathlib's
$\Gamma_0(N)\le\chsixteenSL$ under \lean{mapGL ℝ} (\S\ref{ch16:sec:analytic}); likewise \lean{gamma1GL}.  In
\lean{IsMoonshineGroup Γ N}: \lean{0 < N} is $N\ge1$; \lean{Γ.HasDetOne} is $\Gamma\le\chsixteenSLR$;
\lean{gamma0GL N ≤ Γ} is $\Gamma_0(N)\subseteq\Gamma$; \lean{Subgroup.Commensurable Γ 𝒮ℒ} is commensurability with
$\chsixteenSL$.

\textbf{Each clause, and why it is the intended one.}
\begin{itemize}
\item \emph{Ambient group $\chsixteenGLR$, not $\chsixteenSL$}: necessary, since many Conway--Norton groups contain
  Atkin--Lehner and Fricke elements with irrational entries (Example~\ref{ch16:ex:genus}).
\item \emph{$\Gamma\le\chsixteenSLR$}: Conway--Norton's groups are subgroups of $\chsixteenSLR$ (or of
  $\mathrm{PSL}_2(\chsixteenR)$; note $-1\in\Gamma_0(N)\le\Gamma$ always, so nothing is lost by working with matrices rather
  than M\"obius maps).  Without this clause one could add orientation-reversing or scaling elements; with it, the
  action on $\chsixteenH$ is the usual one.
\item \emph{$\Gamma_0(N)\le\Gamma$ with $N\ge1$}: this is the ``congruence'' clause of the conjecture.  It forces
  $\Gamma$ to contain $\chsixteenmat1101$ (so $T_g$ is $1$-periodic, which it is anyway) and $\chsixteenmat10N1$
  (a non-trivial condition on $T_g$).  For $N=1$ it forces $\Gamma\supseteq\chsixteenSL$.
\item \emph{Commensurable with $\chsixteenSL$}: implies that $\Gamma$ is discrete, of finite covolume, with cusps
  $\chsixteenP(\Q)$; together with $\Gamma_0(N)\le\Gamma$ it says $[\Gamma:\Gamma_0(N)]<\infty$.
\end{itemize}
\textbf{Could the clauses be trivially satisfied?}  $\Gamma=\Gamma_0(1)=\chsixteenSL$, $N=1$ satisfies
\lean{IsMoonshineGroup}; so the definition is not vacuous.  It is also not trivially satisfiable \emph{together with}
\lean{IsNormalizedHauptmodul Γ f}: by the dictionary of Definition~\ref{ch16:def:hauptmodul}, $\pm\Gamma$ is the full
invariance group of $f$, so the content of \lean{IsMoonshineGroup Γ N} is ``the invariance group of $f$ contains
$\Gamma_0(N)$ and is commensurable with $\chsixteenSL$''.  For example $q^{-1}$ alone, or any $f$ invariant only
under translations, admits no moonshine group.

\textbf{Not required (deliberately).}  The definition does not ask that $\Gamma$ normalise $\Gamma_0(N)$, nor that $N$
be related to the order of $g$; see \S\ref{ch16:sec:fidelity}.
\end{dictionary}
\begin{physics}
Write $Z(h,g)$ for the torus partition function with commuting symmetry twists $h$ along the space cycle and $g$
along the time cycle, so that $T_g=Z(1,g)$.  A modular transformation $\chsixteenmat abcd$ relabels the two cycles
and sends $Z(h,g)$ to $Z(h^ag^c,h^bg^d)$ (up to a phase).  For $h=1$ and $g$ of order $N$, the matrices with
$c\equiv0\pmod N$, i.e.\ $\Gamma_0(N)$, keep the space cycle untwisted and send $T_g$ to $T_{g^d}$; those of
$\Gamma_1(N)$ ($d\equiv1$) send it back to $T_g$.  So invariance under $\Gamma_1(N)$ is the expected modularity of a
twisted partition function, and invariance under $\Gamma_0(N)$ adds the statement $T_{g^d}=T_g$.  (The phases are why
Conway and Norton's level is $o(g)h$ rather than $o(g)$.)  Commensurability with $\chsixteenSL$ says that $\Gamma$ is
built from finitely many such relabellings up to finite index; heuristically, this is the content of orbifold
modular invariance.
\end{physics}
\begin{status}
Definitions; \lean{gamma0GL}, \lean{gamma1GL}, \lean{IsMoonshineGroup} printed in
\leanpath{VOA/Moonshine/MoonshineCNAudit3.lean}, standard axioms.
\end{status}

\subsection{The conjecture}

\begin{definition}[The Conway--Norton property and the Conway--Norton conjecture]\label{ch16:def:cn}
$g\in\Aut(\Vnat)$ has the \emph{Conway--Norton property} if there are a subgroup $\Gamma\le\chsixteenGLR$ and
$N\in\N$ such that $\Gamma$ is a moonshine group of level $N$ and $T_g$ is a normalised Hauptmodul for $\Gamma$.  The
\emph{Conway--Norton conjecture} (for $\Aut(\Vnat)$) is: every $g\in\Aut(\Vnat)$ has the Conway--Norton property.  For a
group homomorphism $\iota:G\to\Aut(\Vnat)$, \emph{Conway--Norton via $\iota$} means: $\iota(g)$ has the Conway--Norton
property for every $g\in G$.
\end{definition}
\begin{leancode}
def ConwayNortonFor (g : AutVnat) : Prop :=
  ∃ (Γ : Subgroup (GL (Fin 2) ℝ)) (N : ℕ), IsMoonshineGroup Γ N ∧
    IsNormalizedHauptmodul Γ (mckayThompsonFun g)

def ConwayNortonConjecture : Prop :=
  ∀ g : AutVnat, ConwayNortonFor g

def ConwayNortonVia {𝕄 : Type*} [Group 𝕄] (ι : 𝕄 →* AutVnat) : Prop :=
  ∀ g : 𝕄, ConwayNortonFor (ι g)
\end{leancode}
(\leanpath{VOA/Moonshine/MoonshineCNStatement.lean}.)
\begin{dictionary}
Unfolded, \lean{ConwayNortonConjecture} is
\begin{align*}
  &\forall g\in\Aut(\Vnat)\ \ \exists\,\Gamma\le\chsixteenGLR\ \ \exists\,N\in\N:\\
  &\qquad N\ge1,\ \ \Gamma\le\chsixteenSLR,\ \ \Gamma_0(N)\le\Gamma,\ \ \Gamma\ \text{commensurable with}\ \chsixteenSL,\\
  &\qquad T_g\ \text{satisfies \ref{ch16:H1}--\ref{ch16:H8} for }\Gamma .
\end{align*}
The quantifier order is the right one: $\Gamma$ and $N$ depend on $g$.  \lean{ConwayNortonVia ι} is the same statement
for $\iota(g)$, $g\in G$; for any $\iota$ it follows immediately from \lean{ConwayNortonConjecture}, and it is the form
in which the Monster version is stated in Chapter~\ref{ch:22} ($G=\mathrm{AutO}$ and $\iota$ an isomorphism).  The
full comparison with Conway and Norton's conjecture is \S\ref{ch16:sec:fidelity}.
\end{dictionary}
\begin{physics}
For every global symmetry $g$ of the Monster CFT, the torus partition function with a $g$-twist is a global
coordinate on a genus-zero modular curve, with the vacuum pole at the cusp $\infty$ and finite values at all
other degenerations.
\end{physics}
\begin{status}
Definitions (propositions); printed in \leanpath{VOA/Moonshine/MoonshineCNAudit3.lean}, standard axioms.  The
proposition \lean{ConwayNortonConjecture} is \emph{proved}, with no hypotheses, as \lean{conwayNorton}
(\S\ref{ch16:sec:final}).  The statement file was fixed before the proof was attempted and has not been modified since.
\end{status}

\begin{definition}[The $\Gamma_1$ variant]\label{ch16:def:gammaone}
The \emph{$\Gamma_1$ form of the conjecture}: for every $g\in\Aut(\Vnat)$ \emph{of finite order} there are
$\Gamma\le\chsixteenSLR$ and $N\ge1$ with $\Gamma_1(N)\le\Gamma$, $\Gamma$ commensurable with $\chsixteenSL$, and $T_g$
a normalised Hauptmodul for $\Gamma$.
\end{definition}
\begin{leancode}
def ConwayNortonGammaOne : Prop :=
  ∀ g : AutVnat, IsOfFinOrder g → ∃ (Γ : Subgroup (GL (Fin 2) ℝ)) (N : ℕ),
    0 < N ∧ Γ.HasDetOne ∧ gamma1GL N ≤ Γ ∧ Subgroup.Commensurable Γ 𝒮ℒ ∧
      IsNormalizedHauptmodul Γ (mckayThompsonFun g)
\end{leancode}
(\leanpath{VOA/Moonshine/MoonshineCNStatement.lean}.)
\begin{dictionary}
\lean{IsOfFinOrder g}: $g^n=1$ for some $n\ge1$.  The clauses are those of \lean{IsMoonshineGroup}, with
\lean{gamma1GL N} instead of \lean{gamma0GL N}.  Since $\Gamma_1(N)\le\Gamma_0(N)$, the conclusion is weaker than that
of \lean{ConwayNortonFor}; the extra hypothesis of finite order makes the whole proposition weaker still.  It is the
natural output of the case-free genus-zero argument (Chapter~\ref{ch:19}) and the input of the bridge of
\S\ref{ch16:sec:gammazero}, which upgrades it to \lean{ConwayNortonConjecture}.
\end{dictionary}
\begin{physics}
$\Gamma_1(N)$ is the subgroup that preserves the twist \emph{exactly}, not only up to replacing $g$ by a power
$g^u$, $\gcd(u,N)=1$.  The difference between $\Gamma_1$ and $\Gamma_0$ is therefore the statement that the twisted
partition functions of $g$ and of $g^u$ coincide: a ``Galois'' symmetry of the twisted partition functions, which is
proved in \S\ref{ch16:sec:gammazero} (Theorem~\ref{ch16:thm:galois}).
\end{physics}
\begin{status}
Definition; printed in \leanpath{VOA/Moonshine/MoonshineCNAudit3.lean}.  Proved with no hypotheses as
\lean{conwayNortonGammaOne_holds} (\leanpath{VOA/Moonshine/MoonshineCNRationalTraces.lean}, Chapter~\ref{ch:22}).
\end{status}

\subsection{Is \texttt{ConwayNortonConjecture} the Conway--Norton conjecture?}\label{ch16:sec:fidelity}

This is the fidelity check on which the whole part rests.  We compare the Lean proposition with the conjecture as
stated in \S\ref{ch16:sec:history}, clause by clause, and then with the stronger statements in Conway and Norton's
paper.

\paragraph{(1) The group acting.}  The conjecture is about the Monster acting on FLM's $\Vnat$.  The Lean proposition
quantifies over \lean{AutVnat}, the automorphism group of the library's \lean{VnatVOA}.  Two separate identifications
are therefore involved.
\begin{itemize}
\item \emph{The module.}  \lean{VnatVOA} is the vertex operator algebra on the FLM space constructed in
  Chapters~\ref{ch:10}--\ref{ch:14}; Chapter~\ref{ch:14} records that its identification with FLM's $\Vnat$ (same
  space, possibly a different-looking twisted$\times$twisted vertex operator) is a paper argument (uniqueness of
  simple-current extensions).
\item \emph{The group.}  \lean{AutVnat} is the full automorphism group, whatever it is.  The Lean statement is
  therefore at least as strong as the conjecture for any subgroup, in particular for the Monster \emph{as a group of
  automorphisms of $\Vnat$}.  Chapter~\ref{ch:22} proves, with no hypotheses, that \lean{AutVnat} is isomorphic to the
  subgroup of AxiomMath's $\mathrm{AutO}$ consisting of the elements that act on $\Vnat$, which contains
  $\mathrm{Gmat}\supseteq2^{1+24}.\mathrm{Co}_1$; and, given \lean{GmatEqAutO}, that \lean{AutVnat} $\cong$
  $\mathrm{AutO}$.
\end{itemize}
So \lean{conwayNorton} is the conjecture for the group of all symmetries of the library's $\Vnat$; ``for the
Monster'' requires Chapter~\ref{ch:22} and, for all of $\mathrm{AutO}$, the hypothesis \lean{GmatEqAutO}.

\paragraph{(2) The series.}  \lean{mckayThompsonFun g} is $q^{-1}+\sum_{n\ge1}\tr(g|\Vnat_{n+1})q^n$, exactly
$T_g$ (Definition~\ref{ch16:def:mtfun}), with the conventions $q=e^{2\pi i\tau}$ and weight $n+1$ in degree $n$.  The
\lean{tsum} junk value is harmless (dictionary of Definition~\ref{ch16:def:mtfun}).  For $g=1$ it is $j-744$
(Theorem~\ref{ch16:thm:T1}).

\paragraph{(3) ``Normalised Hauptmodul of a genus-zero group''.}  Encoded by \lean{IsNormalizedHauptmodul}, which is
equivalent to the textbook notion for the groups in question (Proposition~\ref{ch16:prop:haupt}, on paper).  The
normalisation $q^{-1}+0+O(q)$ is Conway--Norton's: it matches $\tr(g|\Vnat_0)=1$ and $\tr(g|\Vnat_1)=0$.

\paragraph{(4) ``Contains $\Gamma_0(N)$ and is commensurable with $\chsixteenSL$''.}  Encoded by
\lean{IsMoonshineGroup}, clause by clause (dictionary of Definition~\ref{ch16:def:moonshinegroup}).  The Lean
statement has $\Gamma_0(N)$, not the weaker $\Gamma_1(N)$; the $\Gamma_1$ form is a separate proposition.

\paragraph{(5) The identity element.}  For $g=1$ the statement says that $J$ is the normalised Hauptmodul of some
moonshine group.  By the dictionary of Definition~\ref{ch16:def:hauptmodul}, $\pm\Gamma$ is then the full invariance
group of $J$, which is $\chsixteenSL$; so on paper the $g=1$ case is exactly the classical fact that $J$ is the
normalised Hauptmodul of $\chsixteenSL$.  This classical fact is not separately stated in the library
(Chapter~\ref{ch:L14} records that the textbook Hauptmodul theorem for $j$ is not formalised as such), but
\lean{conwayNorton 1}, combined with \lean{mckayThompsonFun_one}, gives in Lean that $j-744$ is a normalised
Hauptmodul of \emph{some} moonshine group.

\paragraph{(6) What Conway and Norton said in addition.}  Their paper proposes, for each class, the group $\Gamma_g$ in
the form $\Gamma_0(n|h)+e,f,\dots$, lying between $\Gamma_0(N)$ and its normaliser in $\chsixteenSLR$ with $N=nh$,
$n=o(g)$, $h\mid\gcd(24,n)$; and the explicit Hauptmodul.  The Lean statement asserts neither: it is the
\emph{qualitative} genus-zero statement (``the main conjecture'' in the literature, and the form in which Borcherds'
theorem is usually quoted).  It is therefore \textbf{weaker} than the full table-level conjecture, in three precise
respects: (i) $\Gamma$ is not identified; (ii) $\Gamma$ is not required to lie in the normaliser of $\Gamma_0(N)$;
(iii) $N$ is not related to $o(g)$.  It is \textbf{the same} as the genus-zero form of the conjecture displayed in
\S\ref{ch16:sec:history}.  The project's records state this explicitly: the statement is qualitative only (no identification of $\Gamma_g$ with the
Conway--Norton table).

\paragraph{(7) Could it be vacuous?}  The proposition is a universally quantified statement over \lean{AutVnat}; it
would be weak if \lean{AutVnat} were small.  It is not: Chapter~\ref{ch:22} shows that \lean{AutVnat} contains a copy
of $\mathrm{Gmat}$.  Each instance is strong: for each $g$ it determines the invariance group of $T_g$ up to sign
(dictionary of Definition~\ref{ch16:def:hauptmodul}) and asserts that it is a genus-zero congruence-type group.  There
is no hidden hypothesis: \lean{ConwayNortonConjecture} is a closed proposition, and \lean{conwayNorton} proves it with
no arguments.

\paragraph{Verdict.}  \lean{ConwayNortonConjecture} is the genus-zero (qualitative) form of the Conway--Norton
conjecture, for the full automorphism group of the library's moonshine module.  It is weaker than the table-level
conjecture in the three respects of (6), and it becomes a statement about the Monster through Chapter~\ref{ch:22}.

\begin{warning}
\textbf{Docstrings that say more than the Lean.}  The module docstring of
\leanpath{VOA/Moonshine/MoonshineCNStatement.lean} begins ``For every automorphism $g$ of the moonshine module
$\Vnat$ (the Monster, node M1)''; the parenthesis identifies \lean{AutVnat} with the Monster, which holds only given
\lean{GmatEqAutO} (Chapter~\ref{ch:22}).  The docstring of \lean{ConwayNortonConjecture} is careful (``With M1
(\lean{AutVnat ≃* 𝕄}) this is the conjecture for every element of the Monster'').  The docstring of
\lean{RationalTraces} calls it ``the M1-side input that upgrades $\Gamma_1$ to $\Gamma_0$''; that was the original plan, but the upgrade actually used does not need rationality (\S\ref{ch16:sec:gammazero}), and rationality
is now a corollary.  The Lean statements are unaffected; the statement file was frozen before the proof.
\end{warning}


%======================================================================
\section{Replication, replicable families and rational traces}\label{ch16:sec:replication}
%======================================================================

The statement file also fixes the interface between the Lie-algebra side of the proof (Chapters~\ref{ch:17}--\ref{ch:19},
which produce \emph{replication formulas}) and the analytic side (Chapter~\ref{ch:19} and \S\ref{ch16:sec:gammazero},
which turn replication into genus zero).  We present these definitions here because they are part of the fixed
statement file and because \S\ref{ch16:sec:gammazero} uses them.

\subsection{Hecke operators and replication}

For a $1$-periodic function $f$ on $\chsixteenH$ and $m\ge1$, the (unnormalised) Hecke operator is
\begin{equation}\label{ch16:eq:hecke}
  (f\,|\,T_m)(\tau)=\sum_{ad=m}\ \sum_{0\le b<d} f\Bigl(\frac{a\tau+b}{d}\Bigr),
\end{equation}
the sum over the $\sigma_1(m)$ upper-triangular integer matrices $\chsixteenmat ab0d$ of determinant $m$ with
$0\le b<d$.  If $f=J$, then $J|T_m$ is again $\chsixteenSL$-invariant, holomorphic on $\chsixteenH$, and has
$q$-expansion $q^{-m}+O(q)$: for fixed $a,d$ the sum over $b$ of $q^{k/d}$-terms keeps only exponents with $d\mid k$,
so only integral powers of $q$ survive; the most negative one is $q^{-m}$ (from $a=m$, $d=1$); and the constant term is
$\sigma_1(m)\,c(0)=0$ because $J$ has no constant term.  Hence it is a
polynomial in $J$: there is a unique monic polynomial $P_m$ of degree $m$, the \emph{Faber polynomial}, with
$P_m(J)=J|T_m$.

\emph{Replication} (Norton) generalises this to McKay--Thompson series: the Hecke sum of $T_g$ is formed with the
\emph{replicates} $T_{g^a}$ in the $a$-th block,
\begin{equation}\label{ch16:eq:replication}
  P_{g,m}\bigl(T_g(\tau)\bigr)=\sum_{ad=m}\ \sum_{0\le b<d}T_{g^a}\Bigl(\frac{a\tau+b}{d}\Bigr),
\end{equation}
where $P_{g,m}$ is the monic Faber polynomial of $T_g$ (the unique monic polynomial with $P_{g,m}(T_g)=q^{-m}+O(q)$).
A function satisfying \eqref{ch16:eq:replication} for every $m$, with its replicates, is called \emph{completely
replicable}.

\begin{example}[Replication for $J$ at $m=2$]\label{ch16:ex:rep2}
With $J=\sum_n c(n)q^n$, $c(-1)=1$, $c(0)=0$: $J^2=q^{-2}+2c(1)+2c(2)q+(2c(3)+c(1)^2)q^2+\cdots$, so
$P_2(x)=x^2-2c(1)$.  On the other side, $J(2\tau)+J(\tfrac\tau2)+J(\tfrac{\tau+1}2)=q^{-2}+2c(2)q+(c(1)+2c(4))q^2+\cdots$
(the two half-shifted terms add up to $2\sum_{n\text{ even}}c(n)q^{n/2}$).  Comparing coefficients of $q^2$ gives
Mahler's recursion $c(4)=c(3)+\tfrac12(c(1)^2-c(1))$: indeed
$864299970+\tfrac12(196884^2-196884)=20245856256$.
\end{example}

\begin{definition}[Hecke points and complete replicability]\label{ch16:def:replicable}
For $a,d\ge1$ and $b\in\N$, the \emph{Hecke point} of $\tau\in\chsixteenH$ is $(a\tau+b)/d\in\chsixteenH$.  The family
$(T_g)_{g\in\Aut(\Vnat)}$ is \emph{completely replicable} if for every $g$ and every $m\ge1$ there is a monic
polynomial $P$ of degree $m$ with
\begin{equation*}
  P\bigl(T_g(\tau)\bigr)=\sum_{a\mid m}\ \sum_{0\le b<m/a}T_{g^a}\Bigl(\frac{a\tau+b}{m/a}\Bigr)
  \qquad\text{for all }\tau\in\chsixteenH .
\end{equation*}
\end{definition}
\begin{leancode}
theorem heckePoint_im_pos (a d b : ℕ) (ha : 0 < a) (hd : 0 < d) (τ : ℍ) :
    0 < ((a * (τ : ℂ) + b) / d).im

noncomputable def heckePoint (a d : ℕ) (b : ℕ) (ha : 0 < a) (hd : 0 < d) (τ : ℍ) : ℍ :=
  ⟨(a * (τ : ℂ) + b) / d, heckePoint_im_pos a d b ha hd τ⟩

def CompletelyReplicable : Prop :=
  ∀ (g : AutVnat) (m : ℕ) (hm : 0 < m), ∃ P : Polynomial ℂ, P.Monic ∧ P.natDegree = m ∧
    ∀ τ : ℍ, P.eval (mckayThompsonFun g τ) =
      ∑ a ∈ m.divisors.attach, ∑ b ∈ Finset.range (m / a),
        mckayThompsonFun (g ^ (a : ℕ)) (heckePoint a (m / a) b (Nat.pos_of_mem_divisors a.2)
          (Nat.div_pos (Nat.divisor_le a.2) (Nat.pos_of_mem_divisors a.2)) τ)
\end{leancode}
(\leanpath{VOA/Moonshine/MoonshineCNStatement.lean}.)
\begin{dictionary}
\lean{heckePoint a d b ha hd τ} is $(a\tau+b)/d$ packaged as a point of \lean{ℍ}; the proof argument
\lean{heckePoint_im_pos} supplies $\operatorname{Im}>0$, since $\operatorname{Im}((a\tau+b)/d)=a\operatorname{Im}\tau/d$.
In \lean{CompletelyReplicable}: \lean{m.divisors.attach} runs over the divisors $a$ of $m$ (each packaged with the proof
that it divides $m$, used to prove $a>0$ and $m/a>0$); \lean{Finset.range (m / a)} is $\{0,\dots,d-1\}$ with
$d=m/a$; \lean{g ^ (a : ℕ)} is the replicate $g^a$; \lean{P.Monic ∧ P.natDegree = m} is ``monic of degree $m$'';
\lean{P.eval x} is $P(x)$.  This is \eqref{ch16:eq:replication} with an existential polynomial; as the docstring says,
such a $P$ is necessarily the Faber polynomial (both sides are $q^{-m}+O(q)$ once replicates are normalised), but
uniqueness is not part of the definition and is not needed.

\textbf{Strength.}  The quantifier is over \emph{all} $g\in\Aut(\Vnat)$ (no finite-order hypothesis); every $g$ has
finite order (Chapter~\ref{ch:20}), so this is not a restriction.  The definition is the textbook complete
replicability of the family $\{T_g\}$, stated as an identity of functions on $\chsixteenH$ rather than of formal
$q$-series; for finite-order $g$ the two are equivalent by absolute convergence.
\end{dictionary}
\begin{physics}
Replication is a shadow of the Hecke operators acting on the space of twisted partition functions.  Physically, the
Hecke sum computes the partition function of the symmetric-orbifold-like theory on an $m$-fold cover of the torus:
covering a torus with a $g$-twist by a torus $a$ times longer in one direction produces the twist $g^a$.  This is why
the replicate in the $a$-th block is $T_{g^a}$.  In Borcherds' proof the replication formulas come out of the twisted
denominator identities of the Monster Lie algebra, i.e.\ from the string spectrum.
\end{physics}
\begin{status}
Definitions; printed in \leanpath{VOA/Moonshine/MoonshineCNAudit3.lean}.  Chapter~\ref{ch:19} derives
\lean{CompletelyReplicable} from the twisted denominator identities, given simple-root generation of $\mathfrak m$ and
finite order (\lean{completelyReplicable_of_simpleRootGeneration}); both inputs are theorems (Chapter~\ref{ch:20}),
and the hypothesis-free instance is \lean{completelyReplicable_holds} (file
\leanpath{VOA/Moonshine/MoonshineCNRationalTraces.lean}, Chapter~\ref{ch:22}).
\end{status}

\subsection{Rational traces}

\begin{definition}[Rational traces]\label{ch16:def:rational}
\emph{All traces are rational} if $\tr(g|\Vnat_{n+1})\in\Q$ for all $g\in\Aut(\Vnat)$ and $n\ge0$.
\end{definition}
\begin{leancode}
def RationalTraces : Prop :=
  ∀ (g : AutVnat) (n : ℕ), ∃ r : ℚ, PowerSeries.coeff n (mckayThompson g) = r
\end{leancode}
(\leanpath{VOA/Moonshine/MoonshineCNStatement.lean}.)
\begin{dictionary}
\lean{∃ r : ℚ, ... = r} with the rational number \lean{r} coerced to $\C$.  For $n\ge0$ the coefficient is the trace
on weight $n+1$; weight $0$ (trace $1$) is not included but is trivially rational.  Classically all traces are
\emph{integers}; the Lean proposition asks only rationality, which is what the original route needed.
\end{dictionary}
\begin{physics}
Rationality of the character values of a symmetry on every energy level: the twisted partition functions are
invariant under the Galois group of $\Q(\zeta_{o(g)})/\Q$, equivalently $T_{g^u}=T_g$ for $u$ prime to $o(g)$.
\end{physics}
\begin{status}
Definition; printed in \leanpath{VOA/Moonshine/MoonshineCNAudit3.lean}.  Proved with no hypotheses as
\lean{rationalTraces} (Chapter~\ref{ch:22}); it is a \emph{consequence} of the moonshine theorem
(\S\ref{ch16:sec:gammazero}), not an input to it.
\end{status}

\subsection{Abstract replicable families}

The genus-zero argument of Chapter~\ref{ch:19} does not need $\Vnat$: it works for any family of $q$-series indexed by
the powers of an element of order $N$, dominated by $J$ and satisfying replication.  The statement file therefore
fixes an abstract interface.

\begin{definition}[$q$-series function; replicable family]\label{ch16:def:family}
For $c:\N\to\C$ let $Q_c(\tau)=q^{-1}+\sum_{n\ge1}c(n)q^n$ (the value $c(0)$ is ignored).  A \emph{completely replicable
family of period $N$} is a doubly indexed sequence $c(k,n)\in\C$ ($k,n\in\N$; $k$ indexes the replicate, $n$ the
coefficient) such that
\begin{enumerate}[label=(\roman*)]
\item $c(k+N,n)=c(k,n)$;
\item $|c(k,n)|\le|c(0,n)|$;
\item $c(0,n)=c(n)$, the coefficients of $J$;
\item for all $k$ and $m\ge1$ there is a monic polynomial $P$ of degree $m$ with
  $P(Q_{c(k,\cdot)}(\tau))=\sum_{a\mid m}\sum_{0\le b<m/a}Q_{c(ka,\cdot)}\bigl(\frac{a\tau+b}{m/a}\bigr)$.
\end{enumerate}
\end{definition}
\begin{leancode}
noncomputable def qSeriesFun (c : ℕ → ℂ) (τ : ℍ) : ℂ :=
  (qPar τ)⁻¹ + ∑' n : ℕ, c (n + 1) * qPar τ ^ (n + 1)

structure ReplicableFamily (N : ℕ) where
  c : ℕ → ℕ → ℂ
  periodic : ∀ k n, c (k + N) n = c k n
  dominated : ∀ k n, ‖c k n‖ ≤ ‖c 0 n‖
  zero_eq_J : ∀ n, c 0 n = PowerSeries.coeff (n + 1) JShift
  replicable : ∀ (k m : ℕ) (hm : 0 < m), ∃ P : Polynomial ℂ, P.Monic ∧ P.natDegree = m ∧
    ∀ τ : ℍ, P.eval (qSeriesFun (c k) τ) =
      ∑ a ∈ m.divisors.attach, ∑ b ∈ Finset.range (m / a),
        qSeriesFun (c (k * a)) (heckePoint a (m / a) b (Nat.pos_of_mem_divisors a.2)
          (Nat.div_pos (Nat.divisor_le a.2) (Nat.pos_of_mem_divisors a.2)) τ)
\end{leancode}
(\leanpath{VOA/Moonshine/MoonshineCNStatement.lean}.)
\begin{dictionary}
\lean{qSeriesFun c τ} sums \lean{c (n + 1) * q ^ (n + 1)} over $n\ge0$, i.e.\ $\sum_{m\ge1}c(m)q^m$: the index $0$ is
skipped, which is why ``the value \lean{c 0} is ignored''.  The \lean{tsum} junk convention applies as before.
\lean{ReplicableFamily N} is a structure carrying \emph{data} (the sequence \lean{c}) and four properties: (i)
\lean{periodic}, (ii) \lean{dominated}, (iii) \lean{zero_eq_J} (with $c(0,n)=[q^{n+1}]\mathtt{JShift}$, so $c(0,0)=0$ and
$c(0,n)=c(n)$), (iv) \lean{replicable}, with the replicate index \lean{k * a} in the $a$-th block, mirroring
$T_{(g^k)^a}=T_{g^{ka}}$.  Note that (iv) for $k$ involves only members $ka$, so a family is ``closed under taking
replicates''.  For $N=0$ the periodicity clause is empty; the instances used below have $N=o(g)>0$.

\textbf{Why (ii) and (iii)?}  (ii) gives convergence of every $Q_{c(k,\cdot)}$ on $\chsixteenH$ and the growth bounds of
Chapter~\ref{ch:19}; (iii) pins the family to $J$.  Without (iii) the abstract notion would admit degenerate
families: the project found that $q^{-1}+aq$ satisfies all replication identities with replicates $q^{-1}+a^mq$, for
every $a\in\C$ (\S\ref{ch16:sec:traps}).
\end{dictionary}
\begin{physics}
An abstract stand-in for the collection of twisted partition functions $\{Z[{}^1_{g^k}]\}_k$ of a cyclic symmetry
group, retaining only what the genus-zero argument needs: the twisted partition functions of the powers of $g$, their
domination by the untwisted one, and the Hecke (covering) relations among them.
\end{physics}
\begin{status}
Definitions; printed in \leanpath{VOA/Moonshine/MoonshineCNAudit3.lean}, standard axioms.
\end{status}

\begin{proposition}[The McKay--Thompson family is a replicable family]\label{ch16:prop:mtfamily}
\begin{enumerate}[label=(\alph*)]
\item For every $g$, $T_g=Q_{c_g}$ where $c_g(n)=[X^n]\mathrm{MT}_g$.
\item If the family $\{T_g\}$ is completely replicable and $g$ has finite order $o(g)$, then
  $c(k,n)=[X^n]\mathrm{MT}_{g^k}$ is a completely replicable family of period $o(g)$, and $o(g)>0$.
\end{enumerate}
\end{proposition}
\begin{leancode}
theorem mckayThompsonFun_eq_qSeriesFun (g : AutVnat) :
    mckayThompsonFun g = qSeriesFun (fun n => PowerSeries.coeff n (mckayThompson g))

noncomputable def mckayThompsonReplicableFamily (hrep : CompletelyReplicable)
    (g : AutVnat) (hg : IsOfFinOrder g) : ReplicableFamily (orderOf g) where
  c k n := PowerSeries.coeff n (mckayThompson (g ^ k))

theorem mckayThompson_replicableFamily_exists (hrep : CompletelyReplicable)
    (g : AutVnat) (hg : IsOfFinOrder g) :
    0 < orderOf g ∧ ∃ F : ReplicableFamily (orderOf g),
      ∀ k n, F.c k n = PowerSeries.coeff n (mckayThompson (g ^ k))
\end{leancode}
(\leanpath{VOA/Moonshine/MoonshineCNStatement.lean}.)
\begin{dictionary}
(a) is an equality of functions: \lean{mckayThompsonFun} sums from $n=0$, \lean{qSeriesFun} from $n=1$, and the $n=0$
term vanishes (Proposition~\ref{ch16:prop:mtprops}(b)).  The proof treats the summable and the non-summable case
separately, so (a) holds even with the \lean{tsum} junk convention.  (b) is stated twice: as a definition producing the
structure (with \lean{c k n := PowerSeries.coeff n (mckayThompson (g ^ k))}), and as an existence theorem that also
records $0<o(g)$ (\lean{orderOf g}, Mathlib's order, which is $0$ for elements of infinite order) and the formula for
the coefficients.  Hypotheses: complete replicability and finite order of this $g$; both are theorems (Chapters~\ref{ch:19}
and~\ref{ch:20}).
\end{dictionary}
\paragraph{Proof on paper.}  (a) was explained.  (b) (i) $g^{k+o(g)}=g^k$.  (ii) Theorem~\ref{ch16:thm:tracebounds}(c)
applied to $g^k$ (which has finite order).  (iii) Proposition~\ref{ch16:prop:mtprops}(d).  (iv) Complete
replicability for $g^k$, rewritten with (a) and $(g^k)^a=g^{ka}$. \qed
\begin{physics}
The twisted partition functions of the powers of one symmetry form a closed system under the covering (Hecke)
relations; this is the input to the genus-zero argument.
\end{physics}
\begin{status}
Unconditional (with the stated hypotheses); printed in \leanpath{VOA/Moonshine/MoonshineCNAudit3.lean}, standard
axioms.
\end{status}

%======================================================================
\section{From \texorpdfstring{$\Gamma_1(N)$}{Gamma1(N)} to \texorpdfstring{$\Gamma_0(N)$}{Gamma0(N)}}\label{ch16:sec:gammazero}
%======================================================================

\subsection{Why this step is needed, and the idea}

The case-free genus-zero argument of Chapter~\ref{ch:19} proves that each $T_g$ is the normalised Hauptmodul of a group
$\Gamma\supseteq\Gamma_1(M)$ commensurable with $\chsixteenSL$ (the proposition \lean{ConwayNortonGammaOne}).  The
classical way to get $\Gamma_0(M)$ uses rationality of the traces: $\Gamma_0(M)$ is generated by $\Gamma_1(M)$ and
matrices $\delta$ whose lower-right entry $d$ is a unit mod $M$, and invariance under such $\delta$ amounts to
$T_{g^d}=T_g$, a Galois symmetry.  The original plan obtained it from
a rational form of $\Vnat$ and Tits' theorem; that route turned out to need group-theoretic classification input
(\S\ref{ch16:sec:traps}).

The library instead proves $\Gamma_0(M)\le\pm\Gamma$ \emph{from the $\Gamma_1$ theorem and replication at primes
alone}.  The key idea is the following.  Let $p>M$ be prime and choose
$m,n\in\Z$ with $mp-nM=1$; put
\begin{equation*}
  \sigma_p=\chsixteenmat mnMp\in\Gamma_0(M),\qquad \sigma_p\equiv\chsixteenmat{p^{-1}}{*}0p\pmod M .
\end{equation*}
The replication formula at $p$ involves $f=T_g$ at the $p$ points $(\tau+j)/p$ and the replicate $f'=T_{g^p}$ at
$p\tau$.  A coset computation for the single matrix $L_M=\chsixteenmat10M1$ shows that $\psi=f\circ\sigma_p-f'$ is
invariant under $\tau\mapsto\tau+1$ and under $L_{M/p}=\chsixteenmat10{M/p}1$ with $0<M/p<1$.  By Shimizu's lemma these
two maps generate a non-discrete group; an invariant holomorphic function is then constant.  So
$f(\sigma_p\tau)=f'(\tau)+c$; comparing with the normalisation at the cusps and using \ref{ch16:H4} gives $f'=f$ and
$\pm\sigma_p\in\Gamma$.  Dirichlet's theorem supplies, for every $\delta\in\Gamma_0(M)$, a prime $p>M$ with
$p\equiv d\pmod M$, and then $\delta\in\Gamma_1(M)\sigma_p\subseteq\pm\Gamma$.  As a by-product, $T_{g^u}=T_g$ for $u$
prime to $o(g)$.

We now go through the steps, each with its Lean statement.  The files are
\leanpath{VOA/Moonshine/MoonshineCNGammaZeroShimizu.lean} (Shimizu and constancy),
\leanpath{VOA/Moonshine/MoonshineCNGammaZeroCosets.lean} (coset matching),
\leanpath{VOA/Moonshine/MoonshineCNGammaZeroBridge.lean} (asymptotics, $\pm\Gamma$, the bridge),
\leanpath{VOA/Moonshine/MoonshineCNGammaZeroGalois.lean} (Galois invariance) and
\leanpath{VOA/Moonshine/MoonshineCNGammaZeroFamily.lean} (the abstract family version).  All are in the namespace
\lean{VOA.Moonshine.ConwayNorton}.

\subsection{Shimizu's lemma and constancy}

\begin{definition}[Unipotent matrices]\label{ch16:def:unipotent}
For $x,y\in\chsixteenR$: $T_x=\chsixteenmat1x01$ (the translation $\tau\mapsto\tau+x$) and $L_y=\chsixteenmat10y1$ (the map
$\tau\mapsto\tau/(y\tau+1)$), both in $\chsixteenSLR$.
\end{definition}
\begin{leancode}
noncomputable def glOfDetOne (a b c d : ℝ) (h : a * d - b * c = 1) : GL (Fin 2) ℝ :=
  GeneralLinearGroup.mkOfDetNeZero !![a, b; c, d] (by rw [det_fin_two_of, h]; exact one_ne_zero)

noncomputable def glUpperUnipotent (x : ℝ) : GL (Fin 2) ℝ := glOfDetOne 1 x 0 1 (by ring)

noncomputable def glLowerUnipotent (y : ℝ) : GL (Fin 2) ℝ := glOfDetOne 1 0 y 1 (by ring)
\end{leancode}
(\leanpath{VOA/Moonshine/MoonshineCNGammaZeroShimizu.lean}.)
\begin{dictionary}
\lean{glOfDetOne a b c d h} is the matrix $\chsixteenmat abcd$ as an element of $\chsixteenGLR$, built by Mathlib's
\lean{GeneralLinearGroup.mkOfDetNeZero} from the fact that its determinant $ad-bc=1$ (\lean{h}) is non-zero;
\lean{!![a, b; c, d]} is Mathlib's matrix literal.  The definitions are the textbook
matrices.
\end{dictionary}
\begin{physics}
$T_x$ shifts the torus modulus; $L_y$ is a ``Dehn twist in the other direction'' when $y\in\Z$, and a non-integral
shear otherwise.  A modular-invariant quantity is invariant under $T_1$ and $L_1$; invariance under a shear $L_y$
with $0<y<1$ is never expected of a physical partition function, and the lemma below shows it forces triviality.
\end{physics}
\begin{status}
Definitions; used by the audited theorems below.
\end{status}

\begin{theorem}[Constancy under a non-discrete set; Shimizu]\label{ch16:thm:shimizu}
Let $\psi:\chsixteenH\to\C$ be holomorphic.
\begin{enumerate}[label=(\alph*)]
\item If $C_n\in\chsixteenGLR$ converge entrywise to the identity, each $C_n$ moves some point of $\chsixteenH$, and
  $\psi(C_n\tau)=\psi(\tau)$ for all $n,\tau$, then $\psi$ is constant.
\item If $\psi(\tau+1)=\psi(\tau)$ and $\psi(L_y\tau)=\psi(\tau)$ for all $\tau$, for some $0<y<1$, then $\psi$ is
  constant.
\end{enumerate}
\end{theorem}
\begin{leancode}
theorem eq_of_invariant_of_entries_tendsto_one {ψ : ℍ → ℂ}
    (hψ : MDifferentiable 𝓘(ℂ) 𝓘(ℂ) ψ) {C : ℕ → GL (Fin 2) ℝ}
    (hC : ∀ i j, Tendsto (fun n => (C n : Matrix (Fin 2) (Fin 2) ℝ) i j) atTop
      (𝓝 ((1 : Matrix (Fin 2) (Fin 2) ℝ) i j)))
    (hmove : ∀ n, ∃ τ : ℍ, C n • τ ≠ τ) (hinv : ∀ n τ, ψ (C n • τ) = ψ τ) (τ σ : ℍ) :
    ψ τ = ψ σ

theorem eq_of_invariant_upper_lower_unipotent {ψ : ℍ → ℂ} (hψ : MDifferentiable 𝓘(ℂ) 𝓘(ℂ) ψ)
    {y : ℝ} (hy0 : 0 < y) (hy1 : y < 1) (hT : ∀ τ, ψ (glUpperUnipotent 1 • τ) = ψ τ)
    (hL : ∀ τ, ψ (glLowerUnipotent y • τ) = ψ τ) (τ σ : ℍ) : ψ τ = ψ σ
\end{leancode}
(\leanpath{VOA/Moonshine/MoonshineCNGammaZeroShimizu.lean}.)
\begin{dictionary}
\lean{hC}: for all entries $(i,j)$, $(C_n)_{ij}\to\delta_{ij}$ as $n\to\infty$ (\lean{atTop} on $\N$).  \lean{hmove}:
each $C_n$ acts non-trivially on $\chsixteenH$.  \lean{hinv}: invariance.  Conclusion: $\psi(\tau)=\psi(\sigma)$ for all
$\tau,\sigma$, i.e.\ $\psi$ is constant.  In (b), \lean{hT}, \lean{hL} are invariance under $T_1$ and $L_y$, with
\lean{0 < y < 1}.  Note that (b) does not require $\det>0$ or discreteness anywhere: it is a statement about any
holomorphic function with these two invariances.  A sibling \lean{eq_of_invariant_of_tendsto_one} states (a) with
convergence in the topology of $\chsixteenGLR$.  Both are the textbook statements.
\end{dictionary}
\paragraph{Proof on paper.}  (a) For large $n$, $\det C_n>0$, so $C_n$ is a M\"obius map; a non-trivial real
M\"obius map of positive determinant fixes at most one point of $\chsixteenH$ (the fixed-point equation
$c\tau^2+(d-a)\tau-b=0$ has at most one root in $\chsixteenH$ unless $c=b=0$, $a=d$).  Hence each such $C_n$ moves
$i$ or $2i$, and one of the two, say $\tau_0$, is moved by infinitely many $C_n$.  Since $C_n\to1$, $C_n\tau_0\to\tau_0$
with $C_n\tau_0\ne\tau_0$ infinitely often, and $\psi(C_n\tau_0)=\psi(\tau_0)$: the level set
$\{\psi=\psi(\tau_0)\}$ accumulates at $\tau_0$.  By the identity theorem $\psi\equiv\psi(\tau_0)$.  (This avoids the
inverse function theorem used in the paper blueprint.)

(b) \emph{Shimizu's sequence.}  Let $B_0=L_y$ and $B_{n+1}=B_nT_1B_n^{-1}$.  If $B_n=\chsixteenmat{a_n}{*}{c_n}{*}$ then
a direct computation gives
\begin{equation*}
  B_{n+1}=\chsixteenmat{1-a_nc_n}{a_n^2}{-c_n^2}{1+a_nc_n},
  \qquad\text{so}\qquad c_{n+1}=-c_n^2,\ |c_n|=y^{2^n}\to0,
\end{equation*}
and $a_{n+1}=1-a_nc_n$, from which $a_n\to1$ (the corrections are summable).  Hence $B_n\to T_1$ and
$C_n=B_{n+1}T_1^{-1}\to1$, while the lower-left entry of $C_n$ is $-c_n^2\neq0$, so $C_n$ moves points.  Every $C_n$
lies in the group generated by $T_1$ and $L_y$, under which $\psi$ is invariant; apply (a). \qed
\begin{physics}
A function on moduli space invariant under a non-discrete group of modular transformations cannot vary: there is no
``partition function'' with such a symmetry except a constant.  Shimizu's lemma says: a discrete group containing the
unit translation cannot contain a shear $L_y$ with $0<|y|<1$.
\end{physics}
\begin{status}
Unconditional; both printed in \leanpath{VOA/Moonshine/MoonshineCNGammaZeroAudit.lean} (47 reports, all
\lean{[propext, Classical.choice, Quot.sound]}).
\end{status}

\subsection{The diamond lift and Hecke coset matching}

\begin{definition}[Diamond lift; level sum]\label{ch16:def:diamond}
Let $M\ge1$, $p$ prime, and $m,n\in\Z$ with $mp-nM=1$.  The \emph{diamond lift} is
$\sigma=\chsixteenmat mnMp\in\Gamma_0(M)$.  For $f:\chsixteenH\to\C$ the \emph{level sum} is
\begin{equation*}
  S(\tau)=\sum_{j=0}^{p-1}f\Bigl(\frac{\tau+j}{p}\Bigr)+f\bigl(\sigma(p\tau)\bigr).
\end{equation*}
\end{definition}
\begin{leancode}
def diamondLift (M p : ℕ) (m n : ℤ) (h : m * p - n * M = 1) : SL(2, ℤ) :=
  slOfDetOne m n M p h

noncomputable def heckeLevelSum (f : ℍ → ℂ) (σ : SL(2, ℤ)) (p : ℕ) (hp : 0 < p) (τ : ℍ) : ℂ :=
  ∑ j ∈ Finset.range p, f (heckePoint 1 p j one_pos hp τ) +
    f (mapGL ℝ σ • heckePoint p 1 0 hp one_pos τ)
\end{leancode}
(\leanpath{VOA/Moonshine/MoonshineCNGammaZeroCosets.lean}.)
\begin{dictionary}
\lean{slOfDetOne m n M p h} is $\chsixteenmat mnMp\in\chsixteenSL$ (same file).  In \lean{heckeLevelSum},
\lean{heckePoint 1 p j} is $(\tau+j)/p$, \lean{heckePoint p 1 0} is $p\tau$, and \lean{mapGL ℝ σ • ·} is the M\"obius
action of $\sigma$.  The lemma \lean{diamondLift_mem_Gamma0} records $\sigma\in\Gamma_0(M)$.
\end{dictionary}
\begin{physics}
$\sigma$ is an element of $\Gamma_0(M)$ whose action on a $g$-twisted torus replaces the twist by its $p$-th power
($\sigma\equiv\operatorname{diag}(p^{-1},p)$ mod $M$); the level sum is the Hecke-covering sum of $f$ in which the
$p$-fold-scaled term is transported by $\sigma$.
\end{physics}
\begin{status}
Definitions; used by the audited theorems below (\lean{diamondLift_mem_Gamma0} is printed in
\leanpath{VOA/Moonshine/MoonshineCNGammaZeroAudit.lean}).
\end{status}

\begin{theorem}[Hecke coset matching]\label{ch16:thm:cosets}
Let $f$ be $\Gamma_1(M)$-invariant, $p$ prime and $mp-nM=1$.  Then:
\begin{enumerate}[label=(\alph*)]
\item $f\circ\sigma$ is $1$-periodic;
\item $S(L_M\tau)=S(\tau)$ for all $\tau$.
\end{enumerate}
\end{theorem}
\begin{leancode}
theorem diamond_periodic {f : ℍ → ℂ} {M p : ℕ} {m n : ℤ} (h : m * p - n * M = 1)
    (hf : ∀ γ ∈ Gamma1 M, ∀ τ : ℍ, f (mapGL ℝ γ • τ) = f τ) (τ : ℍ) :
    f (mapGL ℝ (diamondLift M p m n h) • (glUpperUnipotent 1 • τ)) =
      f (mapGL ℝ (diamondLift M p m n h) • τ)

theorem heckeLevelSum_lowerSL {f : ℍ → ℂ} {M p : ℕ} (hp : p.Prime) {m n : ℤ}
    (hmn : m * p - n * M = 1) (hf : ∀ γ ∈ Gamma1 M, ∀ τ : ℍ, f (mapGL ℝ γ • τ) = f τ) (τ : ℍ) :
    heckeLevelSum f (diamondLift M p m n hmn) p hp.pos (mapGL ℝ (lowerSL M) • τ) =
      heckeLevelSum f (diamondLift M p m n hmn) p hp.pos τ
\end{leancode}
(\leanpath{VOA/Moonshine/MoonshineCNGammaZeroCosets.lean}.)
\begin{dictionary}
\lean{hf}: $f(\gamma\tau)=f(\tau)$ for all $\gamma\in\Gamma_1(M)$.  \lean{lowerSL M} is $L_M=\chsixteenmat10M1$ and
\lean{glUpperUnipotent 1} is $T_1$.  No hypothesis $p>M$ or $M>0$ is needed here.  The statements are exactly (a), (b).
\end{dictionary}
\paragraph{Proof on paper.}  (a) $\sigma T_1=\gamma\sigma$ with $\gamma=\chsixteenmat{1-mM}{m^2}{-M^2}{1+mM}\in\Gamma_1(M)$
(multiply out, using $mp-nM=1$), so $f(\sigma(\tau+1))=f(\gamma\sigma\tau)=f(\sigma\tau)$.

(b) Reduce modulo $p$: since $mp-nM=1$, $nM\equiv-1$, so $M$ is invertible mod $p$ and $-M^{-1}\equiv n$.  Let
$j_0\equiv n$ and $j^*\equiv-n$ in $[0,p)$.  There are three matrix identities, each checked by multiplying out:
\begin{itemize}
\item \emph{generic $j\ne j_0$}: with $j'\equiv j/(1+jM)\pmod p$ in $[0,p)$ and $e=(j-(1+jM)j')/p\in\Z$,
  \begin{equation*}
    \chsixteenmat1j0p\,L_M=\gamma_j\chsixteenmat1{j'}0p,\qquad \gamma_j=\chsixteenmat{1+jM}{e}{pM}{1-Mj'}\in\Gamma_1(M);
  \end{equation*}
\item \emph{exceptional $j_0=n+kp$}: $\chsixteenmat1{j_0}0p\,L_M=T_k\,\sigma\,\chsixteenmat p001$;
\item \emph{extra term}: with $e=(n-(mp+nM)j^*)/p\in\Z$,
  $\sigma\chsixteenmat p001L_M=\gamma^*\chsixteenmat1{j^*}0p$, $\gamma^*=\chsixteenmat{mp+nM}{e}{2Mp}{1-2Mj^*}\in\Gamma_1(M)$
  (the $\Gamma_1$ condition holds because $mp+nM\equiv mp\equiv1\pmod M$).
\end{itemize}
The map $j\mapsto j'$ is the M\"obius map $x\mapsto x/(1+Mx)$ of $\chsixteenP(\mathbb F_p)$, which sends $j_0=-1/M$ to
$\infty$ and $\infty$ to $1/M=j^*$; so it is a bijection $[0,p)\setminus\{j_0\}\to[0,p)\setminus\{j^*\}$.  Applying
$f$ and $\Gamma_1(M)$-invariance ($T_k$, $\gamma_j$, $\gamma^*\in\Gamma_1(M)$):
\begin{equation*}
  S(L_M\tau)=\sum_{j'\ne j^*}f\Bigl(\frac{\tau+j'}p\Bigr)+f\bigl(\sigma(p\tau)\bigr)+f\Bigl(\frac{\tau+j^*}p\Bigr)=S(\tau).
  \qquad\qed
\end{equation*}
\begin{physics}
The Hecke correspondence at level $p$ permutes its $p+1$ ``covering tori'' under a modular transformation of the base;
the computation identifies the permutation explicitly for the generator $L_M$, with one covering torus exchanged
between the ``$(\tau+j)/p$'' family and the ``$p\tau$'' term.
\end{physics}
\begin{status}
Unconditional; printed in \leanpath{VOA/Moonshine/MoonshineCNGammaZeroAudit.lean} together with the three coset
identities \lean{heckeShift_lowerSL}, \lean{heckeShift_lowerSL_exceptional}, \lean{diamondScale_lowerSL} and the
index bijection lemmas.
\end{status}

\begin{theorem}[The replicate differs from the diamond image by a constant]\label{ch16:thm:diamondconst}
Let $0<M<p$, $p$ prime, $mp-nM=1$, $f$ holomorphic and $\Gamma_1(M)$-invariant, $f'$ holomorphic and $1$-periodic, and
suppose that for some function $\Phi:\C\to\C$
\begin{equation*}
  \Phi\bigl(f(\tau)\bigr)=\sum_{j=0}^{p-1}f\Bigl(\frac{\tau+j}p\Bigr)+f'(p\tau)\qquad(\tau\in\chsixteenH).
\end{equation*}
Then $\psi=f\circ\sigma-f'$ is invariant under $L_{M/p}$, and there is $c\in\C$ with $f(\sigma\tau)=f'(\tau)+c$ for all
$\tau$.
\end{theorem}
\begin{leancode}
theorem diamond_sub_replicate_lowerUnipotent {f f' : ℍ → ℂ} {Φ : ℂ → ℂ} {M p : ℕ}
    (hp : p.Prime) {m n : ℤ} (hmn : m * p - n * M = 1)
    (hf : ∀ γ ∈ Gamma1 M, ∀ τ : ℍ, f (mapGL ℝ γ • τ) = f τ)
    (hrep : ∀ τ : ℍ, Φ (f τ) = ∑ j ∈ Finset.range p, f (heckePoint 1 p j one_pos hp.pos τ) +
      f' (heckePoint p 1 0 hp.pos one_pos τ)) (z : ℍ) :
    f (mapGL ℝ (diamondLift M p m n hmn) • (glLowerUnipotent ((M : ℝ) / p) • z)) -
        f' (glLowerUnipotent ((M : ℝ) / p) • z) =
      f (mapGL ℝ (diamondLift M p m n hmn) • z) - f' z

theorem diamond_eq_replicate_add_const {f f' : ℍ → ℂ} {Φ : ℂ → ℂ} {M p : ℕ} (hM : 0 < M)
    (hp : p.Prime) (hMp : M < p) {m n : ℤ} (hmn : m * p - n * M = 1)
    (hfhol : MDifferentiable 𝓘(ℂ) 𝓘(ℂ) f)
    (hf : ∀ γ ∈ Gamma1 M, ∀ τ : ℍ, f (mapGL ℝ γ • τ) = f τ)
    (hf'hol : MDifferentiable 𝓘(ℂ) 𝓘(ℂ) f')
    (hf'per : ∀ τ : ℍ, f' (glUpperUnipotent 1 • τ) = f' τ)
    (hrep : ∀ τ : ℍ, Φ (f τ) = ∑ j ∈ Finset.range p, f (heckePoint 1 p j one_pos hp.pos τ) +
      f' (heckePoint p 1 0 hp.pos one_pos τ)) :
    ∃ c : ℂ, ∀ τ : ℍ, f (mapGL ℝ (diamondLift M p m n hmn) • τ) = f' τ + c
\end{leancode}
(\leanpath{VOA/Moonshine/MoonshineCNGammaZeroCosets.lean}.)
\begin{dictionary}
\lean{hrep} is the replication identity at $p$ in ``two-block'' form, with an \emph{arbitrary} function $\Phi$ of
$f$ on the left: monicity and degree of the Faber polynomial are not used.  \lean{hfhol}, \lean{hf'hol}: holomorphy;
\lean{hf}: $\Gamma_1(M)$-invariance; \lean{hf'per}: $f'(\tau+1)=f'(\tau)$; \lean{hM}, \lean{hMp}: $0<M<p$.  The first
theorem is the invariance of $\psi$ under $L_{M/p}$ (\lean{glLowerUnipotent ((M : ℝ) / p)}), the second the
constancy.  Statements as in the theorem.
\end{dictionary}
\paragraph{Proof on paper.}  By Theorem~\ref{ch16:thm:cosets}(b) $S$ is $L_M$-invariant; $\Phi\circ f$ is too, since
$L_M\in\Gamma_1(M)$.  The replication identity says $S(\tau)-\Phi(f(\tau))=f(\sigma(p\tau))-f'(p\tau)=\psi(p\tau)$.  Hence
$\psi(p\,L_M\tau)=\psi(p\tau)$.  Now $p\cdot L_M\tau=p\tau/(M\tau+1)=L_{M/p}(p\tau)$, and $\tau\mapsto p\tau$ is onto
$\chsixteenH$; so $\psi$ is $L_{M/p}$-invariant.  It is also $1$-periodic (Theorem~\ref{ch16:thm:cosets}(a) and
periodicity of $f'$) and holomorphic.  Since $0<M/p<1$, Theorem~\ref{ch16:thm:shimizu}(b) shows $\psi$ is constant. \qed
\begin{physics}
Replication relates the $p$-th Hecke transform of $T_g$ to the replicate $T_{g^p}$; the computation shows that the
replicate is literally the twisted partition function $T_g$ seen through the modular transformation $\sigma$, up to a
constant.
\end{physics}
\begin{status}
Unconditional (with the stated hypotheses); both printed in \leanpath{VOA/Moonshine/MoonshineCNGammaZeroAudit.lean},
standard axioms.
\end{status}

\subsection{Asymptotic comparison and the main lemma}

\begin{lemma}[Asymptotic comparison]\label{ch16:lem:affine}
Let $f,f'$ be normalised ($f-q^{-1}\to0$ and $f'-q^{-1}\to0$ as $\operatorname{Im}\tau\to\infty$), and suppose
$f(a\tau+b)=f'(\tau)+c$ for all $\tau$, with $a>0$, $b\in\chsixteenR$, $c\in\C$.  Then $a=1$, $b\in\Z$ and $c=0$.
\end{lemma}
\begin{leancode}
theorem affine_eq_of_normalized {f f' : ℍ → ℂ}
    (hf : Tendsto (fun τ : ℍ => f τ - (qPar τ)⁻¹) atImInfty (𝓝 0))
    (hf' : Tendsto (fun τ : ℍ => f' τ - (qPar τ)⁻¹) atImInfty (𝓝 0))
    {φ : ℍ → ℍ} {a b : ℝ} (ha : 0 < a) (hφ : ∀ τ : ℍ, (φ τ : ℂ) = a * τ + b) {c : ℂ}
    (heq : ∀ τ : ℍ, f (φ τ) = f' τ + c) :
    a = 1 ∧ (∃ k : ℤ, b = k) ∧ c = 0
\end{leancode}
(\leanpath{VOA/Moonshine/MoonshineCNGammaZeroBridge.lean}.)
\begin{dictionary}
The affine map is given as a function \lean{φ : ℍ → ℍ} with \lean{(φ τ : ℂ) = a * τ + b} (\lean{hφ}); \lean{heq} is
$f(\varphi\tau)=f'(\tau)+c$.  Conclusion \lean{a = 1 ∧ (∃ k : ℤ, b = k) ∧ c = 0}.  As stated.
\end{dictionary}
\paragraph{Proof on paper.}  Write $E(\tau)=q(\varphi\tau)^{-1}-q(\tau)^{-1}$; the hypotheses give $E\to c$.  But
$|q(\varphi\tau)^{-1}|=e^{2\pi a\operatorname{Im}\tau}$ and $|q(\tau)^{-1}|=e^{2\pi\operatorname{Im}\tau}$, and if
$a\ne1$ their difference tends to infinity; so $a=1$.  Then $q(\varphi\tau)=q(\tau)\,\omega$ with $\omega=e^{2\pi ib}$,
and $E=q^{-1}(\omega^{-1}-1)$, which is unbounded unless $\omega=1$, i.e.\ $b\in\Z$, and then $E=0$, so $c=0$. \qed
\begin{physics}
The vacuum pole $q^{-1}$ fixes the normalisation of the modulus at the cusp: no rescaling, non-integral shift or
additive constant is compatible with two partition functions having the same leading vacuum term and no constant term.
\end{physics}
\begin{status}
Unconditional; printed in \leanpath{VOA/Moonshine/MoonshineCNGammaZeroAudit.lean}, standard axioms.
\end{status}

\begin{theorem}[Replicate equals Hauptmodul; diamond in $\pm\Gamma$]\label{ch16:thm:main}
Let $f$ be a normalised Hauptmodul for $\Gamma\le\chsixteenSLR$ with $\Gamma_1(M)\le\Gamma$, $M\ge1$.  Let $p>M$ be prime,
$mp-nM=1$, and $f'$ holomorphic, $1$-periodic and normalised, with
$\Phi(f(\tau))=\sum_{j<p}f((\tau+j)/p)+f'(p\tau)$ for some $\Phi:\C\to\C$.  Then $f'=f$, and $\sigma\in\Gamma$ or
$-\sigma\in\Gamma$.
\end{theorem}
\begin{leancode}
theorem replicate_eq_and_diamond_mem {Γ : Subgroup (GL (Fin 2) ℝ)} {f f' : ℍ → ℂ}
    {Φ : ℂ → ℂ} {M p : ℕ} (hH : IsNormalizedHauptmodul Γ f) (hdet : Γ.HasDetOne)
    (hΓ1 : gamma1GL M ≤ Γ) (hM : 0 < M) (hp : p.Prime) (hMp : M < p) {m n : ℤ}
    (hmn : m * p - n * M = 1) (hf'hol : MDifferentiable 𝓘(ℂ) 𝓘(ℂ) f')
    (hf'per : ∀ τ : ℍ, f' (glUpperUnipotent 1 • τ) = f' τ)
    (hf'norm : Tendsto (fun τ : ℍ => f' τ - (qPar τ)⁻¹) atImInfty (𝓝 0))
    (hrep : ∀ τ : ℍ, Φ (f τ) = ∑ j ∈ Finset.range p, f (heckePoint 1 p j one_pos hp.pos τ) +
      f' (heckePoint p 1 0 hp.pos one_pos τ)) :
    f' = f ∧ (mapGL ℝ (diamondLift M p m n hmn) ∈ Γ ∨ -mapGL ℝ (diamondLift M p m n hmn) ∈ Γ)
\end{leancode}
(\leanpath{VOA/Moonshine/MoonshineCNGammaZeroBridge.lean}.)
\begin{dictionary}
\lean{hH}: \lean{IsNormalizedHauptmodul Γ f}; \lean{hdet}: $\Gamma\le\chsixteenSLR$; \lean{hΓ1}: $\Gamma_1(M)\le\Gamma$;
\lean{hM}: $M>0$; \lean{hp}, \lean{hMp}: $p$ prime, $p>M$; \lean{hmn}: $mp-nM=1$; the three hypotheses on $f'$; and
\lean{hrep}, the replication identity at $p$ with arbitrary $\Phi$.  Commensurability of $\Gamma$ is \emph{not} needed.
The conclusion \lean{f' = f ∧ (σ ∈ Γ ∨ -σ ∈ Γ)}, with \lean{σ = mapGL ℝ (diamondLift M p m n hmn)}.  As stated.
\end{dictionary}
\paragraph{Proof on paper.}  By Theorem~\ref{ch16:thm:diamondconst}, $f(\sigma\tau)=f'(\tau)+c$; the right-hand side is
unbounded as $\operatorname{Im}\tau\to\infty$ ($f'\sim q^{-1}$).  The point $\sigma\infty=m/M$ is a cusp of
$\Gamma_1(M)\le\Gamma$.  If it were not $\Gamma$-equivalent to $\infty$, it would be an other cusp, and by \ref{ch16:H4}
$f(\sigma\tau)$ would tend to a finite value: contradiction.  So $\gamma\infty=\sigma\infty$ for some $\gamma\in\Gamma$,
and $\delta=\gamma^{-1}\sigma$ fixes $\infty$ and has determinant $1$: $\delta=\chsixteenmat{a}{b}{0}{a^{-1}}$, acting by
$\tau\mapsto a^2\tau+ab$.  By invariance, $f(\sigma\tau)=f(\delta\tau)$, so $f(a^2\tau+ab)=f'(\tau)+c$.
Lemma~\ref{ch16:lem:affine} gives $a^2=1$, $ab=k\in\Z$, $c=0$.  Hence $\delta=\pm T_k$, $f'(\tau)=f(\tau+k)=f(\tau)$ (as
$T_k\in\Gamma_1(M)$), and $\sigma=\gamma\delta=\pm\gamma T_k\in\pm\Gamma$. \qed
\begin{physics}
If the twisted partition function of $g$ is a Hauptmodul, then the twisted partition function of $g^p$ (for a large
prime $p$) is the same function, and the modular transformation that turns the $g$-twist into a $g^p$-twist is a
symmetry of it.  The only physical input is the absence of a second vacuum-like pole at the other cusps.
\end{physics}
\begin{status}
Unconditional (with the stated hypotheses); printed in \leanpath{VOA/Moonshine/MoonshineCNGammaZeroAudit.lean},
standard axioms.
\end{status}

\subsection{\texorpdfstring{$\Gamma_0(M)\le\pm\Gamma$}{Gamma0(M) in ±Gamma} by Dirichlet's theorem}

\begin{definition}[Replication at a prime]\label{ch16:def:repprime}
$f$ \emph{replicates at the prime $p$ with replicate $f'$} if there is $\Phi:\C\to\C$ with
$\Phi(f(\tau))=\sum_{j<p}f((\tau+j)/p)+f'(p\tau)$ for all $\tau$.
\end{definition}
\begin{leancode}
def ReplicatesAtPrime (f f' : ℍ → ℂ) (p : ℕ) (hp : p.Prime) : Prop :=
  ∃ Φ : ℂ → ℂ, ∀ τ : ℍ, Φ (f τ) = ∑ j ∈ Finset.range p, f (heckePoint 1 p j one_pos hp.pos τ) +
    f' (heckePoint p 1 0 hp.pos one_pos τ)
\end{leancode}
(\leanpath{VOA/Moonshine/MoonshineCNGammaZeroBridge.lean}.)
\begin{dictionary}
Exactly the definition.  The lemma \lean{heckeSum_prime} (same file) shows that the divisor-sum of
\lean{CompletelyReplicable} at $m=p$ prime has exactly these two blocks ($a=1$: $j=0,\dots,p-1$; $a=p$: one term),
and \lean{replicatesAtPrime_of_replication} turns complete replicability of $g$ at $p$ into
\lean{ReplicatesAtPrime (T g) (T (g ^ p)) p}.
\end{dictionary}
\begin{physics}
The $p$-th Hecke covering relation, in its simplest (prime) form.
\end{physics}
\begin{status}
Definition; used by the audited theorems below.
\end{status}

\begin{theorem}[$\Gamma_0$ from $\Gamma_1$ and replication]\label{ch16:thm:gamma0}
Let $f$ be a normalised Hauptmodul for $\Gamma\le\chsixteenSLR$ with $\Gamma_1(M)\le\Gamma$, $M\ge1$, and suppose that for
every prime $p>M$ there is a holomorphic, $1$-periodic, normalised $f'$ such that $f$ replicates at $p$ with replicate
$f'$.  Then:
\begin{enumerate}[label=(\alph*)]
\item $\Gamma_0(M)\le\pm\Gamma$;
\item $f$ is a normalised Hauptmodul for $\pm\Gamma$;
\item if moreover $\Gamma$ is commensurable with $\chsixteenSL$, then $\pm\Gamma$ is a moonshine group of level $M$.
\end{enumerate}
\end{theorem}
\begin{leancode}
theorem gamma0_le_adjoinNegOne_of_replicable {Γ : Subgroup (GL (Fin 2) ℝ)} {f : ℍ → ℂ} {M : ℕ}
    (hH : IsNormalizedHauptmodul Γ f) (hdet : Γ.HasDetOne) (hΓ1 : gamma1GL M ≤ Γ) (hM : 0 < M)
    (hrep : ∀ (p : ℕ) (hp : p.Prime), M < p → ∃ f' : ℍ → ℂ,
      MDifferentiable 𝓘(ℂ) 𝓘(ℂ) f' ∧ (∀ τ : ℍ, f' (glUpperUnipotent 1 • τ) = f' τ) ∧
      Tendsto (fun τ : ℍ => f' τ - (qPar τ)⁻¹) atImInfty (𝓝 0) ∧ ReplicatesAtPrime f f' p hp) :
    gamma0GL M ≤ Γ.adjoinNegOne

theorem IsNormalizedHauptmodul.adjoinNegOne {Γ : Subgroup (GL (Fin 2) ℝ)} {f : ℍ → ℂ}
    (hH : IsNormalizedHauptmodul Γ f) : IsNormalizedHauptmodul Γ.adjoinNegOne f where

theorem isMoonshineGroup_adjoinNegOne_of_replicable {Γ : Subgroup (GL (Fin 2) ℝ)} {f : ℍ → ℂ}
    {M : ℕ} (hH : IsNormalizedHauptmodul Γ f) (hdet : Γ.HasDetOne) (hΓ1 : gamma1GL M ≤ Γ)
    (hM : 0 < M) (hcomm : Subgroup.Commensurable Γ 𝒮ℒ)
    (hrep : ∀ (p : ℕ) (hp : p.Prime), M < p → ∃ f' : ℍ → ℂ,
      MDifferentiable 𝓘(ℂ) 𝓘(ℂ) f' ∧ (∀ τ : ℍ, f' (glUpperUnipotent 1 • τ) = f' τ) ∧
      Tendsto (fun τ : ℍ => f' τ - (qPar τ)⁻¹) atImInfty (𝓝 0) ∧ ReplicatesAtPrime f f' p hp) :
    IsMoonshineGroup Γ.adjoinNegOne M ∧ IsNormalizedHauptmodul Γ.adjoinNegOne f
\end{leancode}
(\leanpath{VOA/Moonshine/MoonshineCNGammaZeroBridge.lean}.)
\begin{dictionary}
\lean{Γ.adjoinNegOne} is $\pm\Gamma$.  \lean{hrep} quantifies over primes $p>M$ and provides, for each, $f'$ with the
three properties and \lean{ReplicatesAtPrime f f' p hp}.  (c) also needs \lean{hcomm}, commensurability of $\Gamma$.
The conclusions are (a), (b), (c) exactly.  Note that the output group is $\pm\Gamma$, not $\Gamma$; this is harmless
because $-1$ acts trivially on $\chsixteenH$ and on cusps, and $\Gamma_0(M)$ contains $-1$ anyway.
\end{dictionary}
\paragraph{Proof on paper.}  (a) Let $\delta=\chsixteenmat abcd\in\Gamma_0(M)$.  Then $ad\equiv1\pmod M$, so $d$ is a unit
mod $M$.  By Dirichlet's theorem on primes in arithmetic progressions (Mathlib's
\lean{Nat.forall_exists_prime_gt_and_eq_mod}) there is a prime $p>M$ with $p\equiv d\pmod M$; then $\gcd(p,M)=1$ and we
choose $m,n$ with $mp-nM=1$.  Theorem~\ref{ch16:thm:main} gives $\pm\sigma\in\Gamma$.  The matrix
$\varepsilon=\delta\sigma^{-1}$ has integer entries, determinant $1$, lower-left entry $cp-dM\equiv0$, and diagonal
entries $\equiv ap\equiv ad\equiv1$ and $\equiv dm\equiv dp^{-1}\equiv1\pmod M$; so $\varepsilon\in\Gamma_1(M)\le\Gamma$ and
$\delta=\varepsilon\sigma\in\pm\Gamma$.  (b) $-1$ acts trivially on $\chsixteenH$ and on $\chsixteenP(\chsixteenR)$, so each of
the eight fields transfers from $\Gamma$ to $\pm\Gamma$ (the cusps and other cusps are the same).  (c) $\pm\Gamma$ has
determinant one, contains $\Gamma_0(M)$ by (a), and is commensurable with $\Gamma$ (index at most $2$), hence with
$\chsixteenSL$. \qed
\begin{physics}
Dirichlet's theorem supplies enough primes to realise every residue class of twists $g\mapsto g^d$; each realisation is
a symmetry of the Hauptmodul by Theorem~\ref{ch16:thm:main}.  The upshot is the full $\Gamma_0(M)$ invariance expected of
an order-$M$ twisted partition function whose twist is ``Galois-invariant''.
\end{physics}
\begin{status}
Unconditional (with the stated hypotheses); all three printed in \leanpath{VOA/Moonshine/MoonshineCNGammaZeroAudit.lean},
standard axioms.
\end{status}

\subsection{The bridge}

\begin{theorem}[The bridge $\Gamma_1\Rightarrow\Gamma_0$]\label{ch16:thm:bridge}
\begin{enumerate}[label=(\alph*)]
\item If the $\Gamma_1$ form holds and $g$ has finite order and replicates at every prime (with some polynomial), then
  $g$ has the Conway--Norton property.
\item If every automorphism of $\Vnat$ has finite order, the $\Gamma_1$ form holds, and the finite-order McKay--Thompson
  series are completely replicable, then the Conway--Norton conjecture holds.
\item The same with \lean{CompletelyReplicable} in place of its finite-order restriction.
\end{enumerate}
\end{theorem}
\begin{leancode}
def CompletelyReplicableOfFinOrder : Prop :=
  ∀ (g : AutVnat), IsOfFinOrder g → ∀ m : ℕ, 0 < m → ∃ P : Polynomial ℂ, P.Monic ∧
    P.natDegree = m ∧ ∀ τ : ℍ, P.eval (mckayThompsonFun g τ) =
      ∑ a ∈ m.divisors.attach, ∑ b ∈ Finset.range (m / a),
        mckayThompsonFun (g ^ (a : ℕ)) (heckePoint a (m / a) b (Nat.pos_of_mem_divisors a.2)
          (Nat.div_pos (Nat.divisor_le a.2) (Nat.pos_of_mem_divisors a.2)) τ)

theorem conwayNortonFor_of_gammaOne (h1 : ConwayNortonGammaOne) {g : AutVnat}
    (hg : IsOfFinOrder g)
    (hrep : ∀ p : ℕ, p.Prime → ∃ P : Polynomial ℂ, ∀ τ : ℍ, P.eval (mckayThompsonFun g τ) =
      ∑ a ∈ p.divisors.attach, ∑ b ∈ Finset.range (p / a),
        mckayThompsonFun (g ^ (a : ℕ)) (heckePoint a (p / a) b (Nat.pos_of_mem_divisors a.2)
          (Nat.div_pos (Nat.divisor_le a.2) (Nat.pos_of_mem_divisors a.2)) τ)) :
    ConwayNortonFor g

theorem conwayNorton_of_gammaOne_of_finOrder (hfin : ∀ g : AutVnat, IsOfFinOrder g)
    (h1 : ConwayNortonGammaOne) (hrep : CompletelyReplicableOfFinOrder) :
    ConwayNortonConjecture

theorem conwayNorton_of_gammaOne (hfin : ∀ g : AutVnat, IsOfFinOrder g)
    (h1 : ConwayNortonGammaOne) (hrep : CompletelyReplicable) : ConwayNortonConjecture
\end{leancode}
(\leanpath{VOA/Moonshine/MoonshineCNGammaZeroBridge.lean}.)
\begin{dictionary}
\lean{CompletelyReplicableOfFinOrder} is \lean{CompletelyReplicable} restricted to $g$ with \lean{IsOfFinOrder g}.
In (a), \lean{hrep} asks for replication at every prime $p$ with a polynomial $P$ that need not be monic.  The
hypotheses of (b) and (c) are named propositions: \lean{hfin : ∀ g, IsOfFinOrder g}, \lean{h1 : ConwayNortonGammaOne},
\lean{hrep}.  The conclusions are \lean{ConwayNortonFor g} and \lean{ConwayNortonConjecture}.  These theorems are
\textbf{conditional} on their hypotheses; \texttt{\#print axioms} does not show them.  All three hypotheses are
theorems of the library (Chapters~\ref{ch:19} and~\ref{ch:20}); their combination is \lean{conwayNorton}.
\end{dictionary}
\paragraph{Proof on paper.}  (a) The $\Gamma_1$ form for $g$ gives $(\Gamma,M)$.  For a prime $p>M$, the replicate is
$f'=T_{g^p}$: it is $1$-periodic by definition, and holomorphic and normalised because the $\Gamma_1$ form also applies
to $g^p$ (which has finite order).  Replication of $g$ at $p$ is \lean{ReplicatesAtPrime} by \lean{heckeSum_prime}.
Theorem~\ref{ch16:thm:gamma0} gives $\pm\Gamma$, a moonshine group of level $M$ for $T_g$.  (b), (c) Apply (a) to every
$g$. \qed
\begin{physics}
Exact twist-preserving modular invariance ($\Gamma_1$) plus the covering relations between twisted partition functions
imply invariance up to Galois twists ($\Gamma_0$).
\end{physics}
\begin{status}
Conditional on the explicit hypotheses \lean{ConwayNortonGammaOne}, (finite-order) complete replicability and finite
order; all printed in \leanpath{VOA/Moonshine/MoonshineCNGammaZeroAudit.lean} with standard axioms.  The final theorem
\lean{conwayNorton} uses \lean{conwayNorton_of_gammaOne} with all three hypotheses discharged (\S\ref{ch16:sec:final}).
\end{status}

\subsection{Galois invariance}

\begin{theorem}[Galois invariance of McKay--Thompson series]\label{ch16:thm:galois}
Under the hypotheses of Theorem~\ref{ch16:thm:bridge}(a) (respectively (c)), $T_{g^u}=T_g$ for every $u$ coprime to the
order of $g$ (respectively for every $g$ and every such $u$).
\end{theorem}
\begin{leancode}
theorem mckayThompsonFun_pow_eq_of_coprime (h1 : ConwayNortonGammaOne) {g : AutVnat}
    (hg : IsOfFinOrder g)
    (hrep : ∀ p : ℕ, p.Prime → ∃ P : Polynomial ℂ, ∀ τ : ℍ, P.eval (mckayThompsonFun g τ) =
      ∑ a ∈ p.divisors.attach, ∑ b ∈ Finset.range (p / a),
        mckayThompsonFun (g ^ (a : ℕ)) (heckePoint a (p / a) b (Nat.pos_of_mem_divisors a.2)
          (Nat.div_pos (Nat.divisor_le a.2) (Nat.pos_of_mem_divisors a.2)) τ))
    {u : ℕ} (hu : Nat.Coprime u (orderOf g)) :
    mckayThompsonFun (g ^ u) = mckayThompsonFun g

theorem mckayThompsonFun_pow_eq_of_coprime_of_completelyReplicable
    (hfin : ∀ g : AutVnat, IsOfFinOrder g) (h1 : ConwayNortonGammaOne)
    (hrep : CompletelyReplicable) (g : AutVnat) {u : ℕ} (hu : Nat.Coprime u (orderOf g)) :
    mckayThompsonFun (g ^ u) = mckayThompsonFun g
\end{leancode}
(\leanpath{VOA/Moonshine/MoonshineCNGammaZeroGalois.lean}.)
\begin{dictionary}
\lean{Nat.Coprime u (orderOf g)}: $\gcd(u,o(g))=1$.  The conclusion is an equality of functions on $\chsixteenH$.
Conditional on the stated hypotheses (all theorems of the library; the unconditional version is
\lean{mckayThompsonFun_pow_eq_of_coprime_holds}, Chapter~\ref{ch:22}).
\end{dictionary}
\paragraph{Proof on paper.}  Let $(\Gamma,M)$ come from the $\Gamma_1$ form for $g$.  By Dirichlet there is a prime
$p>M$ with $p\equiv u\pmod{o(g)}$, so $g^p=g^u$.  Theorem~\ref{ch16:thm:main} with $f'=T_{g^p}$ gives
$T_{g^u}=T_{g^p}=T_g$. \qed

Comparing $q$-coefficients, $\tr(g^u|\Vnat_n)=\tr(g|\Vnat_n)$, and averaging over $u\in(\Z/o(g))^\times$ shows that the
traces are rational (Ramanujan sums); this is how \lean{rationalTraces} is proved in Chapter~\ref{ch:22}.
\begin{physics}
The twisted partition function of $g$ depends only on the cyclic group generated by $g$: twisting by any generator
$g^u$ gives the same function.  Equivalently the character values of $g$ on every energy level are rational.
\end{physics}
\begin{status}
Conditional on the stated hypotheses; both printed in \leanpath{VOA/Moonshine/MoonshineCNGammaZeroAudit.lean},
standard axioms.
\end{status}

\subsection{The family version}

\begin{theorem}[$\Gamma_0$ for an abstract replicable family]\label{ch16:thm:family}
Let $F$ be a completely replicable family, $k\in\N$, and suppose the member $Q_{c(k,\cdot)}$ is a normalised Hauptmodul
for $\Gamma\le\chsixteenSLR$ with $\Gamma_1(M)\le\Gamma$, $M\ge1$.
\begin{enumerate}[label=(\alph*)]
\item If $\Gamma$ is commensurable with $\chsixteenSL$ and each replicate $Q_{c(kp,\cdot)}$, $p>M$ prime, is holomorphic
  and normalised, then $\pm\Gamma$ is a moonshine group of level $M$ and $Q_{c(k,\cdot)}$ is a normalised Hauptmodul for
  it.
\item For each prime $p>M$ with $Q_{c(kp,\cdot)}$ holomorphic and normalised, $Q_{c(kp,\cdot)}=Q_{c(k,\cdot)}$.
\end{enumerate}
\end{theorem}
\begin{leancode}
theorem isMoonshineGroup_of_gammaOne (k : ℕ) {Γ : Subgroup (GL (Fin 2) ℝ)} {M : ℕ}
    (hM : 0 < M) (hdet : Γ.HasDetOne) (hΓ1 : gamma1GL M ≤ Γ)
    (hcomm : Subgroup.Commensurable Γ 𝒮ℒ) (hH : IsNormalizedHauptmodul Γ (qSeriesFun (F.c k)))
    (hhol : ∀ p : ℕ, p.Prime → M < p → MDifferentiable 𝓘(ℂ) 𝓘(ℂ) (qSeriesFun (F.c (k * p))))
    (hnorm : ∀ p : ℕ, p.Prime → M < p →
      Tendsto (fun τ : ℍ => qSeriesFun (F.c (k * p)) τ - (qPar τ)⁻¹) atImInfty (𝓝 0)) :
    IsMoonshineGroup Γ.adjoinNegOne M ∧
      IsNormalizedHauptmodul Γ.adjoinNegOne (qSeriesFun (F.c k))

theorem qSeriesFun_mul_prime_eq (k : ℕ) {Γ : Subgroup (GL (Fin 2) ℝ)} {M : ℕ}
    (hM : 0 < M) (hdet : Γ.HasDetOne) (hΓ1 : gamma1GL M ≤ Γ)
    (hH : IsNormalizedHauptmodul Γ (qSeriesFun (F.c k))) {p : ℕ} (hp : p.Prime) (hMp : M < p)
    (hhol : MDifferentiable 𝓘(ℂ) 𝓘(ℂ) (qSeriesFun (F.c (k * p))))
    (hnorm : Tendsto (fun τ : ℍ => qSeriesFun (F.c (k * p)) τ - (qPar τ)⁻¹) atImInfty (𝓝 0)) :
    qSeriesFun (F.c (k * p)) = qSeriesFun (F.c k)
\end{leancode}
(\leanpath{VOA/Moonshine/MoonshineCNGammaZeroFamily.lean}, namespace \chsixteenlean{VOA.Moonshine.ConwayNorton.ReplicableFamily},
with \lean{F : ReplicableFamily N} a variable.)
\begin{dictionary}
The same statements as Theorems~\ref{ch16:thm:gamma0} and~\ref{ch16:thm:main}, for member $k$ of an abstract family, with
replicate \lean{qSeriesFun (F.c (k * p))}.  The lemma \lean{ReplicableFamily.replicatesAtPrime} extracts the two-block
replication identity from field \lean{replicable}, and \lean{qSeriesFun_periodic} supplies periodicity.  These
theorems are not on the path to \lean{conwayNorton}; they were provided so that the genus-zero track could work on the
abstract interface alone, and (b) says that the abstract family is automatically ``Galois invariant'' once its members
are Hauptmoduln.
\end{dictionary}
\paragraph{Proof on paper.}  As for Theorems~\ref{ch16:thm:gamma0} and~\ref{ch16:thm:main}. \qed
\begin{physics}
The same statement for an abstract system of twisted partition functions.
\end{physics}
\begin{status}
Conditional on the stated hypotheses; printed in \leanpath{VOA/Moonshine/MoonshineCNGammaZeroAudit.lean}, standard
axioms.
\end{status}

\subsection{Supporting declarations}

The five files contain, besides the statements above, the following supporting material; none of it introduces a new
mathematical notion.
\begin{itemize}
\item \leanpath{VOA/Moonshine/MoonshineCNGammaZeroShimizu.lean}: the invariance subgroup
  \chsixteenlean{precompInvariantSubgroup ψ} $=\{g:\psi\circ g=\psi\}$; real M\"obius maps of positive determinant fixing two
  points of $\chsixteenH$ are scalar (\chsixteenlean{entries_of_two_fixed}, \chsixteenlean{smul_eq_self_of_scalar}), hence a non-trivial
  one moves $i$ or $2i$ (\chsixteenlean{moves_I_or_twoI}, \chsixteenlean{exists_smul_ne_of_lowerLeft_ne_zero}); entrywise convergence
  $C_n\to1$ gives $C_n\tau\to\tau$ and eventually positive determinant; the identity-theorem step
  \chsixteenlean{eq_of_invariant_of_frequently_moved}; Shimizu's sequence \chsixteenlean{shimizuPair}, \chsixteenlean{shimizuMatrix},
  \chsixteenlean{shimizuWitness} with its recursion, convergence and non-triviality.
\item \leanpath{VOA/Moonshine/MoonshineCNGammaZeroCosets.lean}: integer and real matrix constructors
  (\chsixteenlean{slOfDetOne}, \chsixteenlean{upperSL}, \chsixteenlean{lowerSL}, \chsixteenlean{glUpperTri}, \chsixteenlean{heckeShiftGL} $=\chsixteenmat1j0p$,
  \chsixteenlean{heckeScaleGL} $=\chsixteenmat p001$) and their actions on Hecke points; $\sigma T_1\sigma^{-1}\in\Gamma_1(M)$
  (\chsixteenlean{diamondLift_mul_upperSL}); M\"obius arithmetic in a field and the index maps \chsixteenlean{heckeIndexMap},
  \chsixteenlean{heckeIndexInv} with their bijectivity and integrality lemmas; the three coset identities; $p\,L_M\tau=L_{M/p}(p\tau)$
  (\chsixteenlean{heckeScale_lowerSL}) and surjectivity of $\tau\mapsto p\tau$.
\item \leanpath{VOA/Moonshine/MoonshineCNGammaZeroBridge.lean}: $|q^{-1}|=e^{2\pi\operatorname{Im}\tau}\to\infty$
  (\chsixteenlean{norm_inv_qPar}, \chsixteenlean{tendsto_norm_inv_qPar_atTop}); periodicity of $q$ and of $T_g$
  (\chsixteenlean{mckayThompsonFun_periodic}); an elementary inequality $|u-v|\le|e^u-e^v|$ for $u,v\ge0$; $-g$ and $g$ have the
  same action on cusps, so other cusps of $\Gamma$ and $\pm\Gamma$ agree (\chsixteenlean{neg_smul_cuspInfty},
  \chsixteenlean{isOtherCusp_adjoinNegOne_iff}); \chsixteenlean{hasDetOne_adjoinNegOne}; \chsixteenlean{heckePointAux} (a proof-free variant of
  \chsixteenlean{heckePoint}), \chsixteenlean{heckeSum_prime}, \chsixteenlean{replicatesAtPrime_of_replication},
  \chsixteenlean{completelyReplicableOfFinOrder_of_completelyReplicable}.
\item \leanpath{VOA/Moonshine/MoonshineCNGammaZeroFamily.lean}: \chsixteenlean{qSeriesFun_periodic},
  \chsixteenlean{ReplicableFamily.replicatesAtPrime}.
\end{itemize}
All of these are printed in \leanpath{VOA/Moonshine/MoonshineCNGammaZeroAudit.lean} or used by declarations printed
there.

\begin{warning}
\textbf{Deviations from the paper blueprint, none a defect}.
(i) Constancy uses the identity theorem, not local injectivity of $\psi$.  (ii) That $\sigma\infty$ is a cusp of $\Gamma$
comes from arithmeticity of $\Gamma_1(M)$, not from the parabolic $\sigma T\sigma^{-1}$.  (iii) Coset matching holds for
every prime $p$; $p>M$ is used only for $0<M/p<1$.  (iv) Replication is used with an arbitrary function $\Phi$; monicity
and degree are never used.  (v) The output group is $\pm\Gamma$.  (vi) Holomorphy and normalisation of the replicate
$T_{g^p}$ are taken from the $\Gamma_1$ theorem applied to $g^p$, so the bridge uses the $\Gamma_1$ theorem at $g$
\emph{and} at all $g^p$.
\end{warning}


%======================================================================
\section{The final theorem and the audit modules}\label{ch16:sec:final}
%======================================================================

\subsection{The theorem}

\begin{theorem}[Conway--Norton for $\Aut(\Vnat)$; Borcherds]\label{ch16:thm:conwayNorton}
For every automorphism $g$ of $\Vnat$ there are a subgroup $\Gamma_g\le\chsixteenSLR$ and $N\ge1$ such that
$\Gamma_0(N)\le\Gamma_g$, $\Gamma_g$ is commensurable with $\chsixteenSL$, and the McKay--Thompson series
$T_g=q^{-1}+\sum_{n\ge1}\tr(g|\Vnat_{n+1})q^n$ is the normalised Hauptmodul of $\Gamma_g$ in the sense of Definition~\ref{ch16:def:cn}; in particular $\Gamma_g$ has
genus zero (this last consequence is Proposition~\ref{ch16:prop:haupt}, an argument on paper that is not formalized:
the Lean statement asserts the defining properties of a normalised Hauptmodul, not the topology of
$\Gamma_g\backslash\mathbb H^{*}$).
\end{theorem}
\begin{leancode}
theorem conwayNorton : ConwayNortonConjecture :=
  M8.conwayNortonConjecture_holds

example (g : AutMonster.AutVnat) : ∃ (Γ : Subgroup (GL (Fin 2) ℝ)) (N : ℕ), IsMoonshineGroup Γ N ∧
    IsNormalizedHauptmodul Γ (mckayThompsonFun g) :=
  conwayNorton g
\end{leancode}
(\leanpath{VOA/Moonshine/MoonshineCNFinalAudit.lean}, \lean{VOA.Moonshine.ConwayNorton.conwayNorton}.)
\begin{dictionary}
\lean{ConwayNortonConjecture} is Definition~\ref{ch16:def:cn}; the \lean{example} unfolds it one step, and Lean accepts
\lean{conwayNorton g} as its proof, so the unfolded form is literally what is proved.  The theorem has \textbf{no
hypotheses}: no arguments, no section variables, no type-class assumptions beyond the fixed objects.  Its proof term
\lean{M8.conwayNortonConjecture_holds} is the assembly of Chapter~\ref{ch:20}.  The full fidelity discussion is
\S\ref{ch16:sec:fidelity}: this is the genus-zero form of the Conway--Norton conjecture, for the full automorphism group
of the library's $\Vnat$.  ``Genus zero'' is read off on paper: Lean proves that $T_g$ has the defining properties of
a normalised Hauptmodul (\lean{IsNormalizedHauptmodul}), and Proposition~\ref{ch16:prop:haupt} shows, without Lean, that
a group admitting such a function has genus zero.
\end{dictionary}
\paragraph{Proof on paper; where each step lives.}  Unwinding the proof term, the theorem is the composite of the
following implications (names in \chsixteenlean{VOA.Moonshine.ConwayNorton}).
\begin{enumerate}[label=(\arabic*)]
\item \chsixteenlean{M8.monsterLieSimpleRootGeneration}: the Monster Lie algebra is generated by its simple root spaces
  (Chapter~\ref{ch:20}, from the product identity for $J(p)-J(q)$ of Chapter~\ref{ch:19}).
\item \chsixteenlean{M8.autVnat_isOfFinOrder}: every automorphism of $\Vnat$ has finite order (Chapter~\ref{ch:20}).
\item \chsixteenlean{completelyReplicable_of_simpleRootGeneration}: from (1), (2) and the twisted denominator identities,
  \chsixteenlean{CompletelyReplicable} (Chapter~\ref{ch:19}).
\item \chsixteenlean{conwayNortonGammaOne_of_completelyReplicable}: from complete replicability, the prime Hecke identity and the
  non-injectivity of the replicable families modulo translations (both proved, Chapter~\ref{ch:19}), the case-free
  genus-zero argument gives \chsixteenlean{ConwayNortonGammaOne}.
\item \chsixteenlean{conwayNorton_of_gammaOne} (Theorem~\ref{ch16:thm:bridge}): from (2), (4) and (3), \chsixteenlean{ConwayNortonConjecture}.
\end{enumerate}
The links are \chsixteenlean{M8.conwayNortonConjecture_holds} $=$ \chsixteenlean{M8.conwayNorton_of_isOfFinOrder} (2), which applies
\chsixteenlean{conwayNorton_of_simpleRootGeneration} to (1); that theorem calls
\chsixteenlean{M8.conwayNorton_of_finOrder_of_completelyReplicable} with (3), which in turn calls
\chsixteenlean{conwayNorton_of_completelyReplicable}, i.e.\ (4) followed by (5). \qed
\begin{physics}
Every global symmetry of the Monster conformal field theory has a twisted torus partition function that is the
Hauptmodul of a genus-zero modular group.  In the string-theory language of Borcherds' proof: the physical spectrum of
the $26$-dimensional bosonic string compactified on the moonshine orbifold times a Lorentzian torus organises the
twisted partition functions into a completely replicable system, and complete replicability forces genus zero.
\end{physics}
\begin{status}
\textbf{Unconditional.}  \texttt{\#print axioms} for \lean{conwayNorton}, \lean{M8.conwayNortonConjecture_holds},
\lean{M8.monsterLieSimpleRootGeneration} and \lean{M8.autVnat_isOfFinOrder} (the four commands at the end of
\leanpath{VOA/Moonshine/MoonshineCNFinalAudit.lean}): \lean{[propext, Classical.choice, Quot.sound]}, no
\lean{sorryAx}, as recorded in the build record of the release (Chapter~\ref{ch:15}, Section~\ref{ch15:sec:record}): the umbrella
\lean{VOA.Moonshine}, which imports this module, builds green (9,504 jobs), and \texttt{run\_axiom\_audits.py} passes
(5,983 reports across 25 probes, all within \lean{[propext, Classical.choice, Quot.sound]}, with 0 \lean{sorryAx}).
\end{status}

\subsection{The audit modules}

An \emph{audit module} is a Lean file containing \texttt{\#print axioms} commands (and sometimes \lean{example}s and
\texttt{\#check}s) for the declarations of a batch of files; building it makes Lean print the axioms on which each
declaration depends.  The moonshine development kept one cumulative chain of audits for the \texttt{MoonshineCN} files,
each importing the previous one.  The following are assigned to this chapter; the material they audit is documented in
the chapters indicated.  ``Reports'' is the number of \texttt{\#print axioms} lines in the file (counted in the source
files of the release).  By the axiom audit of the release, every report is \lean{[propext, Classical.choice, Quot.sound]}, with no
\lean{sorryAx}.
\begin{center}\small
\begin{longtable}{p{0.31\textwidth}p{0.07\textwidth}p{0.43\textwidth}p{0.08\textwidth}}
\toprule
file (\texttt{VOA/Moonshine/}) & reports & what is audited & chapter\\
\midrule
\endhead
\texttt{MoonshineCNAudit.lean} & 53 & graded traces and McKay--Thompson series (\texttt{Trace}); the lattice $I\!I_{1,1}$, its
  oscillators, the vertex-algebra target & \ref{ch:16}, \ref{ch:17}\\
\texttt{MoonshineCNAudit2.lean} & 383 & the complex-frame construction of $V_{I\!I_{1,1}}$: charges, currents, exponentials,
  conformal vector, Sugawara, Virasoro, reconstruction, grading & \ref{ch:17}\\
\texttt{MoonshineCNAudit3.lean} & 105 & trace bounds; \textbf{the statement file}; tensor fields, conformal structure,
  grading, momentum; Fock adjoint; lattice and tensor forms, nondegeneracy & \ref{ch:16}, \ref{ch:17}\\
\texttt{MoonshineCNAudit4.lean} & 97 & abstract string module: generators, orthogonals, grading, critical dimension,
  longitudinal, physical and transverse states (no-ghost core) & \ref{ch:18}\\
\texttt{MoonshineCNAudit5.lean} & 117 & quotient isometries, string symmetry, hyperbolic Fock space, commutant, Sugawara,
  form, straightening & \ref{ch:17}, \ref{ch:18}\\
\texttt{MoonshineCNAudit6.lean} & 98 & hyperbolic light-cone splitting: negative modes, filtrations, PBW maps and
  coordinates & \ref{ch:17}\\
\texttt{MoonshineCNAudit7.lean} & 94 & central scalar, mixed modes, positive creation, Sugawara adjoint, tensor light cone
  and PBW & \ref{ch:17}\\
\texttt{MoonshineCNAudit8.lean} & 89 & the concrete no-ghost isometries for $\Vnat$: sector grading and laws,
  $\Vnat$ string module, longitudinal quotient, zero weight & \ref{ch:18}\\
\texttt{MoonshineCNAudit9.lean} & 105 & near-VOA mode identities: Goddard, translation, $n$-th products, iterate,
  Borcherds identity, skew symmetry; sector form convention & \ref{ch:17}\\
\texttt{MoonshineCNAudit10.lean} & 83 & physical Lie algebra, Monster Lie algebra, momentum grading, sector frame,
  $\Aut(\Vnat)$ action, momentum reversal & \ref{ch:18}\\
\texttt{MoonshineCNAudit11.lean} & 76 & root spaces, root radical, descended bilinear form and pairing & \ref{ch:18}\\
\texttt{MoonshineCNAudit12.lean} & 66 & generator adjoints; Borcherds roots, data, $\mathfrak{sl}_2$, parabolic & \ref{ch:17}, \ref{ch:18}\\
\texttt{MoonshineCNAudit13.lean} & 60 & charge-sensitive nilpotence, opposite modes, adjoint reduction & \ref{ch:17}\\
\texttt{MoonshineCNAudit14.lean} & 111 & Borcherds algebras: augmentation, generation, Levi, parabolic and string actions,
  relative induction and PBW, Verma multiplication & \ref{ch:18}\\
\texttt{MoonshineCNAudit15.lean} & 80 & locally finite Casimir, root cutoff and duality, resolution range & \ref{ch:18}\\
\texttt{MoonshineCNAudit16.lean} & 43 & opposite iterate identity; hypothesis-free Lie invariance of the form on
  $\mathfrak m$ & \ref{ch:17}, \ref{ch:18}\\
\texttt{MoonshineCNAudit17.lean} & 141 & Verma resolution, tensor algebra $\cong$ enveloping algebra, freeness of the
  positive part, Cartan action, $\mathfrak m$ as Borcherds data (given generation) & \ref{ch:18}\\
\texttt{MoonshineCNAudit18.lean} & 188 & twisted denominator identities, formal and analytic replication, complete
  replicability; Conway--Norton given generation and finite order; imports the $\Gamma_0$ and genus-zero audits &
  \ref{ch:19}\\
\texttt{MoonshineCNGammaZeroAudit.lean} & 47 & the $\Gamma_1\Rightarrow\Gamma_0$ bridge (\S\ref{ch16:sec:gammazero}) &
  \ref{ch:16}\\
\texttt{MoonshineCNFinalAudit.lean} & 4 & \chsixteenlean{conwayNorton} and its three last inputs & \ref{ch:16}\\
\bottomrule
\end{longtable}
\end{center}
The audit modules prove nothing new; they exist to make the axiom check reproducible.  A clean audit line certifies
that the declaration uses no axiom beyond Lean's three and no \lean{sorry}; it does not certify that the declaration is
unconditional (several audited theorems have explicit hypotheses, e.g.\ \lean{conwayNorton_of_gammaOne}), and it does
not certify that the statement means what it should: that is the job of the fidelity checks.  The other final audits of
the part (\texttt{MoonshineCNMonsterFinalAudit}, \texttt{MoonshineCNSymmetryFinalAudit},
\texttt{MoonshineCNRationalAudit}, \texttt{MoonshineCNGenusZeroAudit*}, \texttt{MoonshineLRLiftAudit} and the
\texttt{MoonshineLR*Audit} files) are documented in Chapters~\ref{ch:19}--\ref{ch:22}.

%======================================================================
\section{Recorded traps and dead ends}\label{ch16:sec:traps}
%======================================================================

\begin{warning}
\textbf{Rationality does not follow from the replication identities alone}.  For
every $a\in\C$, the function $f=q^{-1}+aq$ satisfies all replication identities with replicates $f^{(k)}=q^{-1}+a^kq$.
Certificate: in the Hecke sum $\sum_{ad=m}\sum_{b<d}f^{(a)}((a\tau+b)/d)$ the sum over $b$ kills both the $q^{-1}$ and
the $q$ term unless $d=1$, leaving $q^{-m}+a^mq^m$; and the Dickson polynomial $D_m(x,a)$, defined by
$D_m(u+a/u,a)=u^m+(a/u)^m$, is monic of degree $m$ with $D_m(q^{-1}+aq,a)=q^{-m}+a^mq^m$.  With $a=\zeta_N$ the family
$k\mapsto q^{-1}+\zeta_N^kq$ is periodic of period $N$ and not rational.  \lean{ReplicableFamily} excludes it only
through the clause \lean{zero_eq_J}.  Hence Galois invariance must come from the Hauptmodul property, as in
\S\ref{ch16:sec:gammazero}, not from formal identities.
\end{warning}

\begin{warning}
\textbf{Rationality through a $\Q$-form of $\Vnat$ (the original plan) is circular or needs classification}.
The plan was: build a rational form of
$\Vnat$, show every automorphism preserves it, conclude rationality of traces by Galois conjugation.  That every $g$
preserves $\Vnat_2\cap V_\Q$ needs either $\Aut(\Vnat)=\langle C,\sigma\rangle$ (Tits-level group theory) or
rationality of every Ising vector (which needs the Monster).  The Galois twist $g\mapsto\tau g\tau^{-1}$ only gives
$\chi_n\circ c_\tau=\tau\circ\chi_n$ for an unknown automorphism $c_\tau$, which is circular.  Superseded by the bridge
of \S\ref{ch16:sec:gammazero}, which makes rationality a corollary.  Shimura reciprocity on $q$-coefficients (D4) was
also considered and rejected as unnecessary and expensive.
\end{warning}

\begin{warning}
\textbf{The index shift} between \lean{mckayThompson} (coefficient $n$ = trace on weight $n+1$) and \lean{JShift}
($=qJ$, coefficient $n$ = dimension of weight $n$) is easy to get wrong.  The library records it as the negative lemma
\lean{mckayThompson_one_ne_JShift} and the positive lemmas \lean{mckayThompson_one} and \lean{mckayThompsonShift_one}
(Proposition~\ref{ch16:prop:mtprops}).
\end{warning}

\begin{warning}
\textbf{Sign convention for momenta} (Chapters~\ref{ch:17}--\ref{ch:19}).  In Borcherds'
convention the root space of $\mathfrak m$ at $(m,n)$ is $\Vnat_{mn+1}$, for the form $\langle(m,n),(m,n)\rangle=-2mn$.
The library's lattice \lean{lorentzianLattice} has norm $+2xy$, so Borcherds' $(m,n)$ is the library's $(m,-n)$, and
the root space at library momentum $(x,y)$ is $\Vnat_{1-xy}$.  A product over library momenta $x,y\ge1$ is wrong: there
$\Vnat_{1-xy}$ is $\Vnat_0=\C\mathbf 1$ at $(1,1)$ (the real simple root) and zero for $xy\ge2$, so such a product sees
almost none of $\mathfrak m$.
\end{warning}

\begin{warning}
\textbf{``Every element has finite order'' is not ``the group is finite''}.  The
moonshine theorem needs only the former (\lean{autVnat_isOfFinOrder}); the latter, \lean{M8.finite_autVnat}, was proved
separately (Chapter~\ref{ch:20}).  Neither is used to identify $\Aut(\Vnat)$ with the Monster.
\end{warning}

\begin{warning}
\textbf{Literature steps that fail as printed} (Chapter~\ref{ch:19}).  Cummins--Gannon's
Prop.~5.1 is not usable verbatim for $\Gamma_1$ (its prime-clearing step uses exponents not in the admissible set), and
the last step of Cummins (2002), Prop.~2.1, is wrong as printed (right translations cannot change the lower-right entry
modulo the level); the library's genus-zero argument uses corrected versions.  On the Monster side
(Chapter~\ref{ch:22}), a first attempt at the commutant lemma used two code automorphisms that preserve the MOG column
partition, which cannot give transitivity on octads; a third, column-breaking involution was needed.
\end{warning}

%======================================================================
\section{Referee's checklist}\label{ch16:sec:checklist}
%======================================================================

All names are in \lean{VOA.Moonshine.ConwayNorton} unless stated; ``std.'' means \texttt{\#print axioms} $=$
\lean{[propext, Classical.choice, Quot.sound]}.
\begin{center}\small
\begin{longtable}{p{0.27\textwidth}p{0.32\textwidth}p{0.19\textwidth}p{0.11\textwidth}}
\toprule
declaration & intended meaning & hypotheses; status & check\\
\midrule
\endhead
\chsixteenlean{AutMonster.VnatVOA}, \chsixteenlean{AutMonster.AutVnat} & the moonshine VOA of Ch.~\ref{ch:14}; its automorphism group &
  definitions & Def.~\ref{ch16:def:vnatvoa}\\
\chsixteenlean{gradedAutAction}, \chsixteenlean{gradedAutTrace} & $g|_{V_n}$ and $\tr(g|V_n)$ & definitions & Def.~\ref{ch16:def:gradedtrace}\\
\chsixteenlean{gradedAutTrace_one}, \chsixteenlean{gradedAutTrace_conjugate}, \chsixteenlean{vnatAut_preserves_grading} & trace of $1$ is $\dim$;
  class function; $g\Vnat_n=\Vnat_n$ & none; std. & Prop.~\ref{ch16:prop:traces}\\
\chsixteenlean{mckayThompson}, \chsixteenlean{mckayThompsonShift} & $\sum\tr(g|\Vnat_{n+1})X^n$; $qT_g$ & definitions &
  Def.~\ref{ch16:def:mckayThompson}\\
\chsixteenlean{mckayThompson_coeff}, \chsixteenlean{mckayThompson_coeff_zero}, \chsixteenlean{mckayThompson_conjugate}, \chsixteenlean{mckayThompson_one},
  \chsixteenlean{mckayThompsonShift_one}, \chsixteenlean{mckayThompson_one_ne_JShift} & coefficients; $\tr(g|\Vnat_1)=0$; class function;
  $T_1=J$ formally & none; std. & Prop.~\ref{ch16:prop:mtprops}\\
\chsixteenlean{finiteOrderTrace_norm_le}, \chsixteenlean{gradedAutTrace_norm_le}, \chsixteenlean{mckayThompson_coeff_norm_le} &
  $|\tr|\le\dim$ for finite order & finite order; std. & Thm.~\ref{ch16:thm:tracebounds}\\
\chsixteenlean{qPar}, \chsixteenlean{mckayThompsonFun} & $q=e^{2\pi i\tau}$; $T_g$ on $\chsixteenH$ & definitions & Def.~\ref{ch16:def:mtfun}\\
\chsixteenlean{mckayThompsonFun_one} & $T_1=j-744$ & none; std. & Thm.~\ref{ch16:thm:T1}\\
\chsixteenlean{SendsInftyTo}, \chsixteenlean{HasCuspValue}, \chsixteenlean{IsOtherCusp} & cusp data & definitions & Def.~\ref{ch16:def:cuspdata}\\
\chsixteenlean{IsNormalizedHauptmodul}, \chsixteenlean{IsGenusZero} & normalised Hauptmodul; genus zero (width-one cusp) & definitions;
  equivalence with textbook notion on paper & Def.~\ref{ch16:def:hauptmodul}, Prop.~\ref{ch16:prop:haupt}\\
\chsixteenlean{gamma0GL}, \chsixteenlean{gamma1GL}, \chsixteenlean{IsMoonshineGroup} & $\Gamma_0(N)$, $\Gamma_1(N)$ in $\chsixteenGLR$; moonshine
  group & definitions & Def.~\ref{ch16:def:moonshinegroup}\\
\chsixteenlean{ConwayNortonFor}, \chsixteenlean{ConwayNortonConjecture}, \chsixteenlean{ConwayNortonVia} & the conjecture (genus-zero form) for
  $\Aut(\Vnat)$ & definitions & Def.~\ref{ch16:def:cn}, \S\ref{ch16:sec:fidelity}\\
\chsixteenlean{ConwayNortonGammaOne} & $\Gamma_1$ form, finite-order $g$ & definition & Def.~\ref{ch16:def:gammaone}\\
\chsixteenlean{heckePoint}, \chsixteenlean{CompletelyReplicable} & complete replicability & definitions & Def.~\ref{ch16:def:replicable}\\
\chsixteenlean{RationalTraces} & all traces rational & definition & Def.~\ref{ch16:def:rational}\\
\chsixteenlean{qSeriesFun}, \chsixteenlean{ReplicableFamily} & abstract replicable family & definitions & Def.~\ref{ch16:def:family}\\
\chsixteenlean{mckayThompsonFun_eq_qSeriesFun}, \chsixteenlean{mckayThompsonReplicableFamily}, \chsixteenlean{mckayThompson_replicableFamily_exists}
  & $T_g=Q_{c_g}$; the McKay--Thompson family & complete replicability, finite order; std. & Prop.~\ref{ch16:prop:mtfamily}\\
\chsixteenlean{eq_of_invariant_of_entries_tendsto_one}, \chsixteenlean{eq_of_invariant_upper_lower_unipotent} & constancy; Shimizu &
  none; std. & Thm.~\ref{ch16:thm:shimizu}\\
\chsixteenlean{diamond_periodic}, \chsixteenlean{heckeLevelSum_lowerSL} & Hecke coset matching & $\Gamma_1(M)$-invariance; std. &
  Thm.~\ref{ch16:thm:cosets}\\
\chsixteenlean{diamond_sub_replicate_lowerUnipotent}, \chsixteenlean{diamond_eq_replicate_add_const} & $f\circ\sigma-f'$ constant &
  replication at $p$, $0<M<p$; std. & Thm.~\ref{ch16:thm:diamondconst}\\
\chsixteenlean{affine_eq_of_normalized} & asymptotic comparison & none; std. & Lem.~\ref{ch16:lem:affine}\\
\chsixteenlean{replicate_eq_and_diamond_mem} & $f'=f$, $\pm\sigma\in\Gamma$ & Hauptmodul for $\Gamma\supseteq\Gamma_1(M)$,
  replication at $p>M$; std. & Thm.~\ref{ch16:thm:main}\\
\chsixteenlean{gamma0_le_adjoinNegOne_of_replicable}, \chsixteenlean{IsNormalizedHauptmodul.adjoinNegOne},
  \chsixteenlean{isMoonshineGroup_adjoinNegOne_of_replicable} & $\Gamma_0(M)\le\pm\Gamma$ & replication at all $p>M$; std. &
  Thm.~\ref{ch16:thm:gamma0}\\
\chsixteenlean{conwayNortonFor_of_gammaOne}, \chsixteenlean{conwayNorton_of_gammaOne_of_finOrder}, \chsixteenlean{conwayNorton_of_gammaOne} &
  the bridge $\Gamma_1\Rightarrow\Gamma_0$ & \textbf{conditional}: $\Gamma_1$ form, replication, finite order; std. &
  Thm.~\ref{ch16:thm:bridge}\\
\chsixteenlean{mckayThompsonFun_pow_eq_of_coprime}, \chsixteenlean{..._of_completelyReplicable} & $T_{g^u}=T_g$ & \textbf{conditional} as
  above; std. & Thm.~\ref{ch16:thm:galois}\\
\chsixteenlean{ReplicableFamily.isMoonshineGroup_of_gammaOne}, \chsixteenlean{ReplicableFamily.qSeriesFun_mul_prime_eq} & family version &
  \textbf{conditional}; std. & Thm.~\ref{ch16:thm:family}\\
\chsixteenlean{conwayNorton} & \textbf{Conway--Norton (genus-zero form) for every automorphism of $\Vnat$} & \textbf{none};
  std. & Thm.~\ref{ch16:thm:conwayNorton}, \S\ref{ch16:sec:fidelity}\\
\chsixteenlean{GmatEqAutO} & AxiomMath's $\mathrm{Gmat}=\mathrm{AutO}$ & named hypothesis (a \chsixteenlean{Prop}), not an axiom &
  \S\ref{ch16:sec:gmat}\\
\bottomrule
\end{longtable}
\end{center}
\noindent\textbf{Files of this chapter and how they are covered.}  \leanpath{MoonshineCNStatement.lean},
\leanpath{MoonshineCNTrace.lean}, \leanpath{MoonshineCNTraceBounds.lean}, \leanpath{MoonshineCNFinalAudit.lean}: every
declaration individually.  \leanpath{MoonshineCNGammaZeroShimizu.lean}, \leanpath{...Cosets.lean},
\leanpath{...Bridge.lean}, \leanpath{...Galois.lean}, \leanpath{...Family.lean}: headline declarations individually,
supporting lemmas collectively (\S\ref{ch16:sec:gammazero}).  \leanpath{MoonshineCNGammaZeroAudit.lean} and
\leanpath{MoonshineCNAudit.lean}, \leanpath{MoonshineCNAudit2.lean}--\leanpath{MoonshineCNAudit18.lean}: as a list
(\S\ref{ch16:sec:final}).

%======================================================================
\section{Notes and references}\label{ch16:sec:notes}
%======================================================================

\paragraph{History.}  J.~G.~Thompson, \emph{Some numerology between the Fischer--Griess Monster and the elliptic modular
function}, Bull.\ London Math.\ Soc.\ 11 (1979) 352--353, and \emph{Finite groups and modular functions}, ibid.\
347--351 (McKay's observation is reported there).  J.~H.~Conway and S.~P.~Norton, \emph{Monstrous moonshine}, Bull.\
London Math.\ Soc.\ 11 (1979) 308--339: the conjecture, the tables of groups $\Gamma_0(n|h)+e,f,\dots$ and the first
coefficients.  A.~P.~Ogg, \emph{Automorphismes de courbes modulaires}, S\'eminaire Delange--Pisot--Poitou (1974/75),
for the supersingular primes.  R.~L.~Griess, \emph{The friendly giant}, Invent.\ Math.\ 69 (1982) 1--102.

\paragraph{The moonshine module.}  I.~B.~Frenkel, J.~Lepowsky, A.~Meurman, \emph{A natural representation of the
Fischer--Griess Monster with the modular function $J$ as character}, Proc.\ Natl.\ Acad.\ Sci.\ USA 81 (1984)
3256--3260; \emph{Vertex Operator Algebras and the Monster}, Pure and Applied Mathematics 134, Academic Press, 1988.
The construction in this library is documented in Chapters~\ref{ch:02}--\ref{ch:14}.

\paragraph{Borcherds' proof.}  R.~E.~Borcherds, \emph{Monstrous moonshine and monstrous Lie superalgebras}, Invent.\
Math.\ 109 (1992) 405--444.  The no-ghost theorem: P.~Goddard and C.~B.~Thorn, \emph{Compatibility of the dual Pomeron
with unitarity and the absence of ghosts in the dual resonance model}, Phys.\ Lett.\ B 40 (1972) 235--238.

\paragraph{Replication and genus zero.}  S.~P.~Norton, \emph{More on moonshine}, in Computational Group Theory
(Durham, 1982), Academic Press, 1984, 185--193 (replication formulae).  D.~Alexander, C.~Cummins, J.~McKay, C.~Simons,
\emph{Completely replicable functions}, in Groups, Combinatorics and Geometry (Durham, 1990), LMS Lecture Note Ser.\
165, CUP, 1992, 87--98.  Y.~Martin, \emph{On modular invariance of completely replicable functions}, in Moonshine, the
Monster, and Related Topics, Contemp.\ Math.\ 193, AMS, 1996, 263--286.  C.~J.~Cummins and T.~Gannon, \emph{Modular
equations and the genus zero property of moduli functions}, Invent.\ Math.\ 129 (1997) 413--443: completely
replicable functions with suitable growth are Hauptmoduln; the design source of the library's case-free route.
C.~J.~Cummins, \emph{Modular equations and discrete, genus-zero subgroups of $\mathrm{SL}(2,\mathbb R)$ containing
$\Gamma(N)$}, Canad.\ Math.\ Bull.\ 45 (2002) 36--45.  H.~Shimizu, \emph{On discontinuous groups operating on the product
of the upper half planes}, Ann.\ of Math.\ 77 (1963) 33--71, for the lemma of Theorem~\ref{ch16:thm:shimizu}.  F.~Diamond
and J.~Shurman, \emph{A First Course in Modular Forms}, GTM 228, Springer, 2005, for congruence subgroups, cusps,
modular curves and the Hecke cosets of $\Gamma_1(N)$.

\paragraph{Surveys.}  T.~Gannon, \emph{Moonshine beyond the Monster}, Cambridge University Press, 2006, the standard
reference (its \S5.2 is the source of the library's automorphism conventions).  J.~F.~R.~Duncan, M.~J.~Griffin,
K.~Ono, \emph{Moonshine}, Research in the Mathematical Sciences 2 (2015), article 11.  On the physics: L.~Dixon,
P.~Ginsparg, J.~Harvey, \emph{Beauty and the beast: superconformal symmetry in a Monster module}, Comm.\ Math.\ Phys.\
119 (1988) 221--241; M.~P.~Tuite, \emph{On the relationship between monstrous moonshine and the uniqueness of the
moonshine module}, Comm.\ Math.\ Phys.\ 166 (1995) 495--532; E.~Witten, \emph{Three-dimensional gravity revisited},
arXiv:0706.3359 (2007).

\paragraph{The Monster in Lean.}  AxiomMath, \emph{CFSG} (a Lean formalisation of parts of the classification of finite
simple groups), commit \texttt{69cb972f587a3ecc035e092253167ff2e303c962}, Lean v4.35; the theorem
\lean{CFSG.Blueprint.NativeSporadic.MonsterModel.gmat_eq_griess} is the content of \lean{GmatEqAutO}.  The fidelity of the port is
documented in \texttt{Notes/monster\_aut/CFSG\_PORT\_FIDELITY.md}.

\paragraph{The development record.}  The development was recorded in the authors' working notes, which are not part of
the release.  The statement file was written first and frozen before the proof.  Where this chapter states what the
machine confirmed, it relies on the build record and axiom audit of the release (Chapter~\ref{ch:15}).

\endgroup
